\documentclass[11pt]{article}
\usepackage[a4paper,margin=22mm]{geometry}
\usepackage{fontspec}
\setmainfont{DejaVu Serif}
\usepackage{amsmath,amssymb,amsthm,mathtools,graphicx,xcolor,hyperref,xurl}
\setlength{\emergencystretch}{4em}
\newtheorem{theo}{Theorem}[section]
\newtheorem{lemm}[theo]{Lemma}
\newtheorem{prop}[theo]{Proposition}
\newtheorem{coro}[theo]{Corollary}
\newtheorem{corol}[theo]{Corollary}
\newtheorem{conj}[theo]{Conjecture}
\newtheorem{Proposition}[theo]{Proposition}
\newtheorem{theoreme}[theo]{Theorem}
\theoremstyle{definition}
\newtheorem{defi}[theo]{Definition}
\newtheorem{exam}[theo]{Example}
\newtheorem{exem}[theo]{Example}
\newtheorem{rema}[theo]{Remark}
\newtheorem{nota}[theo]{Notation}
\newtheorem*{rema*}{Remark}
\newtheorem*{defi*}{Definition}
\newtheorem*{theo*}{Theorem}
\newtheorem*{lemm*}{Lemma}
\newtheorem*{prop*}{Proposition}
\newcommand{\enoncename}{}
\newtheorem{corpusenonce}[theo]{\enoncename}
\newenvironment{enonce}[1]{\renewcommand{\enoncename}{#1}\begin{corpusenonce}}{\end{corpusenonce}}
\newenvironment{enonce*}[1]{\par\medskip\noindent\textbf{#1.}\itshape}{\par\medskip}
\newenvironment{proofc}{\begin{proof}}{\end{proof}}
\newcommand{\Mc}{\mathcal{M}}
\newcommand{\pointir}{.\ ---\ }

\numberwithin{equation}{section}
\numberwithin{figure}{section}
\numberwithin{table}{section}
\def\endappendix{}
%~Mouliné par MaN_auto v.0.23.0 2020-05-19 15:13:25





\AtBeginDocument{
\newtheorem{fact}[theo]{Fact} }



\newtheorem{thmx}{Theorem}
\newtheorem{corox}[thmx]{Corollary}
\renewcommand*{\thethmx}{\Alph{thmx}}

%%%%%%%%%%Mathop%%%%%%%%%%%
\DeclareMathOperator{\ad}{ad}
\DeclareMathOperator{\Ad}{Ad}
\DeclareMathOperator{\Aut}{Aut}
\DeclareMathOperator{\Conf}{Conf}
\DeclareMathOperator{\GL}{GL}
\DeclareMathOperator{\Hom}{Hom}
\DeclareMathOperator{\iso}{Iso}
\DeclareMathOperator{\Iso}{Iso}
\DeclareMathOperator{\Ker}{Ker}
\DeclareMathOperator{\Nor}{Nor}
\DeclareMathOperator{\PSL}{PSL}
\DeclareMathOperator{\PGL}{PGL}
\DeclareMathOperator{\SL}{SL}
\DeclareMathOperator{\Sol}{SOL}
\DeclareMathOperator{\SO}{SO}
\DeclareMathOperator{\ein}{{\bf Ein}_3}

\DeclareMathOperator{\Omp}{O}%%\Omp
\DeclareMathOperator{\sol}{sol}

\DeclareMathOperator{\Span}{Span}
\DeclareMathOperator{\Sop}{S}
\DeclareMathOperator{\vol}{vol}
\DeclareMathOperator{\rk}{rk}
\DeclareMathOperator{\flatm}{flat}
\DeclareMathOperator{\Is}{Is}
\DeclareMathOperator{\Homeo}{Homeo}
\DeclareMathOperator{\hol}{hol}
\DeclareMathOperator{\Diff}{Diff}
\DeclareMathOperator{\Trace}{Trace}
\DeclareMathOperator{\Fix}{Fix}
\DeclareMathOperator{\rank}{rank}
\DeclareMathOperator{\loc}{loc}
\DeclareMathOperator{\id}{id}
\DeclareMathOperator{\Id}{Id}





\newcommand{\mint}{{M^{\intl}}}
\newcommand{\hmint}{{{\hat M}^{\intl}}}
\DeclareMathOperator{\intl}{int}


\DeclareMathOperator{\slop}{sl}
\newcommand{\sld}{\slop(2,\RR)}
\DeclareMathOperator{\kil}{{\mathfrak{kill}}}
\newcommand{\kiloc}{\kil^{\loc}}
\DeclareMathOperator{\heis}{{\mathfrak{heis}}}
\DeclareMathOperator{\Heis}{Heis}
\DeclareMathOperator{\SOL}{SOL}


%%%%%%%%%%%%%%%%%%%%%%%MATHCAL%%%%%%%%%%%%
\newcommand{\cald}{\mathcal{D}}
\newcommand{\calf}{\mathcal{F}}
\newcommand{\calm}{\mathcal{M}}
\newcommand{\calo}{\mathcal{O}}
\newcommand{\calw}{\mathcal{W}}

%%\newcommand*\O{\mathcal{O}}\CL


%%%%%%%%%%%%%%%%%%%%%%TILDE%%%%%%%%%%%%%
\newcommand{\tcalm}{\tilde{\mathcal{M}}}
\newcommand{\tm}{\tilde{{M}}}
\newcommand{\ddt}{\frac{\partial}{\partial t}}
\newcommand{\tX}{\tilde{{X}}}
\newcommand{\tY}{\tilde{{Y}}}
\newcommand{\tZ}{\tilde{{Z}}}
\newcommand{\tcalf}{\tilde{{\mathcal{F}}}}
\newcommand{\tFd}{\tilde{{F}}_{\Delta}}
\newcommand{\tF}{\tilde{{F}}}
\newcommand{\tD}{\tilde{{D}}}
\newcommand{\tth}{\tilde{h}}
\newcommand{\tcalfd}{\tilde{{\mathcal{F}_{\Delta}}}}

\newcommand{\calfd}{{{\mathcal{F}}}_{\Delta}}
\newcommand{\tcald}{\tilde{{\mathcal{D}}}}

\newcommand{\hcalm}{\hat{\mathcal{M}}}
\newcommand*\dk{{\mathcal{D}\kappa}}
\newcommand*\D{{\Delta}}

%%%%%%%%%%%%%%%%%%%MATHbb%%%%%%%%%
\newcommand*\NN{\mathbb{N}}
\newcommand*\TT{\mathbb{T}}
\newcommand*\RR{\mathbb{R}}
\newcommand*\ZZ{\mathbb{Z}}

%%%%%%%%%%%%%%%%%%%MATHFRAK%%%%%%%%%%%%%%%%
\newcommand{\lien}{{\mathfrak{n}}}
\newcommand{\lieg}{{\mathfrak{g}}}
\newcommand{\lieh}{{\mathfrak{h}}}
\newcommand{\liea}{{\mathfrak{a}}}
\newcommand{\tlien}{\tilde{\mathfrak{n}}}
\newcommand{\lies}{{\mathfrak{s}}}
\newcommand{\oo}{{\mathfrak{o}}}







\newcommand{\R}{{\bf R}}
\newcommand{\BP}{{\bf P}}


\DeclareMathOperator{\POop}{PO}
\newcommand{\PO}{{\POop(2,3)}}

\newcommand{\hm}{{\hat{M}}}
\newcommand{\hx}{{\hat{x}}}
\newcommand{\hz}{{\hat{z}}}


\newcommand{\delx}{{\frac{\partial}{\partial x_1}}}
\newcommand{\dely}{{\frac{\partial}{\partial x_2}}}
\newcommand{\delz}{{\frac{\partial}{\partial x_3}}}

\newcommand{\ka}{{\kappa}}


\let\oldtilde\tilde
\renewcommand*\tilde[1]{\mathchoice{\widetilde{#1}}{\widetilde{#1}}{\widetilde{#1}}{\oldtilde{#1}}}
\let\oldhat\hat
\renewcommand*\hat[1]{\mathchoice{\widehat{#1}}{\widehat{#1}}{\widehat{#1}}{\oldhat{#1}}}



\let\oldbar\bar
\renewcommand*{\bar}[1]{\mathchoice{\widebar{#1}}{\widebar{#1}}{\oldbar{#1}}{\oldbar{#1}}}
\let\oldoverline\overline
\renewcommand*{\overline}{\bar}

\newenvironment{smallpmatrix}{\left(\begin{smallmatrix}}{\end{smallmatrix}\right)}
\newcommand*{\sfrac}[2]{#1/#2}

\newcommand*{\dd}{\mathrm{d}}
\newcommand*{\dx}{\dd x}
\newcommand*{\dt}{\dd t}
\newcommand*{\du}{\dd u}
\newcommand*{\dv}{\dd v}
\newcommand*{\dw}{\dd w}


%%%%%%%%%%%%%%%%%%%%%%%%%%%%%%%%%%%%%%%%%%%%%%%%
\graphicspath{{./figures/}}

\newcommand*{\mk}{\mkern -1mu}
\newcommand*{\Mk}{\mkern -2mu}
\newcommand*{\mK}{\mkern 1mu}
\newcommand*{\MK}{\mkern 2mu}

\hypersetup{urlcolor=purple, linkcolor=blue, citecolor=red}

\newcommand*{\romanenumi}{\renewcommand*{\theenumi}{\roman{enumi}}}
\newcommand*{\Romanenumi}{\renewcommand*{\theenumi}{\Roman{enumi}}}
\newcommand*{\alphenumi}{\renewcommand*{\theenumi}{\alph{enumi}}}
\newcommand*{\Alphenumi}{\renewcommand*{\theenumi}{\Alph{enumi}}}
%%%%%%%%%%%%%%%%%%%%%%%%%%%%%%%%%%%%%%%%%%%%%%%%%%

\begin{document}
\section*{\raggedright 2600: Periodic warped universal covers of closed oriented and time-oriented Lorentz three-manifolds with noncompact isometries and a nonhomogeneous hyperbolic integrability component}
\noindent\textbf{Author:} Charles Frances\par\noindent\textbf{Source:} Lorentz dynamics on closed \$3\$-manifolds. Annales Henri Lebesgue 3 (2020), publisher version, DOI 10.5802/ahl.37. \url{https://ahl.centre-mersenne.org/item/AHL_2020__3__407_0/}
\par\noindent\textbf{License:} CC BY 4.0, \url{https://creativecommons.org/licenses/by/4.0/}. The original publisher PDF has the explicit license notice. See the saved source/license records.
\par\noindent\textbf{Editorial material note (not author text):} The selected problem is Theorem 5.2 for a nonlocally homogeneous hyperbolic component under the exact orientability, time-orientability and noncompact-isometry hypotheses. The complete author article body retains all local Cartan/integrability setup and the complete Anosov-torus, periodic warping, nondegeneration, KAK extension and return-map construction. Other classification results are context and are not additional counted problems. All proof text and formulas below are editable author TeX. Mathematical wording, language and proof order are retained. Only the article class, title matter and bibliography presentation are adapted for independent compilation; no proof has been generated or repaired.
\par\noindent\textbf{Editorial difficulty:} H2: The author extracts a compact Anosov Lorentz torus from generalized-curvature recurrence, proves nondegeneration and completeness of its normal flow using hyperbolic dynamics and a Lorentz isometric-extension theorem, then identifies the return quotient and warped universal cover. This global Lorentz-dynamical construction exceeds ordinary qualifying examinations.
\medskip\hrule\medskip
\noindent\textbf{Author statement, complete proof and local supporting text follow in source order.}
\section{Introduction} 
A celebrated theorem of Myers and Steenrod~\cite{myers-steenrod}, says that the isometry group of an $n$-dimensional Riemannian manifold is always a Lie transformation group of dimension at most $\frac{n(n+1)}{2}$. More precisely, this group is closed in the group of homeomorphisms and the Lie topology coincides with the compact-open topology. This property of the isometry group, including the bound on the dimension, carries over to general pseudo-Riemannian structures (see for instance~\cite{nomizu2} which deals with the more general case of affine connections). For \emph{closed} Riemannian manifolds, Ascoli's theorem readily implies that the isometry group is a compact Lie group. This compactness property is specific to the Riemannian world, and fails for general closed pseudo-Riemannian manifolds. For instance the Lorentz torus $ \RR^n /\ZZ^n$, endowed with the metric induced by $-\dx_1^2+\dx_2^2+\ldots+\dx_n^2$, has isometry group $\Omp(1,n-1)_{\ZZ}\ltimes\TT^n$. This group is noncompact, since $\Omp(1,n-1)_{\ZZ}$ is a lattice in $\Omp(1,n-1)$.

The geometrical implications of those two antagonistic phenomena (noncompact\-ness of the isometry group on the one hand, and compactness of the manifold on the other hand) were much studied in the Lorentzian case, on which we will focus here. A sample of significant results can be found in~\cite{adams.stuck}, \cite{gromov}, \cite{zeghib.isom}, \cite{zeghibisom1}, \cite{zeghibisom2}, \cite{zimmer.iso}, among a lot of other works.

One remarkable point is that the noncompactness of the isometry group is also expected to have strong topological consequences. This was first noticed by M.~Gromov in~\cite{gromov} when the isometry group is ``large'', for instance when it contains a noncompact simple Lie group (see also further developments in~\cite{fisher.zimmer}). Without any extra assumption on the acting group, let us mention the following striking result:



\begin{theo}[\cite{dambra}]\label{thm.dambra}
Let $(M,g)$ be a closed Lorentz manifold. We assume that $M$ and $g$ are real analytic, and $M$ is simply connected. Then $\Iso(M,g)$ is a compact group.
\end{theo}

The analyticity condition is crucial in the proof of Theorem~\ref{thm.dambra}, and when the dimension of the manifold is $> 3$, we actually don't know if the result holds in the smooth category.


The aim of this paper is to focus on the $3$-dimensional situation, and to provide a thorough study of all closed $3$-dimensional manifolds, which can be endowed with a Lorentz metric admitting a noncompact group of isometries.

\subsection{Statement of results}

Let us recall a class of closed $3$-manifolds which will play a prominent role in the sequel, namely the torus bundles over the circle (torus bundles for short). Let $\TT^2$ be a $2$-torus $ \RR^2/ \ZZ^2 $, and let us consider the product $[0,1] \times \TT^2$. We then make the identification $(0,x) \simeq (1,Ax)$, where $A$ is a given element of $ \SL(2,\ZZ)$. The resulting $3$-manifold is denoted $\TT_A^3$. When $A= \id$, we just get the $3$-torus $\TT^3$. If $A \in \SL(2,\ZZ)$ is \emph{hyperbolic}, namely is $\RR$-split with eigenvalues of modulus $\not =1$, we say that $\TT_A^3$ is a \emph{hyperbolic torus bundle}. If $A \in \SL(2,\ZZ)$ is \emph{parabolic}, namely conjugated to a unipotent matrix ($A \not =\id$), we say that $\TT_A^3$ is a \emph{parabolic torus bundle}. Finally, \emph{elliptic torus bundles} are those for which $A$ has finite order.

\subsubsection{A topological classification} Our first result is a topological classification of closed Lorentz $3$-manifolds admitting a noncompact isometry group.

\begin{thmx}\label{thm.topological}
Let $(M,g)$ be a smooth, closed $3$-dimensional Lorentz manifold. Assume that $(M,g)$ is orientable and time-orientable, and that $\Iso(M,g)$ is noncompact. Then $M$ is homeomorphic to one of the following spaces:
\begin{enumerate}
\item \label{list1.1} A quotient $\Gamma \backslash \tilde{\PSL}(2,\RR)$, where $\Gamma \subset \tilde{\PSL}(2,\RR)$ is any uniform lattice.
\item \label{list1.2}A $3$-torus ${\mathbb T}^3$, or a torus bundle ${\mathbb T}_A^3$, where $A \in \SL(2,\ZZ)$ can be any hyperbolic or parabolic element. 
\end{enumerate}
Conversely, any smooth compact $3$-manifold homeomorphic to one of the examples above can be endowed with a smooth Lorentz metric with a noncompact isometry group.
\end{thmx}

We recall that a Lorentz manifold is said to be time-orientable whenever it admits a vector field $X$ which is timelike everywhere, namely $g(X,X)<0$ on $M$. The assumption about orientability and time-orientability of the manifold made in the theorem is not really relevant, and one could drop it (adding a few allowed topological types) with extra case-by-case arguments in our proofs. Notice that any closed $3$-manifold will have a covering of order at most four satisfying the assumptions of Theorem~\ref{thm.topological}. We thus see that a lot of $3$-manifolds do not admit coverings appearing in the list of the theorem, where only four among the eight Thurston's geometries are represented. Hyperbolic manifolds are notably missing, and we can state:

\begin{coro}\label{coro.hyperbolicmanifold}
Let $M$ be a smooth closed $3$-dimensional manifold, which is homeo\-morphic to a complete hyperbolic manifold $\Gamma \backslash {\mathbb H}^3$, where $\Gamma<\Iso({\mathbb H}^3)$ is a torsion-free cocompact lattice. Then for every smooth Lorentz metric $g$ on $M$, the group $\Iso(M,g)$ is compact.
\end{coro}

\subsubsection{Continuous versus discrete isometries}

It is interesting to compare the conclusions of Theorem~\ref{thm.topological} to closely related results, and especially to the work~\cite{zeghibflot}, which was a great source of motivation for the present paper. In~\cite{zeghibflot}, A.~Zeghib studies $3$-dimensional closed manifolds admitting a non equicontinuous isometric flow. This hypothesis is actually equivalent to \emph{the noncompactness of the identity component $\Iso^o(M,g)$}. The classification can be briefly summarized as follows:


\begin{theo}\cite[Theorems~1 and 2]{zeghibflot}\label{thm.ghaniclassification}
Let $(M,g)$ be a smooth, closed $3$-dimensional Lorentz manifold. If the identity component $\Iso^o(M,g)$ is not compact, then:
\begin{enumerate}
\item \label{list1.3.1} Up to a finite cover, the manifold $M$ is homeomorphic either to a torus bundle $\TT_A^3$, with $A \in \SL(2,\ZZ)$ hyperbolic, or to a quotient $\Gamma \backslash \tilde{\PSL}(2,\RR)$, for a uniform lattice $\Gamma \subset \tilde{\PSL}(2,\RR)$.
\item \label{list1.3.2}The manifold $(M,g)$ is locally homogeneous. It is flat when $M$ is a hyperbolic torus bundle, and locally modelled on a Lorentzian, non-Riemannian, left-invariant metric on $\tilde{\PSL}(2,\RR)$ otherwise.
\end{enumerate}
\end{theo}

The definition of Lorentzian, non-Riemannian, left-invariant metrics on $\tilde{\PSL}(2,\RR)$ will be made precise in Section~\ref{sec.exemplessl2}.

Does it make a big difference, putting the noncompactness assumption on $\Iso^o(M,g)$ instead of $\Iso(M,g)$? At the topological level, notice that $3$-tori and parabolic torus bundles do not show up in Theorem~\ref{thm.ghaniclassification}. For Lorentz metrics on those manifolds, $\Iso^o(M,g)$ is always compact, but we will see that there exist suitable metrics $g$, for which the full group $\Iso(M,g)$ is noncompact. It means that for those examples, the noncompactness comes from the \emph{discrete part} $\Iso(M,g)/\Iso^o(M,g)$. Actually, there are instances of $3$-manifolds (see Section~\ref{sec.panorama}), where the isometry group is discrete, isomorphic to $\ZZ$.

To put more emphasis on how the general case may differ from the conclusions of~\cite{zeghibflot}, let us state the following existence result:

\begin{thmx}\label{thm.nonhomogeneous}
Let $M$ be a closed $3$-dimensional manifold which is homeomorphic to a $3$-torus $\TT^3$, or a torus bundle $\TT_A^3$, with $A \in \SL(2,\ZZ)$ hyperbolic or parabolic. Then it is possible to endow $M$ with time-orientable Lorentz metrics $g$ having the following properties:
\begin{enumerate}
\item \label{list2.1}The isometry group $\Iso(M,g)$ is noncompact, but the identity component $\Iso^o(M,g)$ is compact.
\item \label{list2.2}There is no open subset of $(M,g)$ which is locally homogeneous.
\end{enumerate}
\end{thmx}

The second property is equivalent to saying that orbit closures of the pseudo-group of local isometries have empty interior.

Observe that for any $3$-dimensional closed Lorentz manifold $(M,g)$ which is not locally homogeneous, $\Iso^o(M,g)$ is automatically compact by Theorem~\ref{thm.ghaniclassification} above.

The constructions leading to Theorem~\ref{thm.nonhomogeneous} are rather flexible. In particular, on $\TT^3$, or on any hyperbolic or parabolic torus bundle ${\TT_A^3}$, the moduli space of Lorentz metrics admitting a noncompact isometry group is by no means finite dimensional. This is again in sharp contrast with the second point of Theorem~\ref{thm.ghaniclassification}.

\subsubsection{Geometrical results}
The topological classification given by Theorem~\ref{thm.topological} comes as a byproduct of a finer, geometrical understanding of closed Lorentz $3$-manifolds with noncompact isometry group. This is the content of:

\begin{thmx}\label{thm.geometrical}
Let $(M,g)$ be a smooth, closed $3$-dimensional Lorentz manifold. Assume that $(M,g)$ is orientable and time-orientable, and that $\Iso(M,g)$ is noncompact.
\begin{enumerate}
\item \label{list3.1}If $M$ is homeomorphic to $\Gamma \backslash
\tilde{\PSL}(2,\RR)$, then $(M,g)$ is locally homogeneous, modelled on a Lorentzian, non-Riemannian, left-invariant metric on $\tilde{\PSL}(2,\RR)$.
\item \label{list3.2}If $M$ is homeomorphic to $\TT_A^3$, with $A \in \SL(2,\ZZ)$ hyperbolic, then there exists a smooth, positive, periodic function $a:\RR \to (0,\infty)$ such that the universal cover $(\tilde{M}^3,{\tilde{g}})$ is isometric to $\RR^3$ endowed with the metric
\[
{\tilde{g}}=\dt^2+2a(t)\du\dv.
\]
If $g$ is locally homogeneous, it is flat.
\item \label{list3.3}If $M$ is homeomorphic to $\TT_A^3$, with $A \in \SL(2,\ZZ)$ parabolic, then there exists a smooth, positive, periodic function $a:\RR \to (0,\infty)$ such that the universal cover $(\tilde{M}^3,{\tilde{g}})$ is isometric to $\RR^3$ endowed with the metric
\[
{\tilde{g}}=a(v)(\dt^2+2\du\dv).
\]
If $g$ is locally homogeneous, it is either flat or modelled on the Lorentz--Heisenberg geometry.
\item \label{list3.4}If $M$ is homeomorphic to a $3$-torus $\TT^3$, then the universal cover $(\tilde{M}^3,{\tilde{g}})$ is isometric to $\RR^3$ with a metric of type $2)$ or $3)$ above. If the metric $g$ is locally homogeneous, it is flat.
\end{enumerate}
\end{thmx}

The Lorentz--Heisenberg geometry will be described in Section~\ref{sec.lorentzheisenberg}.


We already emphasized that in some examples, the isometry group $\Iso(M,g)$ could be infinite discrete. However, it is worth mentioning that noncompactness of $\Iso(M,g)$ always produces somehow local continuous symmetries, a fact that will play a crucial role in our proofs. It is indeed easy to infer from Theorem~\ref{thm.geometrical} the following result.

\begin{corox}\label{coro.action.universel}
Let $(M,g)$ be a closed $3$-dimensional Lorentz manifold. If $\Iso(M,g)$ is noncompact, then $\Iso^o(\tm,\tilde{g})$ is noncompact. Actually $(\tm,\tilde{g})$ admits an isometric action of the group $\tilde{\PSL}(2,\RR)$, $\Heis$ or $\Sol$.
\end{corox}

All the results we stated so far assume that the Lorentz metric we are considering is smooth. Our proof actually uses a generalized curvature map which needs to be $C^1$, and involves 6 covariant derivatives of the Riemann curvature tensor. It follows that we need a regularity at least $C^9$. We don't know if metrics with low regularity may produce new examples.


\subsection{General strategy of the proof, and organization of the paper}

One aspect of the present work consists of existence results. This is the topic of Section~\ref{sec.panorama}, where we recollect well-known, and probably less known, examples of closed Lorentz $3$-manifolds having a noncompact isometry group. Examples are given, where $\Iso(M,g)$ is infinite discrete, or semi-discrete. This yields the existence part in Theorem~\ref{thm.topological}, and a proof of Theorem~\ref{thm.nonhomogeneous}.

The remaining part of the paper is then devoted to our classification results, namely Theorems~\ref{thm.topological} and~\ref{thm.geometrical}. The point of view we adopted, is that of Gromov's theory of rigid geometric structures~\cite{gromov}.

Section~\ref{sec.gromov} recall the main aspects of the theory, recast in the framework of Cartan geometry as in~\cite{melnick}, \cite{vincent}. The key result is the existence of a dense open subset $\mint \subset M$, called the \emph{integrability locus}, where Killing generators of finite order do integrate into genuine local Killing fields. Using the recurrence properties of the isometry group, this implies the crucial fact that the noncompactness of $\Iso(M,g)$ must produce a lot of local Killing fields (Proposition~\ref{prop.dimension}). Those continuous local symmetries, arising from a potentially discrete $\Iso(M,g)$, will be of great help to understand the geometry of the connected components of $\mint$, which can be roughly classified into three categories: constant curvature, hyperbolic, and parabolic (see Section~\ref{sec.components}). To unravel the global structure of $M$, we must understand how all the components of $\mint$ are patched together (notice that there can be infinitely many such components).

The first, and easiest case to study, is when all the components of $\mint$ are locally homogeneous. Results of~\cite{frances.open} show that $(M,g)$ itself is then locally homogeneous, allowing to understand $(M,g)$ completely. This is done in Section~\ref{sec.homogeneous}.

Section~\ref{sec.hyperbolic} studies the case where one component of $\mint$ is not locally homogeneous and hyperbolic. One then shows that $M$ is a $3$-torus or a hyperbolic torus bundle, and the geometry is that of examples~\eqref{list3.2} and~\eqref{list3.4} of Theorem~\ref{thm.geometrical}. This is summarized in Theorem~\ref{theo.hyperbolic}. The key feature in this case is to show that $(M,g)$ contains a Lorentz $2$-torus, on which an element $h \in \Iso(M,g)$ acts as an Anosov diffeomorphism (Lemma~\ref{lem.anosov.tore}). We then show that it is possible to push this Anosov torus by a kind of normal flow, to recover the topological, as well as geometrical structure of $(M,g)$.

The most tedious case to study is when $(M,g)$ is not locally homogeneous, and there are no hyperbolic components at all. This is the purpose of Sections~\ref{sec.local.geometry}, \ref{sec.geomeinstein} and~\ref{sec.global}. We show there that $M$ is a $3$-torus or a parabolic torus bundle, and the geometry is the one described in cases $3.$ and $4.$ of Theorem~\ref{thm.geometrical}. This is summarized in Theorem~\ref{th.theoplat}. The main observation here is that the manifold $(M,g)$ is conformally flat (Section~\ref{sec.local.geometry}). We then get a developing map $\delta: \tilde{M}^3 \to \ein$, which is a conformal immersion from the universal cover $\tilde{M}^3$ to a Lorentz model space $\ein$, called Einstein's universe. After introducing relevant geometric aspects of $\ein$ in Section~\ref{sec.geomeinstein}, we are in position to study in details the map $\delta: \tilde{M}^3 \to \ein$ in Section~\ref{sec.global}. We show that $\delta$ maps $\tilde{M}^3$ in a one-to-one way onto an open subset of $\ein$, which is conformally equivalent to Minkowski space. We are then reduced to the study of closed, flat, Lorentz $3$-manifolds with noncompact isometry groups, which was already done in Section~\ref{sec.homogeneous}.

All those partial results are recollected in Section~\ref{sec.fin}, where we see how they yield Theorem~\ref{thm.geometrical} and Corollary~\ref{coro.action.universel}.

\subsection*{Acknowledgments}

We warmly thank the referee(s) of this article for providing extremely valuable remarks on the initial draft.

\section{A panorama of examples}\label{sec.panorama}

The aim of this section is to construct a wide range of closed $3$-dimensional Lorentz manifolds $(M,g)$, with noncompact isometry group. Those examples will show that all topologies appearing in Theorem~\ref{thm.topological} do really occur. Moreover, Sections~\ref{sec.exempleshyperboliques}, \ref{sec.exemplesparaboliques} and~\ref{sec.exemplestorus}, prove our Theorem~\ref{thm.nonhomogeneous}. Part of the examples presented here are well known, others like those described in Section~\ref{sec.exempleshyperboliques}, \ref{sec.exemplesparaboliques}, \ref{sec.exemplestorus2} seem less classical, though elementary.


\subsection{Examples on quotients \texorpdfstring{$\Gamma \backslash \tilde{\PSL}(2,\RR)$}{Gamma backslash PSL(2, RR)}}\label{sec.exemplessl2}

The Lie group $\tilde{\PSL}(2,\RR)$, universal cover of ${\PSL}(2,\RR)$, admits a lot of interesting left-invariant Lorentzian metric. The most symmetric one is the \emph{anti-de Sitter} metric $g_{AdS}$. It is obtained by left-translating (a positive multiple of) the Killing form of the Lie algebra $\mathfrak{sl}(2,\RR)$. The space $(\tilde{\PSL}(2,\RR),g_{AdS})$ is a complete Lorentz manifold with constant sectional curvature $-1$, called anti-de Sitter space $\tilde{\bf{AdS}}_3$. Because the Killing form is $\Ad$-invariant, the metric $g_{AdS}$ is invariant by left and right multiplications of $\tilde{\PSL}(2,\RR)$ on itself. It follows that for any uniform lattice $\Gamma \subset \tilde{\PSL}(2,\RR)$, the metric $g_{AdS}$ induces a Lorentz metric $\overline{g}_{AdS}$ on the quotient manifold $\Gamma \backslash \tilde{\PSL}(2,\RR)$, with a noncompact isometry group coming from the right action of $\tilde{\PSL}(2,\RR)$ on $\Gamma \backslash \tilde{\PSL}(2,\RR)$.

There are other metrics than $g_{AdS}$ on $\tilde{\PSL}(2,\RR)$, which allow the same kind of constructions. They are obtained as follows. Exponentiating the linear space spanned by the matrix 
$\begin{smallpmatrix}
0&1\\
0&0\\
\end{smallpmatrix}$
(resp. $
\begin{smallpmatrix}
1&0\\
0&-1\\
\end{smallpmatrix}$), one gets a unipotent (resp. $\RR$-split) $1$-parameter group $\{\tilde{u}^t\}$ (resp. $\{ \tilde{h}^t \}$) in $\tilde{\PSL}(2,\RR)$. The adjoint action of each of those flows, admits invariant Lorentz scalar products on $\mathfrak{sl}(2,\RR)$, which are not equal to a multiple of the Killing form. One can left-translate those scalar products and get metrics $g_u$ and $g_h$ on $\tilde{\PSL}(2,\RR)$ which are respectively $\tilde{\PSL}(2,\RR)
\times \{\tilde{u}^t\}$ and $\tilde{\PSL}(2,\RR)
\times \{\tilde{h}^t\}$-invariant. Actually there are families of such metrics $g_u$ and $g_h$ which are not pairwise isometric. Now, for each uniform lattice $\Gamma \subset \tilde{\PSL}(2,\RR)$, the quotient $\Gamma \backslash \tilde{\PSL}(2,\RR)$ can be endowed with induced metrics $\overline{g}_u$ or $\overline{g}_h$ carrying an isometric, noncompact action of $\RR$, coming from the right actions of, respectively, $\{\tilde{u}^t\}$ and $\{\tilde{h}^t\}$ on $\tilde{\PSL}(2,\RR)$.

In the sequel, the metric $g_{AdS}$ and metrics of the form $g_u$ or $g_h$, will be referred to as \emph{Lorentzian, non-Riemannian, left-invariant metrics on} $\tilde{\PSL}(2,\RR)$. Those are the only left-invariant metrics on $\tilde{\PSL}(2,\RR)$, the isometry group of which does not preserve a Riemannian metric.



\subsection{Examples on hyperbolic torus bundles}\label{sec.exempleshyperboliques}


Let us start with the space $\RR^3$ endowed with coordinates $(x_1,x_2,t)$ associated to a basis $(e_1,e_2,e_t)$. We consider a hyperbolic matrix $A$ in $\SL(2,\ZZ)$. Hyperbolic means that $A$ has two distinct real eigenvalues $\lambda$ and $\lambda^{-1}$ different from $\pm 1$.


Let us consider the group $\Gamma$ generated by $\gamma_1=T_{e_1}$ (the translation of vector $e_1$), $\gamma_2=T_{e_2}$ and the affine transformation $\gamma_3=
\begin{smallpmatrix} A & 0\\ 0 & 1
\end{smallpmatrix} +
\begin{smallpmatrix} 0\\0\\ 1
\end{smallpmatrix}$.
It is clear that $\Gamma$ is discrete, acts freely properly and discontinuously on $\RR^3$, giving a quotient manifold $\Gamma \backslash \RR^3$ diffeomorphic to the hyperbolic torus bundle ${\mathbb T}_A^3$.


We see $A$ as a linear transformation of $\Span(e_1,e_2)$. This transformation is of the form $(u,v) \mapsto (\lambda u, \lambda^{-1}v)$ in suitable coordinates $(u,v)$. For any smooth function $a: \RR \to (0,\infty)$, which is $1$-periodic, the group $\Gamma$ acts isometrically for the Lorentz metric
\[
g_a=\dt^2+2a(t)\du\dv
\]
on $\RR^3$. Hence the metric $g_a$ induces a Lorentz metric $\overline{g}_a$ on $M={\mathbb T}_A^3$.

The flow of translations $T_{e_3}^t$ acts on ${\mathbb T}_A^3$ as an Anosov flow. When $a$ is a constant, the metric $\overline{g}_a$ is flat, and up to finite index, the isometry group of $({\mathbb T}_A^3,\overline{g}_a)$ coincides with this flow. It is thus noncompact.

In the general case of a $1$-periodic function $a: \RR \to (0,\infty)$, the linear transformation $
\begin{smallpmatrix} A & 0\\ 0 & 1
\end{smallpmatrix}$ induces an isometry $f$ of
$({\mathbb T}_A^3,\overline{g}_a)$ which preserves individually the Lorentz tori $t=t_0$ on ${\mathbb T}_A^3$, and acts on them by an Anosov diffeomorphism.

Interesting examples arise if one imposes a genericity condition on the function $t \mapsto a(t)$.


\begin{lemm}\label{lem.calcul.killing}
Assume that the function $a: \RR \to (0,\infty)$ is $1$-periodic, and that there is no sub-interval of $\RR$ where it takes the form $a(t)=Ae^{Bt}$, for some real numbers $A$ and $B$. Then all Killing fields of $g_a$ are tangent to the hyperplanes $t=t_0$. In particular, there is no nonempty open subset where the metric $g_a$ (resp. $\overline{g}_a$) is locally homogeneous. Moreover, the isometry group $\Iso({\mathbb T}_A^3,\overline{g}_a)$ virtually coincides with the subgroup $\langle f \rangle \simeq \ZZ$ generated by $f$. It is thus infinite discrete.
\end{lemm}

\begin{proof}
Let us consider a local Killing field $T$ for $g_a$, defined on some open subset $U \subset \RR^3$ that we may assume to be a product of open intervals. We write $T=\alpha \partial_t + \beta \partial_u + \gamma \partial_v$, where $\alpha, \beta, \gamma$ are smooth functions defined on $U$. The Lie derivative $L_Tg$ vanishes identically, what can be written:
\begin{equation}\label{eq.kil1}
0=L_Tg(\partial_i,\partial_j)=T.g(\partial_i,\partial_j)+g([\partial_i,T],\partial_j)+ g([\partial_j,T],\partial_i)
\end{equation}

Equation~\eqref{eq.kil1} when the pair $(i,j)$ is equal to $(t,t)$, $(u,u)$ and $(v,v)$ respectively leads to:
\begin{equation}\label{eq.kil2}
\frac{\partial \alpha}{\partial t}=0, \quad \frac{\partial \gamma}{\partial u}=0, \quad \frac{\partial \beta}{\partial v}=0.
\end{equation}

Equation~\eqref{eq.kil1} for the pair $(u,v)$ yields:
\begin{equation}\label{eq.kil3}
2 \alpha a'(t)+2 a(t)\left(\frac{\partial \beta}{\partial u}+ \frac{\partial \gamma}{\partial v}\right)=0.
\end{equation}

Finally, the pairs $(t,u)$ and $(t,v)$ lead to:
\begin{equation}\label{eq.kil4}
2 a(t) \frac{\partial \gamma}{\partial t}+ \frac{\partial \alpha}{\partial u}=0
\end{equation}
and
\begin{equation}\label{eq.kil5}
2 a(t) \frac{\partial \beta}{\partial t}+ \frac{\partial \alpha}{\partial v}=0.
\end{equation}

Deriving~\eqref{eq.kil4} with respect to $u$ and~\eqref{eq.kil5} with respect to $v$, we find $\frac{\partial^2 \alpha}{\partial u^2}=0$ and $\frac{\partial^2 \alpha}{\partial v^2}=0$. This leads to $\alpha(u,v)=\alpha_1 u + \alpha_2 v+ \alpha_3$, for some real numbers $\alpha_1, \alpha_2, \alpha_3$. We can now integrate equations~\eqref{eq.kil4} and~\eqref{eq.kil5}. We find $\gamma(t,v)=A_1(t)+B(v)$ and $\beta(t,u)=A_2(t)+C(u)$ for some functions $A_1$, $A_2$, $B$ and $C$. Plugging those expressions into~\eqref{eq.kil3}, we end up with $C'(u)+B'(v)=-\frac{a'(t)}{a(t)}\alpha(u,v)$. Under our assumption that $\frac{a'(t)}{a(t)}$ is constant on no sub-interval, this forces $\alpha$ to be identically zero. The Killing field $T$ is tangent to the hyperplanes $t=t_0$, as announced.

Let us now determine $\Iso({\mathbb T}_A^3,\overline{g}_a)$. The $1$-periodic function $a: \RR \to (0,\infty)$ induces a smooth function $\overline{a}: \Sop^1 \to (0,\infty)$. The value $\overline{a}(t_0)$ has the following geometric meaning. If $F_{t_0}$ denotes the fiber of $t_0 \in \Sop^1$ in the fibration ${\mathbb T}_A^3 \to \Sop^1$, then $\overline{a}(t_0)=\lambda_{\Gamma}\vol(F_{t_0})$, where $\vol(F_{t_0})$ is the Lorentz volume of $F_{t_0}$ and $\lambda_{\Gamma}$ is a positive constant depending only on $\Gamma$. It follows that $\Iso({\mathbb T}_A^3,\overline{g}_a)$ leaves invariant the fibers of $\overline{a}$. Since $\overline{a}$ is not constant, there exists a finite fiber for $\overline{a}$, so that a finite index subgroup of $\Iso({\mathbb T}_A^3,\overline{g}_a)$ leaves invariant a Lorentz torus $F_{t_0}$. We now make two extra observations. The first is that $\langle f \rangle$ has finite index in the isometry group of $F_{t_0}$. The second is and that the subgroup of $\Iso({\mathbb T}_A^3,\overline{g}_a)$ fixing $F_{t_0}$ pointwise is finite (this is just because a lorentz isometry is completely determines by its $1$-jet at a point, and that the subgroup of $\Omp(1,n-1)$ fixing pointwise a Lorentz hyperplane is finite). Those remarks show that $\langle f \rangle$ has finite index in $\Iso({\mathbb T}_A^3,\overline{g}_a)$.
\end{proof}



The examples described in the previous Lemma~\ref{lem.calcul.killing} prove Theorem~\ref{thm.nonhomogeneous} for hyperbolic torus bundles.


\subsection{Examples on parabolic torus bundles}\label{sec.exemplesparaboliques0}

\subsubsection{Flat, or non locally homogeneous examples}\label{sec.exemplesparaboliques}
We consider now $\RR^3$ with coordinates $(u,t,v)$. Let us call $H$ the $3$-dimensional Lie group given by the affine transformations
\[
\begin{pmatrix} 1& z & -\frac{\;\;z^2}{2}\\ 0 & 1 & -z \\ 0& 0 & \;\;\;1
\end{pmatrix}
+
\begin{pmatrix} r\\s\\ z
\end{pmatrix}
\]
where $r,s,z$ describe $\RR$. Observe that $H$ is a subgroup isomorphic to the $3$-dimensional Heisenberg group $\Heis$. The action of $H$ on $\RR^3$ is free and transitive. Observe also that $H$ acts isometrically for the flat Lorentz metric $h_{\flatm}=\dt^2+2\du\dv$. Let $a:\RR \to (0,\infty)$ be a smooth function, which is $1$-periodic, and let us consider the metric
\[
h_a=a(v)(\dt^2+2\du\dv).
\]
When $a$ is not constant, it is no longer true that $h_a$ is $H$-invariant. But it remains true that $h_a$ is invariant under the action of the discrete subgroup $\Gamma \subset H$, comprising transformations of the form
\[
\begin{pmatrix} 1& m & -\frac{\;\;m^2}{2}\\ 0 & 1 & -m \\ 0& 0 & \;\;1
\end{pmatrix}
+ \begin{pmatrix} \frac{n}{2}\\\frac{l}{2}\\ m
\end{pmatrix}
\]
where $m,n,l$ describe $\ZZ$. The gluing map between planes $v=0$ and $v=1$ is made by the matrix $A= 
\begin{smallpmatrix} 1& 1\\0 & 1
\end{smallpmatrix}$.
Thus the quotient $\Gamma \backslash \RR^3 $ is diffeomorphic to $\TT_A^3$, with $A$ the unipotent matrix above. Torus bundles $\TT_A^3$ are characterized up to homeomorphism by the conjugacy class in $\SL(2,\ZZ)$ of the gluing matrix $A$. It follows that all parabolic torus bundles are obtained for gluing maps of the form $A_k= \begin{smallpmatrix} 1& k\\0 & 1
\end{smallpmatrix}$, $k \in \NN^*$, hence
by considering finite index subgroups of $\Gamma$.

The metric $h_a$ induces a Lorentz metric $\overline{h}_a$ on the parabolic torus bundle $\TT_A^3$, and the linear maps $ B=
\begin{smallpmatrix} 1& m & -\sfrac{m^2}{2}\\ 0 & 1 & -m \\ 0& 0 & \;\;1
\end{smallpmatrix}$, $m \in \ZZ$, normalize $\Gamma$, hence induce a group of isometries in $(\TT_A^3,\overline{h}_a)$. It is readily checked that this group does not have compact closure in $\Iso(\TT_A^3, \overline{h}_a)$.



We now make the following observation. Let $X=X_1 \frac{\partial}{\partial u} +X_2 \frac{\partial}{\partial t}+X_3 \frac{\partial}{\partial v}$ be a local conformal Killing field for the flat metric $h_{\flatm}$ (namely the local flow of $X$ preserves the conformal class of $h_{\flatm}$). Then $L_Xh_{\flatm}=\alpha_X h_{\flatm}$ for a smooth function $\alpha_X$. Assume that $X_3$ is nonzero on a small open set, then $X$ will be a Killing field for $h_a$ if and only if
\begin{equation}\label{equa.kill}
\frac{a'(v)}{a(v)}=-\frac{\alpha_X(x)}{X_3(x)}
\end{equation}

Because we are in dimension $>2$, the set of local conformal Killing fields for $h_{\flatm}$ is finite dimensional, hence for a generic choice of smooth, $1$-periodic $a$, the relation~\eqref{equa.kill} will not be satisfied, whatever the conformal Killing field $X$ we are considering. It follows that for such a generic set of functions, there will not be any open subset of $\TT_A^3$ (resp. of $\TT^3$) where the metric $\overline{h}_a$ will be locally homogeneous.

These examples prove Theorem~\ref{thm.nonhomogeneous} for parabolic torus bundles.

\subsubsection{Examples modelled on Lorentz--Heisenberg geometry}\label{sec.lorentzheisenberg}
We denote by $\heis$ the $3$-dimensional Heisenberg Lie algebra, and $\Heis$ the connected, simply connected, associated Lie group. Recall that $\heis$ admits a basis $X,Y,Z$, for which the only nontrivial bracket relation is $[X,Y]=Z$. Let $B \in \SL(2,\ZZ)$ be a hyperbolic matrix, and consider the automorphism $\varphi$ of $\heis$, which in the basis $X,Y,Z$ writes $\begin{smallpmatrix} B & 0\\ 0 & 1\end{smallpmatrix} $. It defines an automorphism $\Phi$ of the Lie group
$\Heis$.

The matrix $B$ is diagonal in some basis $X',Y'$ of $\Span(X,Y)$, with eigenvalues $\lambda,\lambda^{-1}$. The Lorentz scalar product defined by $\langle X',Y'\rangle=1$, $\langle Z,Z\rangle=1$, and all other products are zero, can be left-translated on $\Heis$ to give an homogeneous Lorentz metric $g_{LH}$ called \emph{the Lorentz--Heisenberg} metric on $\Heis$. The reader will find more details and further references about this geometry in~\cite[Section~4.1]{sorin.ghani}. One can actually show that different choices of the hyperbolic matrix $B$ will produce isometric spaces.

By construction, $\Phi$ acts isometrically on $(\Heis,g_{LH})$, and so do left translations. It is explained in~\cite[Section~4.1]{sorin.ghani} that the identity component $\Iso^o(\Heis, g_{LH})$ is $4$-dimensional isomorphic to $\RR \ltimes \Heis$. The $\RR$-factor corresponds a $1$-parameter group of automorphisms of $\heis$ containing $\Phi$.

We now consider the following lattice in $\Heis$:
\[
H_{\ZZ}:=\{ \exp(aX+bY+cZ) \, | \, (a,b,c) \in \ZZ^3 \}.
\]
The quotient $H_{\ZZ} \backslash \Heis$ is homeomorphic to a parabolic torus bundle, on which the Lorentz--Heisenberg metric induces a metric $\overline{g}_{LH}$. The automorphism $\Phi$ preserves $H_{\ZZ}$, hence induces an isometry $\overline{\Phi}$ on $(H_{\ZZ} \backslash \Heis,\overline{g}_{LH})$, and $\overline{\Phi}$ generates a noncompact group. As in Section~\ref{sec.exemplesparaboliques}, those examples arise on a parabolic torus bundle. Notice however that they are geometrically different from examples of Section~\ref{sec.exemplesparaboliques}, which were not locally homogeneous.

\subsection{Some examples on the \texorpdfstring{$3$}{3}-torus \texorpdfstring{$\TT^3$}{TT3}}\label{sec.exemplestorus}



\subsubsection{A flat example}\label{sec.exemplestorus1}
The most classical example, already mentioned in the introduction, comes from the flat metric
\[
g_{\flatm}=-\du^2+\dv^2+\dw^2.
\]

We call $\Omp(1,2)$ the group of linear transformations preserving $g_{\flatm}$, and we introduce $\Gamma$ the discrete subgroup generated by the translations $T_u,T_v,T_w$ of vectors $u,v,w$. The quotient $\Gamma \backslash \RR^3$ inherits an induced metric $\overline{g}_{\flatm}$ from $g_{\flatm}$, and the isometry group of $(\TT^3,\overline{g}_0)$ is $\Omp(1,2)_{\ZZ} \ltimes \TT^3$. Because the quadratic form $-u^2+v^2+w^2$ has rational coefficients, a theorem of Borel and Harish--Chandra ensures that $\Omp(1,2)_{\ZZ}$ is a lattice in $\Omp(1,2)$. In particular, $\Omp(1,2)_{\ZZ} \ltimes \TT^3$ is noncompact. The identity component of the isometry group is however compact in this case.

\subsubsection{Non locally homogeneous examples}\label{sec.exemplestorus2}

These examples are built in the same way as those of Sections~\ref{sec.exempleshyperboliques} and~\ref{sec.exemplesparaboliques}, so that we will be rather sketchy in our description.

We consider the metric $g_a$, introduced in Section~\ref{sec.exempleshyperboliques}, for $a: \RR \to (0,\infty)$ a smooth $1$-periodic function.

The metric $g_a$ is invariant by the discrete group $\Gamma$ generated by the translations of vectors $e_1,e_2$ and $e_t$. Hence $g_a$ induces a metric $\overline{g}_a$ on $\TT^3$. As in Section~\ref{sec.exempleshyperboliques}, for generic choices of the function $a:\RR \to (0,\infty)$, there is no open set on which $\overline{g}_a$ is locally homogeneous. The isometry group is then $\ZZ \ltimes \TT^2$ (the $\ZZ$-factor comes from the transformation $
\begin{smallpmatrix} A & 0\\ 0 & 1
\end{smallpmatrix}$, as in~\ref{sec.exempleshyperboliques}).


We can also consider the metric $h_a$ introduced in Section~\ref{sec.exemplesparaboliques}, and take for $\Gamma$ the discrete subgroup generated by the translations of vectors $(e_u,e_t,e_v)$. This yields a metric $\overline{h}_a$ on $\TT^3=\Gamma \backslash \RR^3$. For generic choices of the $1$-periodic function $a:\RR \to (0,\infty)$, there is no open set on which $\overline{h}_a$ is locally homogeneous, and the isometry group is noncompact, isomorphic to $\ZZ \ltimes \TT^2$.

These examples prove Theorem~\ref{thm.nonhomogeneous} for $3$-dimensional tori.


\section{Curvature, recurrence, and local Killing fields}\label{sec.gromov}
\subsection{Generalized curvature map and integrability locus}\label{sec.curvaturemap}

Let us consider $(M,g)$, a smooth Lorentz manifold of dimension $n \geq 2$. All the material presented below holds actually in the much wider framework of Cartan geometries, but we will not need such a generality.

\subsubsection{Cartan connection associated to the metric}\label{sec.cartan.connection}
Let $\pi: \hm \to M$ denote the bundle of orthonormal frames on $\hm$. This is a principal $\Omp(1,n-1)$-bundle over $M$, and it is classical (see~\cite[Chapter~IV.2]{kobanomi}) that the Levi-Civita connection associated to $g$ can be interpreted as an Ehresmann connection $\alpha$ on $\hm$, with values in the Lie algebra $\oo(1,n-1)$. For the reader's convenience, we briefly recall the link between the two points of view. The kernel of the form $\alpha$ determines a distribution ${\mathcal H}$ on $\hm$, which is transverse to the fibers and $\Omp(1,n-1)$-invariant. Let us consider a curve $\gamma: [0,1] \to M $, and a frame at $x=\gamma(0)$, that we see as a point $\hat{x} \in \hm$. There is a unique lift $\hat{\gamma}$ of $\gamma$ to $\hm$, which starts at $\hx$ and is tangent to ${\mathcal H}$. This curve $t \mapsto \hat{\gamma}(t)$ describes the family of frames obtained by parallel transporting $\hx$ along $\gamma$, for the Levi-Civita connection.
%Levi-Civita est une seule personne avec un nom composé (comme Franco-Nollet)


Let now $\theta$ be the soldering form on $\hm$, namely the $\RR^n$-valued $1$-form on $\hm$, which to every $\xi \in T_{\hx}\hm$ associates the coordinates of the vector $\pi_*(\xi) \in T_xM$ in the frame $\hx$. The sum $\alpha + \theta$ is a $1$-form $\omega: T\hm \to \oo(1,n-1) \ltimes \RR^n$ called the \emph{canonical Cartan connection} associated to $(M,g)$ (see~\cite[Chapter~6]{sharpe} for a nice introduction to Cartan geometries).

In the following, we will denote by $\lieg$ the Lie algebra $\oo(1,n-1) \ltimes \RR^n$. The Cartan connection $\omega$ satisfies the two crucial properties:
\begin{itemize}
\item \label{list3.1.1.1}For every $\hx \in \hm$, $\omega_{\hx}: T_{\hx}\hm \to \lieg$ is an isomorphism of vector spaces.
\item \label{list3.1.1.2}The form $\omega$ is $\Omp(1,n-1)$-equivariant (where $\Omp(1,n-1)$ acts on $\lieg$ via the adjoint action).
\end{itemize}

\subsubsection{Generalized curvature map}\label{sec.generalized.curvature}
The \emph{curvature} of the Cartan connection $\omega$ is a $2$-form $K$ on $\hm$, with values in $\lieg$. If $\hat{\widehat{X}}$ and $\hat{Y}$ are two tangent vectors at a same point of $\hm$, it is given by the relation:
\[
K(\hat{X},\hat{Y})=\dd \omega(\hat{X},\hat{Y}) +[\omega(\hat{X}),\omega(\hat{Y})].
\]
Because $\omega_{\hx}$ establishes an isomorphism between $T_{\hx}\hm$ and $\lieg$ at each point $\hx$ of $\hm$, $\omega$ provides a trivialization of the tangent bundle $T\hm=\hm \times \lieg$. It follows that any field of $k$-linear forms on $\hm$, with values in some vector space ${\mathcal W}$, can be seen as a map from $\hm$ to $ \Hom(\otimes^k \lieg,{\mathcal W})$. This remark applies in particular for the curvature form, yielding a curvature map $\kappa: \hm \to {\mathcal W}_0$, where the vector space ${\mathcal W}_0$ is $\Hom(\wedge^2(\RR^n);\lieg)$ (the curvature is antisymmetric and vanishes when one argument is tangent to the fibers of $\hm$).

We can now differentiate $\kappa$, getting a map $D \kappa : T\hm \to {\mathcal W}_0$. Our previous remark allows us to see $D\kappa$ as a map $D \kappa : \hm \to {\mathcal W}_1$, with ${\mathcal W}_1=\Hom(\lieg,{\mathcal W}_0)$. Applying this procedure $r$ times, we define inductively the $r$-derivative of the curvature $D^r \kappa : \hm \to \Hom(\otimes^r \lieg,{\mathcal W}_r)$ (with ${\mathcal W}_r$ defined inductively by ${\mathcal W}_r=\Hom(\lieg, {\mathcal W}_{r-1})$).

Let us now set $m = \dim \Omp(1,n-1)=\frac{n(n-1)}{2}$. The \emph{generalized curvature map} of $(M, g)$ is the map $\dk =(\kappa,D\kappa, \ldots,D^{m}\kappa)$. The $\Omp(1,n-1)$-module ${\mathcal W}_{1} \oplus \ldots \oplus {\mathcal W}_{m}$ will be simply denoted ${\mathcal W}$ in the following.



\subsubsection{Integrability locus}\label{sec.integrability.locus}
One defines \emph{the integrability locus of $\hm$}, denoted $\hmint$, as the set of points $\hx \in \hm$ at which the rank of $ \dk$ is locally constant. Notice that $ \hmint$ is a $\Omp(1,n-1)$-invariant open subset of $\hm$. Because the rank can only increase locally, this open subset is dense. We define also $ \mint \subset M$, the integrability locus of $M$, as the projection of $\hmint$ on $M$. This is a dense open subset of $M$.



\subsection{Integrability theorem and structure of \texorpdfstring{$\Is^{\loc}$}{Isloc}-orbits}
\label{sec.integration}


For $x \in M$, the $\Is^{\loc}$-orbit of $x$ is the set of points $y \in M$ such that $y=f(x)$ for some local isometry $f:U\subset M \to V \subset M$. The $\kiloc$-orbit of $x$ is the set of points $y \in M$ that can be reached by flowing along (finitely many) successive local Killing fields.

Local flows of isometries on $M$ clearly induce local flows on the bundle of orthonormal frames, which moreover preserve $\omega$. It follows that any local Killing field $X$ on $U \subset M$ lifts to a vector field $\hat{X}$ on $\hm$, satisfying $L_{\hat{X}}\omega=0$. Conversely, local vector fields of $\hm$ such that $L_{\hat{X}}\omega=0$, that we will henceforth call \emph{$\omega$-Killing fields}, commute with the right $\Omp(1,n-1)$-action on $\hm$. Hence, they induce local vector fields $X$ on $M$, which are Killing because their local flow maps orthonormal frames to orthonormal frames. It is easily checked that a $\omega$-vector field which is everywhere tangent to the fibers of the bundle $\hm \to M$ must be trivial. As a consequence, there is a one-to-one correspondence between local $\omega$-Killing fields on $\hm$ and local Killing fields on $M$. We will use this correspondence all along the paper. The same remark holds for local isometries.


Observe finally that if $\hat{X}$ is a $\omega$-Killing field on $\hm$ (namely $L_{\hat{X}} \omega=0$), then the local flow of $\hat{X}$ preserves $\dk$, hence $\hat{X}$ belongs to $\Ker(D_{\hx}\dk)$ at each point. The integrability theorem below says that the converse is true on $\hmint$.


\begin{theo}[Integrability theorem]\label{thm.integrabilite}
Let $(M^n,g)$ be a Lorentz manifold, and $\hmint \subset \hm$ the integrability locus.
\begin{enumerate}
\item \label{list3.1.1}For every $\hx \in \hmint$, and every $\xi \in Ker(D_{\hx}\dk)$, there exists a local $\omega$-Killing field $\hat{X}$ around $\hx$ such that $\hat{X}(\hx)=\xi$.
\item \label{list3.1.2}The $\Is^{\loc}$-orbits in $\mint$ are submanifolds of $\mint$, the connected components of which are $\kiloc$-orbits.
\end{enumerate}
\end{theo}


The deepest, and most difficult part, of the theorem is the first point. Such an integrability result as well as the structure of $\Is^{\loc}$-orbits first appeared in~\cite{gromov}. The results were recast in the framework of Cartan geometry by K.~Melnick in the analytic case (see~\cite{melnick}). The reference~\cite{vincent} gives an alterative approach for smooth Cartan geometries, leading to the statement of Theorem~\ref{thm.integrabilite}. A proof that the integrability property actually holds on the set where the rank of $\dk$ is locally constant (first point of the theorem) can be found in Appendix~A of~\cite{frances.open}.


Let us recall how the second point of Theorem~\ref{thm.integrabilite} easily follows from the first one (see also~\cite[Section~4.3.2]{vincent}). The generalized curvature map $\dk: \hm \to {\mathcal W}$ is invariant under all local isometries. It follows that $\hmint$ is invariant as well. Given $\hx \in \hmint$, and $w=\dk(\hx)$, the $\Is^{\loc}$-orbit $\Is^{\loc}(\hx)$ is contained in $\dk^{-1}(w) \cap \hmint$. Now since $\dk$ has locally constant rank on $\hmint$, $\dk^{-1}(w) \cap \hmint$ is a submanifold of $\hmint$, and the first point of Theorem~\ref{thm.integrabilite} exactly means that the $\kiloc$-orbit of $\hx$ coincides with the connected component of $\dk^{-1}(w) \cap \hmint$ containing $\hx$, hence is a submanifold on $\hmint$. The set $\Is^{\loc}(\hx)$ is a union of such connected components, hence a submanifold too. The point we have to check is that this property remains true when one projects everything on $M$. Observe first that the projection of $\dk^{-1}(w) \cap \hmint$ on $M$ coincides with that of $\dk^{-1}({\calo }.w) \cap \hmint$, where ${\calo}.w$ stands for the $\Omp(1,n-1)$-orbit of $w$ in ${\mathcal W}$. Now, using the constancy of $\rank(\dk)$ on $ \dk^{-1}({\calo }.w) \cap \hmint$, the $\Omp(1,n-1)$-equivariance of $\dk$, and the fact that $\Omp(1,n-1)$-orbits in ${\mathcal W}$ are locally closed, one shows that $\dk^{-1}({\calo }.w) \cap \hmint$ is a submanifold of $\hmint$. By $\Omp(1,n-1)$-invariance of this set, its projection on $\mint$ is a submanifold too.


\subsection{Components of the integrability locus and \texorpdfstring{$\kiloc$}{kiloc}-algebra}\label{sec.intlocus}
Let us recall here classical facts about the behaviour of Killing fields on the integrability locus $\mint$ (see~\cite[Section~5.15]{dambragro} for a general discussion in the framework of rigid geometric structures).

For each $x \in M$, let us consider a sequence $U_i$ of nested connected open neighbourhoods of $x$, such that $\{ x \}=\bigcap_{i \in \NN}U_i$. For each $i \in \NN$, denote by $\kil(U_i)$ the Lie algebra of Killing fields defined on $U_i$. A Killing field of $U_i$ vanishing on an open subset must be identically zero, so that the restriction maps yield a sequence of Lie algebra embeddings $\kil(U_i) \to \kil(U_{i+1})$. The dimension of each $\kil(U_i)$ is finite (bounded by $\frac{n(n+1)}{2}$) and is nondecreasing with $i$, hence it stabilizes for $i \geq i_0$. In other words, all Killing fields defined on $U_i$ for $i \geq i_0$ are restrictions of Killing fields of $U_{i_0}$. We can thus state:
\begin{fact}\label{fact.extension}
Every point $x \in M$ admits an open connected neighbourhood $U(x)$ such that any Killing field defined on some connected open set $V$ containing $x$, can be extended to $U(x)$.
\end{fact}

As a consequence, there is a good notion of Lie algebra of local Killing fields at $x$ that we denote by $\kiloc(x)$: It coincides with $\mathfrak{kill}(U(x))$, where the open set $U(x)$ is given by the previous fact. Observe that for every $y \in U(x)$ there is a natural embedding of Lie algebras $\kiloc(x) \to \kiloc(y)$. It is obtained by restricting Killing fields on $U(x)$ to a small neighbourhood of $y$.


\subsubsection{Components and analytic continuation}\label{sec.analytic}
The integrability locus $\mint$ splits into a union of connected components $\bigcup {\mathcal M}_i$. The ${\mathcal M}_i's$ will be just called \emph{components} in the sequel. If $x$ belongs to a component $\calm$, then the open set $U(x)$ given by Fact~\ref{fact.extension} will be chosen to be included in $\calm$. We are going to see that Killing fields on such components behave as if the structure was analytic.


Let $x \in \mint$ and denote by $\calm$ be the component containing $x$. It follows from Theorem~\ref{thm.integrabilite} that the dimension of the Lie algebra $\kiloc(x)$ is equal to $\dim \hm - \rk(\dk)$. Hence the dimension of $\kiloc(x)$ does not depend on the point $x \in \calm$. We already observed that for $y \in U(x)$, there is a Lie algebra embedding $\kiloc(x) \to \kiloc(y)$. Since the two dimensions coincide this embedding is actually an isomorphism. In other words, any Killing field defined on a connected open subset $V \subset U(x)$ can be extended to a Killing field defined on $U(x)$ (recall that $U(x) \subset \calm$). This extension result is similar to the one obtained by Amores in~\cite{amores}. It allows to perform analytic continuation of local Killing fields along paths contained in $\mint$. As in~\cite{amores}, this leads to the:


\begin{fact}\label{fact.extensionbis}
Let $\calm$ be a component of the integrability locus $\mint$, and $U \subset \calm$ be a connected, simply connected open subset. Then every Killing field defined on a connected open subset of $U$ can be extended to a Killing field on $U$.
\end{fact}

Observe that by the discussion above, the isomorphism type of $\kiloc(x)$ does not depend on $x \in \calm$, and we will sometimes write $\kiloc({\mathcal M})$ instead of $\kiloc(x)$.


\subsubsection{Isotropy algebra}\label{sec.isotropy.algebra}
For $x \in M$, we consider $\mathfrak{Is}(x)$, the isotropy algebra at $x$, namely the Lie algebra of local Killing fields defined in a neighbourhood of $x$ and vanishing at $x$.

\begin{fact}\label{rk.isotropie}
If $x \in \mint$, then the isotropy algebra $\mathfrak{Is}(x)$ is isomorphic to the Lie algebra of the stabilizer of $\dk(x)$ in $\Omp(1,n-1)$.
\end{fact}


\begin{proof}
Let us consider $\hx \in \hm$ in the fiber of $x$. Every local Killing field $X$ around $x$ which vanishes at $x$, lifts to a local $\omega$-Killing field around $\hx$, denoted $\hat{X}$, which is vertical (namely tangent to the fiber) at $\hx$. We call $ev_{\hx}$ the map $\hat{X} \mapsto \omega(\hat{X}(\hx))$. The relation $\varphi_{\hat{X}}^t.\, \hx=\hx.e^{tev_{\hx}({\hat{X}})}$, available for $t$ in a neighbourhood of $0$, together with the invariance of $\dk$ under $\omega$-Killing flows, shows that $ev_{\hx}$ is a linear embedding from $\mathfrak{Is}(x)$ to the Lie algebra of the stabilizer of $\dk(x)$ in $\Omp(1,n-1)$ (this map is one-to-one because a local $\omega$-Killing field on $\hm$ vanishing at a point must be identically zero). Cartan's formula $L_{\hat{X}}=\iota_{\hat{X}} \circ d +d \circ \iota_{\hat{X}}$ shows that whenever $\hat{X}$ and $\hat{Y}$ are two $\omega$-Killing fields around $\hx$, the relation $\omega([\hat{X},\hat{Y}])=K(\hat{X},\hat{Y})-[\omega(\hat{X}),\omega(\hat{Y})]$ holds. When $\hat{X}$ or $\hat{Y}$ is vertical, $K(\hat{X},\hat{Y})=0$, proving that $ev_{\hx}$ is an anti-morphism of Lie algebras. To see that $ev_{\hx}$ is onto, let us consider $\{ e^{t \xi} \}_{t \in \RR}$, a $1$-parameter group of $\Omp(1,n-1)$ fixing $\dk(\hx)$. Clearly, $\xi$ belongs to $\Ker(D_{\hx}\dk)$, so that by Theorem~\ref{thm.integrabilite}, $\omega^{-1}(\xi)$ is the evaluation at $\hx$ of a local $\omega$-Killing field. This $\omega$-Killing field being vertical at $\hx$, its projection yields a local Killing field of ${\mathfrak{Is}}(x)$.
\end{proof}



\subsection{Nontrivial recurrence provides nontrivial Killing fields}\label{sec.recurrence}
We still deal here with $(M^n,g)$ a closed $n$-dimensional Lorentz manifold ($n \geq 2$). There is, as in Riemannian geometry, a notion of Lorentzian volume, which provides a smooth, $\Iso(M^n,g)$-invariant measure on $M$. This measure is finite under our assumption that $M$ is closed. When the group $\Iso(M^n,g)$ is noncompact, Poincaré's recurrence theorem applies and almost every point of $M$ is recurrent for the action of $\Iso(M^n,g)$. Recall that a point $x$ is said to be recurrent when there exists a sequence of isometries $(f_k)$ leaving every compact subset of $\Iso(M^n,g)$, and such that $f_k(x)$ converges to $x$. We are going to see that such a recurrence phenomenon is responsible for the existence of nontrivial continuous local symmetries. The precise statement is:

\begin{prop}\label{prop.dimension}
Let $(M^n,g)$ be a closed, $n$-dimensional Lorentz manifold, and assume that $\Iso(M^n,g)$ is noncompact. Then
\begin{enumerate}
\item \label{list3.5.1}For every $x \in \mint$, the isotropy algebra $\mathfrak{Is}(x)$ generates a noncompact subgroup of ${\operatorname O}(T_xM)$.
\item \label{list3.5.2}For every component ${\mathcal M} \subset \mint$, the Lie algebra $\kiloc({\mathcal M})$ is at least $3$-dimensional.
\end{enumerate}
\end{prop}


\begin{proof}
The proof of the first point is already contained in~\cite[Proposition~5.1]{frances.open}. We summarize here the main arguments for the reader's convenience. Let $x$ be a recurrent point for $\Iso(M^n,g)$ and choose $\hx \in \hm$ in the fiber of $x$. The recurrence hypothesis means that there exists $(f_k)$ tending to infinity in $\Iso(M^n,g)$, and $(p_k)$ a sequence of ${\operatorname O}(1,n-1)$ such that $f_k(\hx).p_k^{-1}$ tends to $\hx$. By equivariance of the generalized curvature map $\dk: \hm \to \mathcal{W}$, we also have
\[
p_k.\dk(\hx) \to \dk(\hx).
\]
Observe that $(p_k)$ tends to infinity in ${\operatorname O}(1,n-1)$, because $\Iso(M^n,g)$ acts properly on $\hm$.

The ${\operatorname O}(1,n-1)$-orbits on $\mathcal{W}$ are locally closed, because the action of ${\operatorname O}(1,n-1)$ is linear, hence algebraic. As a consequence, there exists a sequence $(\epsilon_k)$ in ${\operatorname O}(1,n-1)$ with $\epsilon_k \to \id$ and $\epsilon_k.p_k.\dk(\hx) = \dk(\hx)$. Since $(p_k)$ tends to infinity by properness of the action of $\Iso(M^n,g)$ on $\hm$, so does $(\epsilon_k.p_k)$, proving that the stabilizer $I_{\hx}$ of $\dk(\hx)$ in ${\operatorname O}(1,n-1)$ is noncompact. This group is algebraic, hence the identity component $I_{\hx}^{o}$ is noncompact too. Fact~\ref{rk.isotropie} then ensures that $\mathfrak{Is}(x)$ generates a noncompact subgroup of ${\operatorname O}(T_xM)$ (under the identification of $\mathfrak{Is}(x)$ with a subalgebra of $\mathfrak{o}(T_xM)$ under the isotropy representation), for every recurrent point $x \in \mint$. The property is thus true everywhere on $\mint$, by density of recurrent points on $M$.

To prove the second point, we start with $x \in \mint$, and consider a connected, simply connected neighbourhood $U \subset \mint$ of $x$. By Fact~\ref{fact.extensionbis}, every algebra $\mathfrak{Is}(y)$, $y \in U$, is realized as a Lie algebra of Killing fields defined on $U$. The first point of proposition~\ref{prop.dimension} says that there exists $X$ a nontrivial Killing field on $U$, such that $X(x)=0$. The zero locus of $X$ is a nowhere dense set in $U$ (actually a submanifold of codimension $\geq 1$). We can thus pick $y \in U$ satisfying $X(y) \not =0$, and apply again the first point of the proof at $y$. We get a second nontrivial Killing field $Y$ defined on $U$, and vanishing at $y$. Again, the zero locus of $Y$ is a submanifold of codimension $\geq 1$. The set of vectors in $T_yM$ which are transverse to this submanifold is thus open in $T_yM$, so that we can pick a vector $u \in T_yM$ such that $g_y(u,X(y)) \not =0$, and $u$ is transverse to the zero locus of $Y$. Let $t \mapsto \gamma_u(t)$ be the geodesic passing through $y$ at $t =0$ and such that $\dot{\gamma}_u(0)=u$. Because $X$ is a Killing field, the operator $\nabla X$ is antisymmetric for $g$.

This leads to \emph{Clairault's equation}: $\frac{\dd}{\dt}g(\dot{\gamma}_u(t),X(\gamma(t)))=0$. In particular $g(\dot{\gamma}_u,X)$ is constant on $\gamma_u$, hence does not vanish on $\gamma_u$ by our choice of $u$.

On the other hand, by the same Clairault's equation, $g_{\gamma_u}(\dot{\gamma}_u,Y) =0$ along $\gamma_u$. This implies that $X(\gamma(t))$ and $Y(\gamma_u(t))$ are linearly independent as soon as $Y(\gamma_u(t)) \not =0$, a property satisfied for small and nonzero values of $t$. We obtain an open subset of $U$ where the orbits of the local Killing algebra have dimension $\geq 2$, while for every point $z \in U$, the dimension of $\mathfrak{Is}(z)$ is $\geq 1$. The minoration $\dim \kiloc(U) \geq 3$~follows.
\end{proof}

\begin{rema}\label{rem.pseudoRiemannian}
The proof above does not use the fact that the metric $g$ is Lorentzian, and Proposition~\ref{prop.dimension} actually holds for any pseudo-Riemannian (non Riemannian) metric.
\end{rema}

\subsection{Components of the integrability locus, and their classification}
\label{sec.components}
We now stick to dimension~$3$, and we consider a closed Lorentz manifold $(M,g)$. We assume that the isometry group $\Iso(M,g)$ is noncompact.

On each component ${\mathcal M} \subset \mint$, the discussion of Section~\ref{sec.integration} shows that there is a well-defined Lie algebra $\kiloc({\mathcal M})$ of local Killing fields (beware that some monodromy phenomena may occur), which by Proposition~\ref{prop.dimension} is at least $3$ dimensional. The first point of Proposition~\ref{prop.dimension} ensures that the dimension of $\mathfrak{Is}(x)$ is at least $1$. It can not be $2$, because Fact~\ref{rk.isotropie} would then imply that a vector in $\mathcal{W}$ has a stabilizer of dimension~$2$ under the linear action of ${\operatorname O}(1,2)$ on $\mathcal{W}$. But since ${\operatorname O}(1,2)$ is locally isomorphic to $\SL(2,\RR)$, all finite dimensional linear representation of ${\operatorname O}(1,2)$ are easily described, and on checks that no vector can have a stabilizer of dimension exactly $2$. We conclude that the dimension of $\mathfrak{Is}(x)$ is $1$ or $3$ for every $x \in {\mathcal M}$, and the dimension of $\kiloc({\mathcal M})$ is thus $6$, $4$ or $3$.
\begin{itemize}
\item \label{list3.5.a} When this dimension is $6$, the component has constant sectional curvature. Indeed, at each point the $3$-dimensional isotropy acts transitively on the Grassmannian of Lorentzian (resp. Riemannian) $2$-planes.
\item \label{list3.5.b} When the dimension of $\kiloc({\mathcal M})$ is $4$, it is not hard to check that $\calm$ is locally homogeneous (see for instance~\cite[Lemma~4]{sorin.karin}). The dimension of $\mathfrak{Is}(x)$ is then $1$ at each point $x \in \calm$.
\item \label{list3.5.c} When the dimension of $\kiloc({\mathcal M})$ is $3$, then the dimension of $\mathfrak{Is}(x)$ is $1$ or $3$ at each point $x \in \calm$. The $\kiloc$-orbits have dimension~$0$ or $2$, and the component is nowhere locally homogeneous.
\end{itemize}


Proposition~\ref{prop.dimension} ensures that whenever the dimension of $\mathfrak{Is}(x)$ is $1$, then this algebra generates a hyperbolic or a parabolic flow in ${\operatorname O}(T_xM) \simeq {\operatorname O}(1,2)$. In the first case, we say that $x$ is \emph{a hyperbolic point}, and in the second one we call $x$ \emph{a parabolic point}.

\begin{defi}[Hyperbolic and parabolic components]
A component ${\mathcal M}$ of $\mint$ which is not of constant curvature is said to be hyperbolic when it contains a hyperbolic point. Otherwise, it is called parabolic.
\end{defi}

Observe that this definition allows a priori a hyperbolic component to contain parabolic points (it will turn out later that this does not occur).

To summarize, components of $\mint$ split into three (rough) categories.
\begin{enumerate}
\item \label{list3.7.1}The first category comprises all components having \emph{constant sectional\linebreak curvature}. 
\item \label{list3.7.2}The second category comprises \emph{hyperbolic components}. Those in turn split into two subcategories:
\begin{enumerate}\alphenumi
\item \label{list3.7.2a}The locally homogeneous ones, for which the dimension of $\kiloc({\mathcal M})$ is~$4$.
\item \label{list3.7.2b}The non locally homogeneous ones, for which the dimension of $\kiloc({\mathcal M})$ is $3$.
\end{enumerate}
\item \label{list3.7.3}The remaining components are \emph{parabolic}. They can also be split into:
\begin{enumerate}\alphenumi
\item \label{list3.7.3a}The locally homogeneous ones for which $\dim(\kiloc({\mathcal M}))=4$.
\item \label{list3.7.3b}The non locally homogeneous ones for which $\dim(\kiloc({\mathcal M}))=3$.
\end{enumerate}
\end{enumerate}

Let us notice that components $\calm$ with constant curvature are those for which $\kiloc(\calm)$ has dimension~$6$. Theorem~\ref{thm.integrabilite} says that on $\calm$, $\dim \kiloc(\calm)=\dim \hm - \rk(\dk)$. We infer that points belonging to components of constant curvature are those for which the rank of $\dk$ is locally equal to $0$. In the same way, points belonging to locally homogeneous components are those for which the rank of $\dk$ is locally equal to $2$. Finally, we prove:

\begin{lemm}\label{lem.rank3}
Points $x \in M$ belonging to a component of $\mint$ which is not locally homogeneous are exactly those at which the rank of $\dk$ is $3$.
\end{lemm}


\begin{proof}
Recall the generalized curvature map $\dk : \hm \to {\mathcal W}$ introduced in Section~\ref{sec.integration}. We saw in Proposition~\ref{prop.dimension} that for every $x \in \mint$, $\kiloc(x)$ has dimension $\geq 3$. Because $\dk$ is invariant along $\kiloc$-orbits in $\hm$, the corank of $\dk$ is a least $3$ on $\hmint$, and because $\hm$ has dimension~$6$, the rank of $\dk$ is at most $3$ on the dense set $\hmint$, hence on $\hm$. The rank can only increase locally, hence points where the rank of $\dk$ is $3$ actually stay in $\mint$.
\end{proof}



\section{Locally homogeneous Lorentz manifolds with noncompact isometry group}\label{sec.homogeneous}

In this section, we prove Theorem~\ref{thm.topological} in the case where all the components of $\mint$ are locally homogeneous, implying that $(M,g)$ is locally homogeneous on a dense open set. Our study will also settle the locally homogeneous case of Theorem~\ref{thm.geometrical}.

We observed in the previous section that if all components of $\mint$ are locally homogeneous, and under our standing assumption that $\Iso(M,g)$ is noncompact, the Lie algebra of local Killing vector fields has dimension $\geq 4$ on each component. We can apply the results of~\cite{frances.open}, saying that we must then have $\mint=M$, and the manifold $(M,g)$ is locally homogeneous.

\begin{theo}[{\cite[Theorem~B]{frances.open}}]\label{thm.quasihomogeneous}
Let $(M,g)$ be a smooth $3$-dimensional Lorentz manifold. Assume that on a dense open subset, the Lie algebra of local Killing fields is at least $4$-dimensional. Then $(M,g)$ is locally homogeneous.
\end{theo}

There are a lot of homogeneous $3$-dimensional models for Lorentz manifolds. Fortunately, very few of them can appear as the local geometry of a closed manifold with a noncompact isometry group. We have indeed:


\begin{theo}[{\cite[Theorem~2.1]{sorin.ghani}}]\label{thm.ghanisorin}
Let $(M,g)$ be a closed locally homogeneous Lorentz manifold. Assume that at some (and then at each) point $x \in M$ the isotropy algebra $\mathfrak{Is}(x)$ generates a noncompact subgroup of $\Omp(T_xM)$. Then the metric $g$ is locally isometric to:
\begin{enumerate}
\item \label{list4.2.1}A flat metric.
\item \label{list4.2.2}A Lorentzian, non-Riemannian, left-invariant metric on $\tilde{\PSL}(2,\RR)$.
\item \label{list4.2.3}The Lorentz--Heisenberg metric $g_{LH}$ on the group $\Heis$.
\item \label{list4.2.4}The Lorentz--Sol metric $g_{\sol}$ on the group $\SOL$.
\end{enumerate}
\end{theo}
The Theorem applies in our situation since Proposition~\ref{prop.dimension} ensures that $\mathfrak{Is}(x)$ generates a noncompact subgroup of $\Omp(T_xM)$ for almost every (hence every) point $x \in M$.

Lorentzian, non-Riemannian, left-invariant metrics $g_{AdS},g_u$ and $g_h$ on $\tilde{\PSL}(2,\RR)$ were introduced in Section~\ref{sec.exemplessl2}, while Lorentz--Heisenberg geometry was described in Section~\ref{sec.lorentzheisenberg}. It remains to explain what is the Lorentz--Sol geometry.

The Lie algebra $\mathfrak{sol}$ is the $3$-dimensional Lie algebra with basis $T,X,Y$ and nontrivial bracket relations $[T,Z]=Z$, $[T,X]=-X$. The corresponding connected, simply connected Lie group is denoted $\SOL$. On $\mathfrak{sol}$, we can consider the Lorentz scalar product such that $\langle T,Z\rangle=\langle X,X\rangle=1$, and all other products are $0$. After left-translating this scalar product on $\SOL$, we get a Lorentz metric $g_{\sol}$ on $\SOL$ which is called the Lorentz--Sol metric. The isometry group of $(\Sol,g_{\sol})$ contains $\Sol$ (acting by left-translations), but it is actually $4$-dimensional. The Lie algebra of Killing fields is obtained by adding $Y$ to $T,X,Z$, with bracket relations $[T,Y]=2Y$ and $[X,Y]=Z$ (see~\cite[Section~4.2]{sorin.ghani} for further details).

\subsection{Ruling out Lorentz-\texorpdfstring{$\SOL$}{SOL} geometry}

We first establish:
\begin{prop}\label{prop.compacite.sol}
Let $(M,g)$ be a closed, $3$-dimensional Lorentz manifold locally modelled on $(\SOL,g_{\sol})$. Then $\Iso(M,g)$ is a compact group.
\end{prop}


The key property for proving Proposition~\ref{prop.compacite.sol} is a Bieberbach rigidity theorem for closed manifolds modelled on Lorentz--Sol geometry.

\begin{theo}[{\cite[Theorem~1.2$\MK$(iv), and Proof of Proposition~7.1]{sorin.ghani}}]\label{thm.bieberbach.sol}
Let $(M,g)$ be a closed, $3$-dimensional Lorentz manifold locally modelled on $(\SOL,g_{\sol})$. Then $(M,g)$ is isometric to the quotient of $(\SOL,g_{\sol})$ by a discrete subgroup $\Gamma \subset \Iso(\SOL,g_{\sol})$. Moreover, the intersection of $\Gamma$ with $\Iso^o(\SOL,g_{\sol})$ is a lattice $\Gamma_0 \subset \SOL$ acting by left translations.
\end{theo}



We know that $(M,g)$ is isometric to some quotient $\Gamma \backslash \Sol$, by Theorem~\ref{thm.bieberbach.sol}. Let us denote by $L_{\SOL}$ the subgroup of $\Iso(\Sol,g_{\sol})$ comprising all left-translations by elements of $\Sol$. The group $\Gamma_0=\Gamma \cap \Iso^o(\Sol,g_{\sol})$ is Zariski-dense in $L_{\SOL}$ by Theorem~\ref{thm.bieberbach.sol}. It follows that $\Nor(\Gamma)$, the normalizer of $\Gamma$ in $\Iso(\Sol,g_{\sol})$, must normalize $L_{\SOL}$. But the description of $\mathfrak{Iso}(\Sol,g_{\sol})$ made above shows that $L_{\SOL}$ has finite index in its normalizer. We infer that $\Nor(\Gamma)/\Gamma$ is compact, which proves Proposition~\ref{prop.compacite.sol}.


\subsection{Minkowski and Lorentz--Heisenberg geometries with noncompact isometry group}\label{sec.minkowski}
We now focus on closed Lorentz manifolds locally modelled on Minkowski space, or on Lorentz--Heisenberg geometry. In both cases, one has a Bieberbach type theorem. This is very well known in the flat case, thanks to the works~\cite{fried.goldman} and~\cite{goldman.kami}, and the completeness result of Carrière for closed flat Lorentz manifolds~\cite{carriere}. For manifolds modelled on Lorentz--Heisenberg geometry, this is proved in~\cite[Proposition~8.1]{sorin.ghani}. The precise statement is the following:

\begin{theo}[Bieberbach's theorem for flat and Lorentz--Heisenberg manifolds]\label{thm.bieberbach.minkowski}
Let $(M,g)$ be a closed, $3$-dimensional, Lorentz manifold.
\begin{enumerate}
\item \label{list4.5.1}If $(M,g)$ is flat, there exists a discrete subgroup $\Gamma \subset \Iso(\RR^{1,2})$ such that $(M,g)$ is isometric to the quotient $\Gamma \backslash \RR^{1,2}$. Moreover, there exists a connected $3$-dimensional Lie group $G \subset \Iso(\RR^{1,2})$, which is isomorphic to $\RR^3$, $\Heis$ or $\Sol$, and which acts simply transitively on $\RR^{1,2}$, satisfying that $\Gamma_0=G \cap \Gamma$ has finite index in $\Gamma$ and is a uniform lattice in $G$.
\item \label{list4.5.2}If $(M,g)$ is locally modelled on Lorentz--Heisenberg geometry, then it is isometric to the quotient of $(\Heis,g_{\Heis})$ by a discrete subgroup $\Gamma \subset \Iso(\Heis,g_{\Heis})$. Moreover, there exists a finite index subgroup $\Gamma_0 \subset \Gamma$ which is a lattice $\Gamma_0 \subset \Heis$ acting by left translations.
\end{enumerate}
\end{theo}

Notice that in~\cite{sorin.ghani}, the authors make use of the classification result obtained in~\cite{zeghibflot}, namely Theorem~\ref{thm.ghaniclassification}, to prove this Bieberbach's theorem for Lorentz--Heisenberg manifolds. We will explain at the end of the paper (Appendix~\ref{ann.completeness}) how to adapt the proof of~\cite[Proposition~8.1]{sorin.ghani} in order to avoid the use of~\cite{zeghibflot}. Hence, there is no vicious circle in our arguments, and our main results are genuinely independent of Theorem~\ref{thm.ghaniclassification}.



Theorem~\ref{thm.bieberbach.minkowski} says that \emph{up to finite cover}, a closed Lorentz manifold $(M,g)$ modelled on Minkowski, or Lorentz--Heisenberg geometry, is homeomorphic to $\TT^3$ or to $\TT_A^3$ for $A \in \SL(2,\ZZ)$ hyperbolic or parabolic. Noncompactness of the isometry group allows to be more precise.


\begin{prop}\label{prop.topologie.minkowski}
Let $(M,g)$ be a closed, $3$-dimensional Lorentz manifold, such that $\Iso(M,g)$ is noncompact. We assume that $(M,g)$ is orientable and time-orientable.
\begin{enumerate}
\item \label{list4.6.1}If $(M,g)$ is flat, then $M$ is diffeomorphic either to a torus $\TT^3$, or to a torus bundle $\TT_A^3$ with $A \subset \SL(2,\ZZ)$ hyperbolic, or parabolic.
\item \label{list4.6.2}If $(M,g)$ is modelled on Lorentz--Heisenberg geometry, then $M$ is diffeomorphic to a torus bundle $\TT_A^3$ with $A \in \SL(2,\ZZ)$ parabolic ($A \not =\id$).
\end{enumerate}
\end{prop}

\begin{proof}
The situation provided by Theorem~\ref{thm.bieberbach.minkowski} is the following (both in the flat and Lorentz--Heisenberg case). We have a $3$-dimensional Lie group $G$, which is either $\RR^3$, $\Heis$ or $\Sol$, as well as a left-invariant metric $\mu$ on $G$, and the manifold $(M,g)$ is isometric to a quotient of $(G,\mu)$ by a discrete subgroup $\Gamma \subset \Iso(G,\mu)$. Moreover, if we denote by $L_G$ the group of left-translations by elements of $G$, the intersection $\Gamma_0=\Gamma \cap L_G$ has finite index in $\Gamma$, and is a uniform lattice in $L_G$. Observe that if $\Nor(\Gamma)$ denotes the normalizer of $\Gamma$ in $\Iso(G,\mu)$, then the isometry group $\Iso(M,g)$ coincides with the quotient group $\Gamma \backslash \Nor(\Gamma)$.

An important remark for the following is that the group $\Nor(\Gamma)$ normalizes $L_G$. It is obvious in the case of Lorentz--Heisenberg geometry $(G,\mu)=(\Heis,g_{\Heis})$. In this case, the identity component $\Iso^o(G,\mu)$ is of the form $\RR \ltimes L_G$ (see~\cite[Section~4.1]{sorin.ghani}), and $L_G$ is thus normalized by the full isometry group $\Iso(G,\mu)$.

In the case of Minkowski geometry, one has to remember that the group $G$ (more accurately $L_G$) is the identity component of the \emph{crystallographic hull} of $\Gamma$ (see~\cite[Section~1.4]{fried.goldman}). The last part of~\cite[Theorem~1.4]{fried.goldman} ensures that in the case of Minkowski geometry, the crystallographic hull is unique. It follows that $\Nor(\Gamma)$ must normalize this crystallographic hull, as well as its identity component $L_G$.

As a consequence, elements of $\Nor(\Gamma)$ (in particular elements of $\Gamma$) belong to the group $\Aut(G)_{\mu} \ltimes L_G$, where $\Aut(G)_{\mu}$ denotes the automorphisms of $G$ preserving the metric $\mu$. Let us denote by $\Gamma_l$ the projection of $\Gamma$ on $\Aut(G)_{\mu}$. If this projection is trivial, we get that $\Gamma \subset L_G$. The manifold $M$ is obtained as a quotient of $\RR^3$, $ \Sol$ or $ \Heis$ by a uniform lattice, and we are done.


If this projection is nontrivial, we are going to get a contradiction. Indeed, $\Gamma_l$ must be a finite subgroup of $\Aut(G)_{\mu}$, because $\Gamma_0 \subset L_G$ has finite index in $\Gamma$. The isotropy representation $\rho$ of $\Aut(G)_{\mu}$ at $e$ identifies $\Gamma_l$ with a nontrivial finite subgroup of $\SO(\mu)$, which is actually in $\SO^o(\mu)$ because $(M,g)$ is orientable and time-orientable. Let $\Nor^o(\Gamma)$ be the intersection of $\Nor(\Gamma)$ with $\Aut^o(G)_{\mu} \ltimes L_G$, and let $N_l$ the projection of $\Nor^o(\Gamma)$ on $\Aut^o(G)_{\mu}$. We get that $\rho(N_l)$ normalizes $\rho(\Gamma_l)$. Now, nontrivial finite groups in $\SO^o(\mu)$ have a unique fixed point in ${\mathbb H}^2$, so their normalizer in $\SO^o(\mu)$ are contained in a compact subgroup. It follows that $\Nor^o(\Gamma)$ is contained in a subgroup $K \ltimes L_G$, with $K$ compact in $\Aut^o(G)_{\mu}$. We thus see that $\Nor^o(\Gamma)/\Gamma_0$ is compact, implying the compactness of $\Nor(\Gamma)/ \Gamma$. This in turns implies $\Iso(M,g)$ compact: Contradiction.
\end{proof}

\subsection{Anti-de Sitter structures with noncompact isometry group}

It remains to study closed Lorentz manifolds $(M,g)$ modelled on a Lorentzian, non-Riemannian, left-invariant metric on $\tilde{\PSL}(2,\RR)$. We observe that the identity components $\Iso^o(\tilde{\PSL}(2,\RR),g_u)$ and $\Iso^o(\tilde{\PSL}(2,\RR),g_h)$ are actually contained in $\Iso^o(\tilde{\PSL}(2,\RR),g_{AdS})$. Thus, if $(M,g)$ is a closed, orientable and time-orientable, Lorentz manifold modelled on $(\tilde{\PSL}(2,\RR),g_u)$ or $(\tilde{\PSL}(2,\RR),g_h)$, there exists an anti-de Sitter metric $g'$ on $M$ which is preserved by a finite index subgroup of $\Iso(M,g)$. Hence, it will be enough for us to focus on the topology of closed anti-de Sitter manifolds with noncompact isometry group. In the sequel, we will denote $\tilde{\bf AdS}_3$ the space $(\tilde{\PSL}(2,\RR),g_{AdS})$

\begin{prop}\label{prop.topologieads}
Let $(M,g)$ be a closed, orientable and time-orientable, anti-de Sitter manifold. If $\Iso(M,g)$ is noncompact, then $M$ is homeomorphic to a quotient $\Gamma \backslash \tilde{\PSL}(2,\RR)$, for a uniform lattice $\Gamma \subset \tilde{\PSL}(2,\RR)$.
\end{prop}

It is worth noticing that all closed (orientable and time-orientable) $3$-dimensional anti-de Sitter manifolds are Seifert fiber bundles over hyperbolic orbifolds (a short proof of this fact can be found in~\cite[Corollary~4.3.6]{tholozan}). Conversely, any Seifert fiber bundle over a hyperbolic orbifold, with nonzero Euler number, can be endowed with an anti-de Sitter metric (see~\cite{scott}). The assumption that $\Iso(M,g)$ is noncompact reduces the possibilities for the allowed Seifert bundles. For instance, all nontrivial circle bundles over a closed orientable surface of genus $g \geq 2$ admit anti-de Sitter metrics, but only those for which the Euler number divides $2g-2$ do occur in Proposition~\ref{prop.topologieads}.

It was shown in~\cite{klingler} that closed anti-de Sitter manifolds are complete. It follows that $(M,g)$ as in Proposition~\ref{prop.topologieads} is a quotient of $\tilde{\bf AdS}_3$ by a discrete subgroup $\tilde{\Gamma} \subset \Iso(\tilde{\bf AdS}_3)$. Actually $\tilde{\Gamma} \subset \Iso^o(\tilde{\bf AdS}_3)$ because $(M,g)$ is orientable and time-orientable. The center of $\tilde{\PSL}(2,\RR)$ is infinite cyclic, generated by an element $\xi$. The group $\tilde{\PSL}(2,\RR) \times
\tilde{\PSL}(2,\RR)$ acts on ${\tilde{\bf AdS}_3}$ by left and right translations: $(h_1,h_2).g=h_1gh_2^{-1}$. This yields an epimorphism $\tilde{\PSL}(2,\RR) \times
\tilde{\PSL}(2,\RR) \to \Iso^o({\tilde{\bf AdS}_3})$, with infinite cyclic kernel generated by $(\xi,\xi)$. The group $\Iso^o({\tilde{\bf AdS}_3})$ has a center $Z$ which is generated by the left action of $\xi$. Doing the quotient of $\Iso^o({\tilde{\bf AdS}_3})$ by $Z$ yields an epimorphism $\pi: \Iso^o({\tilde{\bf AdS}_3}) \to \PSL(2,\RR) \times \PSL(2,\RR)$. Notice that $\PSL(2,\RR) \times \PSL(2,\RR)$ coincides with the identity component of the isometries of $\PSL(2,\RR)$ endowed with its anti-de Sitter metric.


An important result, known as \emph{finiteness of level}, says that $\tilde{\Gamma} \cap Z \not = \id$. This was first stated in~\cite{kulkarni.raymond}. A detailed proof can be found in~\cite[Theorem~3.3.2.3]{saleinthese}. Geometrically, this theorem ensures that there exists a finite group of isometries $\Lambda \subset \Iso(M,g)$, which acts freely and centralizes a finite index subgroup of $\Iso(M,g)$, such that the quotient manifold of $(M,g)$ by $\Lambda$ is a quotient of $\PSL(2,\RR)$ by a discrete group $\overline{\Gamma} \subset \PSL(2,\RR) \times \PSL(2,\RR)$. Let us denote by $(\overline{M}^3,\overline{g})$ this new Lorentz manifold, and observe that $\Iso(\overline{M}^3,\overline{g})$ is still noncompact. Observe also that the projection $\pi$ maps $\tilde{\Gamma}$ onto $\overline{\Gamma}$.

The structure of the group $\overline{\Gamma}$ is well understood. Up to conjugacy, there exists $\Gamma_0$ a uniform lattice in $\PSL(2,\RR)$, and a representation $\rho: \Gamma_0 \to \PSL(2,\RR)$ such that
\[
\overline{\Gamma}=\{ (\gamma,\rho(\gamma)) \in \PSL(2,\RR) \times \PSL(2,\RR) \, | \, \gamma \in \Gamma_0 \}.
\]
This was established in~\cite[Theorem~5.2]{kulkarni.raymond} when $\Gamma$ is torsion-free. For a group with torsion, the adapted proof can be found in~\cite[Lemma~4.3.1]{tholozan}.

Because the group $\Iso(\overline{M}^3,\overline{g})$ is noncompact, a result of Zeghib (\cite[Theorem~1.2]{zeghibisom1}) ensures that $(\overline{M}^3,\overline{g})$ must admit a codimension one, transversally Lipschitz lightlike and totally geodesic foliation. Such foliations ${\mathcal F}$ in $\PSL(2,\RR)$ endowed with the anti-de Sitter metric are well understood (see for instance~\cite[Lemme~3.3.2.9]{saleinthese}, or~\cite[Section~15.2]{zeghibisom1}). Let $AG \subset \PSL(2,\RR)$ be the connected $2$-dimensional group corresponding to the upper-triangular matrices (it is isomorphic to the affine group of the line). Then up to conjugacy in $\PSL(2,\RR) \times \PSL(2,\RR)$, the leaves of ${\mathcal F}$ are given by $\{ gAG \ | \ g \in \PSL(2,\RR)\}$ or $\{ AGg \ | \ g \in \PSL(2,\RR)\}$. If $\overline{\Gamma}$ preserves such a foliation ${\mathcal F}$, we infer that the leaves are of the form $\{ gAG \ | \ g \in \PSL(2,\RR)\}$, and $\rho(\Gamma_0)$ normalizes $AG$, namely $\rho(\Gamma_0) \subset AG$.

We now consider the normalizer $H$ of $\overline{\Gamma}$ in $\PSL(2,\RR) \times \PSL(2,\RR)$. The first projection $\pi_1(H)$ must normalize $\Gamma_0$. Since uniform lattices in $\PSL(2,\RR)$ are of finite index in their normalizer, we can replace $H$ by a finite index subgroup and assume $\pi_1(H)=\Gamma_0$. Let us consider $h=(h_1,h_2)$ in $H$, and $\gamma \in \Gamma_0$. Because $H$ normalizes $\overline{\Gamma}$, $(h_1 \gamma h_1^{-1}, h_2 \rho(\gamma)h_2^{-1}) \in \overline{\Gamma}$, which implies $h_2 \rho(\gamma)h_2^{-1}=\rho(h_1)\rho(\gamma) \rho(h_1)^{-1}$. In other words, $h_2^{-1}\rho(h_1)$ centralizes $\rho(\Gamma)$. As a consequence, $\rho(\Gamma)$ can not be Zariski dense in $AG$. Otherwise, $h_2^{-1}\rho(h_1)$ should be trivial implying that $h=(h_1,h_2)$ actually belongs to $\overline{\Gamma}$. We thus would get that $\Iso(\overline{M}^3,\overline{g})$ is finite, a contradiction.

As a result, $\rho(\Gamma)$ is included in a $1$-parameter subgroup of $AG$. This implies that the group $\tilde{\Gamma}$ is contained in a product $\tilde{\PSL}(2,\RR) \times \RR$, where $\tilde{\PSL}(2,\RR)$ acts by left translations, and $\RR \subset \tilde{\PSL}(2,\RR)$ is a $\RR$-split or unipotent $1$-parameter group acting on the right. We consider the projection $\pi_1: \tilde{\Gamma} \to \tilde{\PSL}(2,\RR)$ on the left-factor. The group $\Gamma:=\pi_1(\tilde{\Gamma})$ projects surjectively on $\overline{\Gamma}$ by $\pi$, hence is a uniform lattice in $\tilde{\PSL}(2,\RR)$. Moreover, the kernel of $\pi_1$ must be trivial, otherwise some nontrivial element of $\tilde{\Gamma}$ would belong to $\{\id \} \times \RR$, and $\Nor(\tilde{\Gamma})/\tilde{\Gamma}$ would be compact, contradicting the hypothesis $\Iso(M,g)$ noncompact. It follows that $\tilde{\Gamma}$ is isomorphic to $\Gamma$.

In conclusion, the two manifolds $\Gamma \backslash \tilde{\PSL}(2,\RR)$ and $\tilde{\Gamma} \backslash \tilde{\PSL}(2,\RR)$ are two Seifert bundles over the hyperbolic orbifold $\Gamma_0 \backslash {\mathbb H}^2$. Their Euler number is nonzero, so that they are both large Seifert manifolds (see~\cite[p.~92]{orlik}). Their fundamental groups are isomorphic (to $\Gamma$), hence by~\cite[Theorem~6, p.~97]{orlik} they are homeomorphic.



\subsection{Conclusions}
The previous results show that closed, orientable and time-orientable, Lorentz $3$-dimensional manifolds which are locally homogeneous are homeomorphic to $\TT^3$, a torus bundle $\TT_A^3$ for $A$ hyperbolic or parabolic, or a quotient $\Gamma \backslash \tilde{\PSL}(2,\RR)$. This proves Theorem~\ref{thm.topological} for manifolds such that all components of $\mint$ are locally homogeneous. Our analysis shows moreover that when $M$ is homeomorphic to a hyperbolic torus bundle or to a $3$-torus, the metric must be flat. When $M$ is homeomorphic to a parabolic torus bundle, the metric is either flat, or locally modelled on Lorentz--Heisenberg geometry. When $M$ is homeomorphic to $\Gamma \backslash \tilde{\PSL}(2,\RR)$, the geometry is locally anti-de Sitter or locally modelled on a Lorentzian, non-Riemannian, left-invariant metric on $\tilde{\PSL}(2,\RR)$. This is in accordance to points~\eqref{list3.2}, \eqref{list3.3}, \eqref{list3.4} of Theorem~\ref{thm.geometrical}.



\section{Manifolds admitting a hyperbolic component}\label{sec.hyperbolic}

Our aim in this section is to prove Theorem~\ref{thm.topological} under the assumption that our Lorentz manifold $(M,g)$ admits at least one hyperbolic component ${\mathcal M}$. Actually Theorem~\ref{thm.topological} will be implied by a more precise description provided by Theorem~\ref{theo.hyperbolic} to be stated below.

\subsection{A first reduction}



\begin{prop}\label{prop.hyperbolic}
Assume that the integrability locus $\mint$ contains a hyperbolic component. Then either $(M,g)$ is locally homogeneous, or there exists a hyperbolic component which is not locally homogeneous.
\end{prop}

\begin{proof}
Assume that all hyperbolic components in $\mint$ are locally homogeneous, and consider $\calm \subset \mint$ such a component. It is open by definition, and we are going to show that the boundary $\partial \calm$ is empty, which will yield $\calm=M$ by connectedness of $M$. Local homogeneity of $M$ will follow.

Let us assume that there exists $x \in \partial \calm$. We are going to see that $x \in \mint$, which will yield a contradiction. We pick $x_0 \in \calm$. In the sequel, we use the notation $\dk(z)$ for the $\Omp(1,2)$-orbit of $\dk(\hz)$ in ${\mathcal W}$ (where $\hz$ is any point in the fiber of $z$). By assumption, the rank of $\dk$ is constant equal to $2$ on $\calm$. Moreover, by local homogeneity, $\dk({\calm})$ is a single $\Omp(1,2)$-orbit in ${\mathcal W}$. This orbit is $2$-dimensional because $\calm$ is locally homogeneous, but does not have constant curvature. The stabilizer of points in $\dk(x_0)$ are hyperbolic $1$-parameter subgroups of $\Omp(1,2)$ by assumption that ${\mathcal M}$ is hyperbolic. In every finite-dimensional representation of $\Omp(1,2)$, such hyperbolic orbits are closed. This is a standard fact, the proof of which is, for instance, detailed in~\cite[Appendix~B]{frances.open}. It follows that $\dk(x)=\dk(x_0)$. Hyperbolic $1$-parameter groups are open in the set of $1$-parameter groups of $\Omp(1,2)$. It follows that there is a sufficiently small neighbourhood $U$ of $x$ in $M$ such that the rank of $\dk$ on $U$ is $\geq 2$, and for every $y \in U$, the orbit $\dk(y)$ is $2$-dimensional with hyperbolic $1$-parameter groups as stabilizers of points (notice that $\dk(y)$ can not be $3$-dimensional because of the first point of Proposition~\ref{prop.dimension}). If at some point $y \in U$, the rank of $\dk$ is $3$, then $y$ belongs to a component of $\mint$ which is not locally homogeneous, by Lemma~\ref{lem.rank3}. This component must be hyperbolic because stabilizers in $\dk(y)$ are hyperbolic. Since we assumed that there are no such components, it follows that the rank of $\dk$ is constant equal to $2$ on $U$. But then $x \in \mint$ leading to the desired contradiction.
\end{proof}

The case of locally homogeneous manifolds was already settled in Section~\ref{sec.homogeneous}, so that we will assume \emph{in all the remaining part of this section} that $\mint$ contains a component which is hyperbolic but not locally homogeneous. We will call $\calm$ this component. We are going to show that under these circumstances, $M$ is diffeomorphic to a hyperbolic torus bundle. More precisely, the geometry of $(M,g)$ can be described as follows:

\begin{theo}\label{theo.hyperbolic}
Assume that $(M,g)$ is a closed, orientable and time-orientable~$3$-dimensional Lorentz manifold, such that $\Iso(M,g)$ is noncompact. Assume that $(M,g)$ admits a hyperbolic component which is not locally homogeneous. Then
\begin{enumerate}
\item \label{list5.2.1}The manifold $M$ is diffeomorphic to a $3$ torus ${\mathbb T}^3$, or a torus bundle ${\mathbb T}_A^3$ where $A \in \SL(2,\ZZ)$ is a hyperbolic matrix.
\item \label{list5.2.2}The universal cover $(\tilde{M},\tilde{g})$ is isometric to $\RR^3$ endowed with the metric $\dt^2+2a(t)\du\dv$ for some positive nonvanishing, periodic, smooth function $a: \RR \to (0,+\infty)$.
\item \label{list5.2.3}There is an isometric action of the Lie group $\Sol$ on $(\tilde{M},\tilde{g})$.
\end{enumerate}
\end{theo}

The proof of Theorem~\ref{theo.hyperbolic}, will be the aim of Sections~\ref{sec.anosovtori} to~\ref{sec.conclusionhyperbolic} below.

\subsection{Existence of Anosov tori}\label{sec.anosovtori}
We first prove that under the hypotheses of Theorem~\ref{theo.hyperbolic}, one can find an element $h$ in $\Iso(M,g)$ which acts by an Anosov transformation of a flat Lorentz $2$-torus on $M$ (see Lemma~\ref{lem.anosov.tore}).

\subsubsection{Facts about flat Lorentz surfaces}
We begin by recalling elementary and well known facts about closed flat Lorentz surfaces. Since Lorentz manifolds must have zero Euler characteristic, such surfaces are tori or Klein bottles.
\begin{lemm}\label{lem.lorentzsurface}
A closed, flat Lorentz surface $(\Sigma,g)$ admitting a noncompact isometry group is a torus. Moreover, any group $H \subset \Iso(\Sigma,g)$ which does not have compact closure, contains an element $h$ acting on $\Sigma$ by a hyperbolic linear transformation.
\end{lemm}

\begin{proof}
Closed Lorentz manifolds with constant curvature are geodesically complete (\cite{carriere}, \cite{klingler}). It follows that a closed, flat Lorentz surface $(\Sigma, g)$ is a quotient of the Minkowski plane $\RR^{1,1}$, by a discrete subgroup $\Gamma \subset \Omp(1,1) \ltimes \R^2$, which acts freely and properly on $\RR^{1,1}$. Observe that nontrivial elements of $\Omp(1,1) $ are of two kinds. Either they are hyperbolic (namely have two real eigenvalues of modulus $\not = 1$), or they have order $2$ (an orthogonal symmetry with respect to a spacelike (resp. timelike) line). It is readily checked that $\Gamma$ is either a lattice in $\R^2$, or admits a subgroup of index $2$ which is such a lattice. Assume we are in the first case, and let $H \subset \Iso(\Sigma,g)$ be a subgroup which does not have compact closure. We lift $H$ to $\tilde{H} \subset \Omp(1,1) \ltimes \R^2$. The group $\tilde{H}$ has a nontrivial projection on $\SO(1,1)$, otherwise $H$ would be compact. Thus $\tilde{H}$ contains a conjugate of a hyperbolic element of $\Omp(1,1)$, which acts as an Anosov diffeomorphism on $\Sigma$.

In the second case, where the projection of $\Gamma$ on $\Omp(1,1)$ is an order $2$ subgroup, the normalizer $Nor(\Gamma)$ must have trivial projection on $\SO(1,1)$, which implies that $\Gamma$ is cocompact in $Nor(\Gamma)$. This shows that flat Klein bottles have compact isometry~group.
\end{proof}



\subsubsection{Closed \texorpdfstring{$\kiloc$}{kiloc}-orbits} Our next aim is to exhibit some $\kiloc$-orbits which are closed surfaces.

\begin{lemm}\label{lem.anosov.tore}
In every hyperbolic component ${\mathcal M}$, there exists a $\kiloc$-orbit $\Sigma_0$ which is a flat Lorentz $2$-torus, and such that there exists $h \in \Iso(M,g)$ leaving $\Sigma_0$ invariant, and acting on $\Sigma_0$ as a linear hyperbolic automorphism.
\end{lemm}



Let us prove this lemma. We consider our distinguished component ${\mathcal M}$ that, we recall, is not locally homogeneous, and hyperbolic. We pick $x \in \calm$ a hyperbolic point, and we choose $\hx \in {\hat{\calm}}$ in the fiber of $x$. We already observed in the proof of Lemma~\ref{lem.rank3} that the rank of $\dk$ is at most $3$ on $M$. It is exactly $3$ at $\hx$, still by Lemma~\ref{lem.rank3}, hence remains constant equal to $3$ in a neighbourhood of $\hx$. Hence, if $U \subset \hcalm$ is a small open set around $\hx$, $\dk(U)$ is a $3$-dimensional submanifold of ${\mathcal W}$. If $U$ is chosen small enough, the $\Omp(1,2)$-orbit of every point in $\dk(U)$ will be $2$-dimensional and will have hyperbolic $1$-parameter groups as stabilizers of points. Let us now call $\hat{\Lambda}$ the closed subset of $\hm$ where the rank of $\dk$ is $ \leq2$. By Sard's theorem, the $3$-dimensional Hausdorff measure of $\dk(\hat{\Lambda})$ is zero. We infer the existence of $w \in \dk(U) \setminus \dk(\hat{\Lambda})$. Moving $\hx$ inside $U$, we assume that $w=\dk(\hx)$, and we denote by $\calo(w)$ the $\Omp(1,2)$-orbit of $w$ in $\calw$. By $\Omp(1,2)$-equivariance of $\dk$, the inverse image $\dk^{-1}({\mathcal O}(w))$ avoids $\hat{\Lambda}$, hence the rank of $\dk$ is constant equal to $3$ on $\dk^{-1}({\mathcal O}(w))$. Lemma~\ref{lem.rank3} then leads to the inclusion $\dk^{-1}({\mathcal O}(w)) \subset \hmint$. By the discussion right after Theorem~\ref{thm.integrabilite}, the projection of $\dk^{-1}({\mathcal O}(w))$ on $M$ is a submanifold $N$ of $M$. The stabilizer of $w$ in $\Omp(1,2)$ is hyperbolic, thus as mentioned in the proof of Lemma~\ref{lem.rank3}, the orbit $\calo(w)$ is closed in $\calw$. It follows that $N$ is closed in $M$, hence compact. By (the proof of) Theorem~\ref{thm.integrabilite}, the $\Is^{\loc}$-orbit of $x$ is a union of connected components of $N$, and the connected component of $x$ in $N$, denoted $\Sigma_{0}$, coincides with the $\kiloc$-orbit of $x$. It is a connected compact surface in $M$. Let us show that this surface has Lorentz signature. The Lie algebra $\mathfrak{Is}(x)$ is generated by a local Killing field $X$ around $x$, vanishing at $x$, and such that the flow $\{D_{x}\phi_X^t\}
\subset \Omp(T_{x}M)$ is a hyperbolic $1$-parameter group. Linearizing $X$ around $x$ thanks to the exponential map, we see there are two distinct lightlike directions $u$ and $v$ in ${T}_{x}M$ such that the two geodesics $\gamma_u: s \mapsto \exp(x,su)$ and $\gamma_v: s \mapsto \exp(x,sv)$ are left invariant by $\phi_X^t$. In particular, for $s \not =0$ close to $0$, $\dot{\gamma}_u(s)$ and $\dot{\gamma}_v(s)$ are collinear to $X$, hence tangent to ${\mathcal O}(\gamma_u(s))$ and ${\mathcal O}(\gamma_v(s))$ respectively. By continuity, this property must still hold for $s=0$. We infer that ${T}_{x}({\mathcal O}(x))$ contains the two distinct lightlike directions $u$ and $v$, hence has Lorentz signature. By local homogeneity of the $\kiloc$-orbit $\Sigma_0$, we get that $\Sigma_0$ is Lorentz, and moreover has constant Gauss curvature. The only closed Lorentz surfaces of constant curvature are flat tori or Klein bottles.

Now $\Iso(M,g)$ sends $\Sigma_0$ to components of the $\Is^{\loc}$-orbit of $x$, and there are finitely many such components by compactness of $N$. As a consequence the subgroup $H \subset \Iso(M,g)$ leaving $\Sigma_0$ invariant is noncompact. Observe that if $g_0$ is the metric induced by $g$ on $\Sigma_0$, then the injection $H \to \Iso(\Sigma_0,g_0)$ is proper (see for instance~\cite[Proposition~3.6]{zeghibflot}). It follows that $\Iso(\Sigma_0,g_0)$ is a noncompact group. Lemma~\ref{lem.lorentzsurface} ensures that $(\Sigma_0,g_0)$ is a flat Lorentz torus, and there exists $h \in \Iso(M,g)$ acting on $\Sigma_0$ by a hyperbolic linear automorphism.



\subsection{Pushing Anosov tori along the normal flow}\label{sec.anosov.mapping}
From the $2$-torus $\Sigma_0$ and the diffeomorphism $h \in \Iso(M,g)$ given by Lemma~\ref{lem.anosov.tore}, we are going to recover the topology of the whole manifold $M$, as well as its geometry.



\subsubsection{Preliminary definitions}\label{sec.push}
On the torus $\Sigma_0$, $h$ acts as an Anosov diffeomorphism. It means that there are two $1$-dimensional subbundles ${\mathcal E}^s$ and ${\mathcal E}^u$, inducing a $h$-invariant splitting $T(\TT^2)={\mathcal E}^s \oplus {\mathcal E}^u$, so that vectors in ${\mathcal E}^s$ (resp. in ${\mathcal E}^u$) are exponentially contracted under $Df^n$ as $n \to + \infty$ (resp. $n \to - \infty$). This property and the fact that $h$ acts isometrically for the metric $g$ imply that the bundles ${\mathcal E}^s$ and ${\mathcal E}^u$ are lightlike. We choose a frame field $(E^-,E^+)$ with the property that $E^-$ and $E^+$ are future lightlike, satisfy $g(E^-,E^+)=1$, and generate the ${\mathcal E}^u$ and ${\mathcal E}^s$ respectively. Because $M$ is assumed to be orientable, this defines a smooth normal field $\nu: \Sigma_0 \to T\Sigma_0^{\perp}$ with the property that $(E^-,E^+,\nu)$ is a direct frame of $T_zM$ at each point $z \in \Sigma_0$, and $g(\nu,\nu)=+1$.



In all the rest of the section, we pick once for all $z_0 \in \Sigma_0$ a periodic point of $h$ (recall that the set of periodic points is dense in $\Sigma_0$). This point has period $m_0$, and replacing $h$ by $h^{2m_0}$ if necessary, \emph{we will assume henceforth that $h(z_0)=z_0$ and $h^* \nu =\nu$}.

For every $z \in \Sigma_0,$ we will call $\gamma_z$ the oriented geodesic arc through $z$, with tangent $\nu(z)$ at $z$. Observe that $\gamma_{z_0}$ is a closed spacelike geodesic. Indeed, since $h$ is a Lorentz isometry, the fixed points set $\Fix(h)$ is a closed, totally geodesic submanifold of $M$. The matrix of the differential $D_{z_0}h$, expressed in the basis $(E^-(z_0),E^+(z_0), \nu(z_0))$, is of the form $
\begin{smallpmatrix} \sfrac{1}{\lambda_0}&0&0\\ 0& \lambda_0&0\\ 0&0&1
\end{smallpmatrix}$ with $|\lambda_0|>1$. Linearizing $h$
around $z_0$ thanks to the exponential map, we see that the component of $\Fix(h)$ containing $z_0$, is precisely $\gamma_{z_0}$.



\subsubsection{The normal flow, and an auxiliary pseudo-Riemannian manifold}\label{sec.normalflow}
It will be useful in the sequel to consider the manifold $N=\RR \times
\Sigma_0$. On this manifold, we have the vector field $\frac{\partial}{\partial t}$. Pushing the vector fields $E^-,E^+$ on $\{ 0 \} \times \Sigma_0$ by the flow of $\frac{\partial}{\partial t}$, we get two more vector fields $\tilde{E}^-,\tilde{E}^+$ on $N$. The frame field $(\tilde{E}^-,\tilde{E}^+,\frac{\partial}{\partial t})$ provides $N$ with an orientation.


Let us consider the map $f: (t,z) \mapsto \exp(z,t\nu(z))$. It is well-defined and smooth on some maximal open subset $U_{max} \subset N$. An easy application of the inverse mapping theorem shows that $f: (-\epsilon,\epsilon) \times \Sigma_0 \to M$ is a one-to-one immersion for small $\epsilon >0$.



A key property of the map $f$ is its equivariance with respect to the action of $h$, namely :
\begin{equation}\label{eq.equivariance}
h \circ f(t,z) = f(t,h(z)),
\end{equation}
which is available for $(t,z) \in U_{max}$ (observe that $f(U_{max})$ is left invariant by $h$). Relation~\eqref{eq.equivariance} just follows from the fact that $h$ is an isometry preserving the normal field $\nu$.

In the following, we are going to introduce
\[
\tau_m := \sup\{ s \in (0,\infty) \ | \ f:(0,s) \times \Sigma_0 \subset U_{max} \to M
\text{ is an injective immersion}\}.
\]
It will be sometimes more suggestive to restrict $f$ to $\{0\} \times \Sigma_0$, and consider \emph{the normal flow of $\Sigma_0$}, $\phi^t: \Sigma_0 \to M$ defined by $\phi^t(z):=\exp(z,t \nu(z))$. By what we said before, $\phi^t$ is at least defined on $(0,\tau_m)$, and for all $t \in (0,\tau_m)$, $\phi^t: \Sigma_0 \to M$ is a proper embedding, with image $\Sigma_t \subset M$. Equivariance relation~\eqref{eq.equivariance} shows that $h$ preserves $\Sigma_t$ and acts on it as an Anosov diffeomorphism. In particular, the stable and unstable bundles must be lightlike for the metric $g$, showing that $\Sigma_t$ is a Lorentz torus. We denote by $g_t$ the restriction of $g$ to $\Sigma_t$.



The map $t \mapsto \phi^t(z_0)$ provides a (cyclic) parametrization of the closed geodesic $\gamma_{z_0}$ at speed $+1$. Hence for every $t \in \RR$, the map $z \mapsto \phi^t(z)$ is defined and smooth on some small neighbourhood ${\mathcal U_t} \subset \Sigma_0$ containing $z_0$. We call ${\mathcal E}^{\pm}(t):=D_{z_0}\phi^t(E^{\pm}(z_0))$, and the formula $a(t):=g_{\gamma_{z_0}(t)}({\mathcal E}^{-}(t),{\mathcal E}^{+}(t))$ defines a smooth function $a: \RR \to \RR$. This in turns defines on the manifold $N$ a symmetric $(2,0)$-tensor $\tilde{g}=\dt^2+a(t)g_{0}$ (here $g_0$ is the metric induced by $g$ on $\Sigma_0$). Our main task in the following will be to show that $a(t)$ does not vanish. Doing this, we will prove that $\tilde{g}$ is a genuine Lorentz metric.


\subsubsection{First return time}\label{sec.first}
For any point $z \in \Sigma_0$, the geodesic $\gamma_z$ is defined on some maximal interval $[0,\tau_z^*)$. If $\gamma_z((0,\tau_z^*)) \cap \Sigma_0 \not = \emptyset$, then there exists a smallest $\tau(z) \in (0,\tau_z^*)$ such that $\gamma_z(\tau(z)) \in \Sigma_0$. When $\gamma_z((0,\tau_z^*)) \cap \Sigma_0 = \emptyset$, we just put $\tau(z)=+ \infty$. We introduce \emph{the first return set:}
\[
\Omega_r= \{ z \in \Sigma_0\ | \ \tau(z)<+ \infty \text{ and } \gamma_z'(\tau(z)) \text{ is transverse to }T\Sigma_0 \}
\]
It is clear that $\Omega_r$ is an open set, and it is nonempty because any periodic point of $h$ belongs to $\Omega_r$ (for such points, $\gamma_z'(\tau(z))$ is actually orthogonal to $T\Sigma_0$). Now the map $\varphi: z \mapsto (\tau(z), \gamma_z'(\tau(z)))$ is continuous on $\Omega_r$. When $z$ is periodic for $h$, $\gamma_z'(\tau(z))=\epsilon(z) \nu(\gamma_z(\tau(z)))$, where $\epsilon(z):=\pm 1$. By density of such periodic points, we get that $\varphi$ maps continuously $\Omega_r$ to $\RR_+ \times \{-1,+1\}$. Observe that $\Omega_r$, as well as $\varphi$ are $h$-invariant. Because $h$ is topologically transitive on $\Omega_r$, this implies that $\varphi$ is actually a constant map $z \mapsto (\tau_r,\epsilon)$. We call $\tau_r$ the \emph{first return time} of the normal flow, and $\theta_r=\phi^{\tau_r}: \Omega_r \to \Sigma_0$ \emph{the first return map}.



\subsection{First geometric properties of the normal flow}\label{sec.geometric.properties}
We detail here the main geometric properties of the normal flow $\phi^t$ which, we recall, is defined on $(0,\tau_m)$.


\begin{prop}\label{fact.lorentz}\ \\*[-1.3em]
\begin{enumerate}
\item \label{5.5.1}For each $t \in (0,\tau_m)$, $\phi^t$ is an homothetic transformation from $(\Sigma_0,g_0)$ to $(\Sigma_t,g_t)$. More precisely $a(t)>0$ and $(\phi^t)^*g_t=a(t)g_0$.
\item \label{5.5.2}The tensor $\tilde{g}$ is a Lorentz metric on $(0,\tau_m) \times \Sigma_0$, and $f: ((0,\tau_m) \times \Sigma_0, \tilde{g}) \to (M,g)$ is a one-to-one, orientation preserving, isometric immersion.
\end{enumerate}
\end{prop}


\begin{proof}
Equivariance relation~\eqref{eq.equivariance} implies that $h$ acts as an Anosov diffeomorphism on $\Sigma_t$, and $\mathcal{E}^{-}(t)$ (resp. $\mathcal{E}^{+}(t)$) generates the stable (resp. unstable) bundle of $h$ at $\gamma_{z_0}(t)$. It follows that $\mathcal{E}^{\pm}(t)$ are lightlike, and linearly independent since $D_{z_0}\phi^t$ is one-to-one. This implies $a(t)=g(\mathcal{E}^{-}(t),\mathcal{E}^{+}(t))\not = 0$. Because $a(0)=1$, we get $a(t)>0$ for $t \in (0,\tau_m)$. Relation~\eqref{eq.equivariance} shows that $\varphi^t$ maps the stable (resp. unstable) foliation of $h$ on ${\Sigma}_0$ to the stable (resp. unstable) foliation of $h$ on ${\Sigma}_t$. Hence the differential $D_z\varphi^t$ maps at each point $z$ of ${\Sigma}_0$ the lightcone of ${T}_z{\Sigma}_0$ to the lightcone of ${T}_{\varphi^t(z)}{\Sigma}_t$, which means that $(\varphi^t)^*g_t=\sigma_tg_0$, for some smooth function $\sigma_t : {\Sigma}_0 \to \RR_+^*$. Now, because of~\eqref{eq.equivariance}, the function $\sigma_t$ is $h$-invariant, hence constant since $h$ admits dense orbits on ${\Sigma}_0$. This constant is given by $g(D_{z_0}\varphi^t(E^-(z_0)), D_{z_0}\varphi^t(E^+(z_0)))$, namely ${a}(t)$.


The second point follows easily. Indeed, we already noticed that $a(t)>0$ for $t \in (0,\tau_m)$, which ensures that $\tilde{g}$ is Lorentzian on $(0,\tau_m) \times \Sigma_0$. By the first point, $f$ will be isometric if we prove that $T_{\gamma_z(t)}\Sigma_t$ is orthogonal to $\gamma_z'(t)$ for all $t \in (0,\tau_m)$. Now, observe that for a linear Lorentz transformation $L=\begin{smallpmatrix} \lambda&0&0\\ 0& \sfrac{1}{\lambda}&0\\ 0&0&1\end{smallpmatrix}$ with $|\lambda|>1$, the only Lorentz plane invariant by $L$ is the one generated by the two first basis vectors, namely the orthogonal to the line of fixed point of $L$. This remark shows that if $z \in \Sigma_0$ is a periodic point for $h$, $T_{\gamma_z(t)}\Sigma_t \perp \gamma_z'(t)$ holds. By density of periodic points of $h$ on ${\Sigma}_0$, the property actually holds for all $z \in {\Sigma}_0$. Finally $f$ is orientation preserving because $(E^-,E^+,\nu)$ is positively oriented, and we defined the orientation of $N$ to make $(\tilde{E^-},\tilde{E}^+,\ddt)$ positively oriented.
\end{proof}

Let us also say a few words about the surfaces $\Sigma_t$. We observe that generally, $\Sigma_t$ are not totally geodesic submanifolds of $M$. However, they enjoy the weaker condition:

\begin{fact}\label{fait.lightlike.geodesics}
For every $t \in (0,\tau_m)$, the parametrized lightlike geodesics of $\Sigma_t$ for the metric $g_t$ are parametrized geodesics for the metric $g$.
\end{fact}


This can be checked directly by computation for the metric $\tilde{g}$ on $(0,\tau_m) \times \Sigma_0$. From Fact~\ref{fait.lightlike.geodesics}, we infer the following relation, available for all $z \in \Sigma_0$, $t \in (0,\tau_m)$ and $s \in \RR$~:
\begin{equation}\label{eq.equivariance2}
\phi^t(\exp(z,sE^{\pm}(z)))=\exp(\phi^t(z),sD_z\phi^t(E^{\pm}(z)))
\end{equation}


\subsection{Completeness of the normal flow}
Thanks to the previous section, we understand pretty well the behaviour of the normal flow for $t \in (0,\tau_m)$. Our next step is to show that the flow can be extended for $t \geq \tau_m$.


\subsubsection{Extension of the normal flow at \texorpdfstring{$t=\tau_m$}{t= taum}} An important step toward extending the normal flow for $t=\tau_m$ is to show that the geometry of the Lorentz manifold $((0,\tau_m) \times \Sigma_0, \tilde{g})$ does not degenerate at the boundary. Precisely, we prove :

\begin{lemm}\label{lem.nondegeneracy}
There exists $\epsilon >0$ such that $\tilde{g}$ is a Lorentz metric on the open set $(- \epsilon, \tau_m+\epsilon) \times \Sigma_0 \subset N$.
\end{lemm}

\begin{proof}
We just have to show that $a(\tau_m) \not =0$.

Recall the point $z_0 \in {\Sigma}_0$ we introduced at the beginning of Section~\ref{sec.push}. This point is fixed by $h$, and we already observed that $f(t,z_0)$ exists for all $t \in \RR$. Hence, on ${\mathcal U} \subset {\Sigma}_0$, a small convex neighbourhood of $z_0$ (convex relatively to the metric $g_0$), we have an extended flow $\phi^t: {\mathcal U} \to M$ defined for $t \in (0,\tau_m+\delta)$, $\delta>0$. Saturating ${\mathcal U}$ by the action of $(h^m)_{m \in \ZZ}$, we get a dense open set ${\mathcal V}$ on which $\phi^t$ is defined for $t \in (0,\tau_m+\delta)$. Observe that ${\mathcal V}$ contains the stable and unstable manifolds of $h$ at $z_0$, namely $W^{\pm}:=\{ \exp(z_0,sE^{\pm}) \ | \ s \in \RR \}.$

Along the geodesic $t \mapsto \gamma_{z_0}(t)$, we define two vector fields $(E^-(t),E^+(t))$ by parallel transporting $(E^-(z_0),E^+(z_0))$. For each $t \in (0,\tau_m)$, relation~\eqref{eq.equivariance} yields the existence of two nonzero reals $\lambda_t^{\pm}$ such that ${\mathcal E}^{\pm}(t)=D_{z_0}\phi^t(E^-(z_0))=\lambda_t^{\pm}E^{\pm}(t))$. Observe that ${a}(t)=\lambda_t^+ \lambda_t^-$. Proving $a(\tau_m) \not =0$ amounts to show that $\lambda_t^{\pm}$ are bounded away from $0$ in $(0,\tau_m)$.

Assume for a contradiction that there exists some sequence $(t_k)$ in $(0,\tau_m)$, such that $t_k \to \tau_m$ and $\lambda_{t_k}^- \to 0$. Relation~\eqref{eq.equivariance2} says that for $s \in \RR$, and $k \in \NN$
\[
\phi^{t_k}(\exp(z_0,sE^{-}(z_0)))=\exp(\phi^{t_k}(z_0),s\lambda_{t_k}^-E^-(t_k)).
\]
This implies $\phi^{\tau_m}(\exp(z_0,sE^-(z_0)))=\phi^{\tau_m}(z_0)$ for all $s \in \RR$. In particular, because the unstable manifold $W^-$ is dense in ${\mathcal V}$, we get that $\phi^{\tau_m}(z)=\phi^{\tau_m}(z_0)$ for every $z \in {\mathcal V}$. Let us choose a $h$-periodic point $z_1 \in {\mathcal V}$, $z_1 \not = z_0$, of period $q \in \NN^*$ (such a point exists by density of $h$-periodic points in $\Sigma_0$). By what we just said, $\gamma_{z_0}(\tau_m)=\gamma_{z_1}(\tau_m)$. We observe that $\gamma_{z_0}^{\prime}(\tau_m)$ and $\gamma_{z_1}^{\prime}(\tau_m)$ can not be linearly independent, otherwise $D_{\gamma_{z_0}(\tau_m)}h^q$ would fix pointwise a $2$-dimensional space in $T_{\gamma_{z_0}(\tau_m)}M$, implying that at $\gamma_{z_0}(\tau_m)$, $Dh^q$ is trivial or has order $2$. A Lorentz isometry being completely determined by its first jet at a given point, this situation would lead to $h^{2q}=\id$, a contradiction. We infer that $\gamma_{z_0}'(\tau_m)=-\gamma_{z_1}'(\tau_m)$, and $z_1=\gamma_{z_0}(2 \tau_m)$. Applying the same argument to a periodic point $z_2 \in {\mathcal V}$ different from $z_0$ and $z_1$, we get a contradiction. Exchanging the roles of $W^-$ and $W^+$, the same argument holds if $\lambda_{t_k}^+ \to 0$ for some sequence $t_k \to \tau_m$, and the lemma follows.
\end{proof}

We have shown that the Lorentz metric $\tilde{g}$ on $(0,\tau_m) \times \Sigma_0$ extends to a Lorentz metric on $(-\epsilon, \tau_m+\epsilon) \times \Sigma_0$. Our next goal is to extend our isometric embedding $f$ to a map $\overline{f}: [0,\tau_m] \times \Sigma_0$. This will be done thanks to the following general extension result, which is of independent interest.


\begin{prop}\label{prop.extension}
Let $(L,\tilde{g})$ be a Lorentz manifold, and $\Omega \subset L$ an open subset such that the closure $\overline{\Omega}$ is a manifold with boundary. Assume that the boundary $\partial \Omega$ is a smooth Lorentz hypersurface of $L$. If $(M,g)$ is a closed Lorentz manifold having same dimension as $L$, and if $f: (\Omega,\tilde{g}) \to (M,g)$ is a one-to-one isometric immersion, then $f$ extends to a smooth isometric immersion $\overline{f} : \overline{ \Omega} \to M$.
\end{prop}

In the previous proposition, smooth isometric immersion means that $\overline{f}: \overline{\Omega} \to M$ admits a well defined differential $D_z{\overline{f}}: T_zL \to T_{\overline{f}(z)}M$ for every $z \in \overline{\Omega}$, which is isometric with respect to $\tilde{g}$ and $g$, and varies smoothly with $z$.

\begin{proof}
The main part of the proof is to show the following:

\begin{lemm}\label{lem.extension.locale}
Each point $x \in \partial \Omega$ admits an open neighbourhood $U_x \subset L$ such that
\begin{enumerate}
\item \label{5.9.1}The sets $U_x \cap \Omega$ and ${\mathcal U}_x:=U_x \cap \partial \Omega$ are connected.
\item \label{5.9.2}There exists a smooth injective immersion $\tilde{f}_x: U_x \to M$ such that $\tilde{f}_x$ and $f$ coincide on $U_x \cap \Omega$.
\end{enumerate}
\end{lemm}

\begin{proof}
We consider at $x$, a unit spacelike vector $\nu$ which is normal to $T_x(\partial \Omega)$ and points toward $\Omega$. We consider $\gamma$ a small geodesic segment starting from $x$ and satisfying $\gamma'(0)=\nu$, as well as a sequence $(x_k)$ of points of $\gamma \cap \Omega$ converging to $x$. Since $M$ is compact, we may assume that $f(x_k)$ converges to a point $y \in M$. In small neighbourhoods $U$ and $V$ of $x$ and $y$, we choose two orthonormal frame fields, which yield at each points $z,z'$ of $U$ and $V$ respectively, isometric identifications $i_z: \RR^{1,n-1} \to (T_zL,\tilde{g})$, $i_{z'}:\RR^{1,n-1} \to (T_{z'}M,g)$ (here $\RR^{1,n-1}$ stands for $n$-dimensional Minkowski space). Obviously, one can choose our orthonormal frame fields such that $i_{\gamma(t)}^{-1}(\gamma'(t))$ is a constant vector $\xi \in \RR^{1,n-1}$. Also, there are ${\mathcal U}, {\mathcal V}$ neighbourhoods of the origin in $\RR^{1,n-1}$ (depending only of our initial choice of $U$ and $V$) such that $u \mapsto \exp(z,i_z(u))$, $u \in {\mathcal U}$, and $v \mapsto \exp(z',i_{z'}(v))$, $v \in {\mathcal V}$, make sense and are diffeomorphisms on their images for every $z \in U$ and $z' \in V$. In the trivialization given by the frame fields, the sequence of differentials $(D_{x_k}f)$ becomes a sequence of matrices $(A_k)$ in $\Omp(1,n-1)$. Since $f$ is an isometry, we have the relation
\begin{equation}\label{eq.exponentielle}
f(\exp(x_k,u))=\exp(f(x_k),A_k(u))
\end{equation}
for every $u \in {\mathcal U}$.

We can prove the Lemma~\ref{lem.extension.locale} if we show that the sequence $(A_k)$ is contained in a compact set of $\Omp(1,n-1)$. For if it is the case, we may assume $A_k \to A_{\infty}$, and shrinking maybe $\mathcal{U}$, we will have $A_k(\mathcal{U}) \subset \mathcal{V}$ for all $k \in \NN$. Then, we choose ${\mathcal C} \subset \RR^{1,n-1}$ an open cone with vertex $0$, containing $-a\xi$ (for some small $a>0$) and contained in ${\mathcal U}$. For $k_0$ large enough, $U_x=\exp(x_{k_0}, i_{x_{k_0}}({\mathcal C}))$ contains $x$, and if ${\mathcal C}$ is chosen connected and narrow enough around $-a\xi$, $U_x \cap \Omega$ and $U_x \cap \partial \Omega$ are connected. The map $f_x : U_x \to M$ given by $f_x(\exp(x_{k_0},i_{x_{k_0}}(u)))=\exp(y_{k_0},i_{y_{k_0}}(u))$, $u \in {\mathcal C}$ is a one-to-one immersion which coincides with $f$ on $U_x \cap \Omega$.

It remains to explain why the sequence $(A_k)$ must be bounded. If not, we apply the $KAK$ decomposition of $\Omp(1,n-1)$ to the sequence $(A_k)$, and after considering a subsequence we can write $A_k$ as a product $M_kD_kN_k$ with $M_k \to M_{\infty}$ (resp. $N_k \to N_{\infty}$) in $\Omp(1,n-1)$, and $D_k$ is diagonalisable in a fixed basis $(e_1',\ldots,e_n')$ where it takes the form $
\begin{smallpmatrix} \lambda_k & & \\
& \scalebox{0.4}[0.6]{$\ddots$} & \\
& & \lambda_k^{-1}\\
\end{smallpmatrix} $, $|\lambda_k| \to \infty$. We see that there exists a lightlike hyperplane ${\mathcal H} \subset \RR^{1,n-1}$ (namely the image by $N_{\infty}^{-1}$ of $\Span(e_2',\ldots,e_n')$) with the following dynamical property: For every $u \in {\mathcal H}$, there exists $u_k \to u$ such that after extracting a subsequence, $A_k(u_k) \to u_{\infty}$. Moreover, we see that if $v \not \in {\mathcal H}$, one can find a sequence of real numbers $s_k \to 0$ such that $A_k(s_kv) \to v_{\infty} \not = 0$.

Because ${\mathcal H}$ is lightlike while $T_x(\partial \Omega)$ has Lorentz signature, one can find a nonzero $u \in {\mathcal H} \cap {\mathcal U}$ such that $i_{x}(u) \not \in T_x(\partial \Omega)$ and $i_{x}(u)$ points toward $\Omega$. We choose a sequence $(u_k)$ in $\mathcal{U}$ converging to $u$ such that $A_k(u_k)$ tends to $u_{\infty}$ (after extraction). We can also pick some $v \in {\mathcal{U}}\setminus { \mathcal{H}}$ such that $i_x(v)$ points toward $\Omega$. Then we can find $(s_k)$ a sequence of real numbers tending to $0$ such that $A_k(s_k v)$ converges to $v_{\infty} \not =0$, and $(u_k)$. Observe that $\exp(x_k,u_k)$ and $\exp(x_k, u_k + s_kv)$ belong to $\Omega$ for $k$ large. Now, $f(\exp(x_k,u_k))$ tends to $f(\exp(x,u))$, and relation~\eqref{eq.exponentielle} shows that $f(\exp(x,u))=\exp(y,u_{\infty})$. On the other hand, $f(\exp(x_k, u_k+s_k v))$ should also converge to $f(\exp(x,u))$, because $s_k \to 0$. But relation~\eqref{eq.exponentielle} says that this sequence actually converges to $\exp(y,u_{\infty}+v_{\infty})$. Since $v_{\infty} \not =0$, and because we can rescale $u$ and $v$ so that $u_{\infty}$ and $u_{\infty}+v_{\infty}$ belong to ${\mathcal V}$, we have $\exp(y,u_{\infty}+v_{\infty}) \not = \exp(y,u_{\infty})$, and we get a contradiction.
\end{proof}



Lemma~\ref{lem.extension.locale} easily provides a smooth extension of $f$, $\overline{f}: \overline{\Omega} \to M $, putting $\overline{f}(x)=f(x)$ if $x \in \Omega$, and $\overline{f}(x):=\tilde{f}_x(x)$ for every $x \in \partial \Omega$. This extension map $\overline{f}$ is well defined because if $U_x \cap U_{x'} \not = \emptyset$, then $\tilde{f}_x$ and $\tilde{f}_{x'}$ are equal to $f$ on $U_x \cap U_{x'}\cap \Omega$, hence on $U_x \cap U_{x'}\cap \partial \Omega$. Observe that the relation $f^*g=\tilde{g}$ which is available on $U_x \cap \Omega$ must still hold on $U_x \cap \overline{\Omega}$. This proves that $\overline{f}$ is an isometric immersion.
\end{proof}


\subsubsection{The normal flow at \texorpdfstring{$\tau_m$}{taum} realizes the first return map} We apply Proposition~\ref{prop.extension} choosing for $L$ the Lorentz manifold $((- \epsilon, \tau_m + \epsilon) \times \Sigma_0, \tilde{g})$ and for ${\Omega}$ the product $(0,\tau_m) \times \Sigma_0$. We get a smooth, orientation preserving, extension $\overline{f} : ([0,\tau_m] \times \Sigma_0, \tilde{g}) \to (M,g)$ which is an isometric immersion coinciding with $f$ on $[0,\tau_m) \times \Sigma_0$.

We recall the first return set $\Omega_r \subset \Sigma_0$ and the first return map $\theta_r: \Omega_r \to \Sigma_0$ introduced at the end of Section~\ref{sec.first}.

\begin{prop}
The extension ${\overline{f}}$ maps $\{ \tau_m\} \times {\Sigma}_{0}$ diffeomorphically and isometrically onto $\Sigma_0$. In other words, the first return time $\tau_r$ coincides with $\tau_m$, the first return set $\Omega_r$ coincides with $\Sigma_0$, and $\phi^{\tau_m}: \Sigma_0 \to \Sigma_0$ realizes the first return map.
\end{prop}

\begin{proof}
In the proof, we are going to write $\tilde{\Sigma}_{\tau_m}$ (resp. $\tilde{\Sigma}_0$) instead of
$\{\tau_m \} \times \Sigma_0$ (resp. $\{0\} \times \Sigma_0$). We recall the notations $\frac{\partial}{\partial t}, \tilde{E}^-,\tilde{E}^+$ from Section~\ref{sec.normalflow}. We first show that the restriction of ${\overline{f}}$ to $\tilde{\Sigma}_{\tau_m}$ is one-to-one. Assume for a contradiction that it is not the case. We get two points $\tilde{z}_1=(\tau_m,z_1)$ and $\tilde{z}_2=(\tau_m,z_2)$ on $\tilde{\Sigma}_{\tau_m}$ such that ${\overline{f}}(\tilde{z}_1)={\overline{f}}(\tilde{z}_2)$. Observe that $D_{\tilde{z}_1}{\overline{f}}(T_{\tilde{z}_1}\tilde{\Sigma}_{\tau_m})=D_{\tilde{z}_2}{\overline{f}}(T_{\tilde{z}_2}\tilde{\Sigma}_{\tau_m})$ because a transverse intersection of those two subspaces would not be compatible with injectivity of ${\overline{f}}$ on $(0,\tau_m) \times \Sigma_0$. Looking at orthogonal subspaces, and because ${\overline{f}}$ is an isometric immersion, we get $D_{\tilde{z}_1}{\overline{f}}(\frac{\partial}{\partial t})=\pm D_{\tilde{z}_2}{\overline{f}}(\frac{\partial}{\partial t})$. Again, $D_{\tilde{z}_1}{\overline{f}}(\frac{\partial}{\partial t})= D_{\tilde{z}_2}{\overline{f}}(\frac{\partial}{\partial t})$ would violate the injectivity of ${\overline{f}}$ on $(0,\tau_m) \times \Sigma_0$. We can reformulate equality $D_{\tilde{z}_1}{\overline{f}}(\frac{\partial}{\partial t})=- D_{\tilde{z}_2}{\overline{f}}(\frac{\partial}{\partial t})$, saying that $\gamma_{z_1}(\tau_m)=\gamma_{z_2}(\tau_m)$, and $\gamma_{z_1}'(\tau_m)=-\gamma_{z_2}'(\tau_m)$. It follows that $z_1$ and $z_2$ are in the first return set $\Omega_r$, $z_2=\theta_r(z_1)$, and the first return time $\tau_r$ equals $2 \tau_m$. As a consequence, we get for every $z \in \Omega_r$, the identity:
\[
{\overline{f}}(\tau_m,z)={\overline{f}}(\tau_m,\theta_r(z)).
\]
Defining $\tilde{\theta}_r(t,z):=(t,\theta_r(z))$, this identity yields:
\[
D_{(\tau_m,z)}{\overline{f}}(\tilde{E}^{-})=D_{\tilde{\theta}_r(\tau_m,z))}{\overline{f}}(D_{(\tau_m,z)}\tilde{\theta}_r(\tilde{E}^-)).
\]
Because $\theta_r$ commutes with $h$ by equation~\eqref{eq.equivariance} in Section~\ref{sec.normalflow}, there exists $\alpha(z) \not = 0$ such that $D_{(\tau_m,z)}\tilde{\theta}_r(\tilde{E}^-)=\alpha(z)\tilde{E}^-$. We thus obtain:
\[
D_{(\tau_m,z)}{\overline{f}}(\tilde{E}^{-})= \alpha(z) D_{\tilde{\theta}_r(\tau_m,z))}{\overline{f}}(\tilde{E}^-).
\]
For the same reasons, there exists $\beta(z) \not =0$ such that
\[
D_{(\tau_m,z)}{\overline{f}}(\tilde{E}^{+})= \beta(z) D_{\tilde{\theta}_r(\tau_m,z))}{\overline{f}}(\tilde{E}^+).
\]

Going back to $z=z_0$, $\theta_r(z)=z_1$, we see that $(D_{(\tau_m,z_1)}{\overline{f}})^{-1} \circ D_{(\tau_m,z_0)}{\overline{f}}$ is a linear isometry, preserving the orientation, and sending the direct frame $(\tilde{E}^-,\tilde{E}^+,\ddt)$ at $(\tau_m,z_0)$ to the frame $(\alpha(z_0) \tilde{E}^-, \beta(z_0) \tilde{E}^+,-\ddt)$ at $(\tau_m,z_1)$. The isometric condition yields $\alpha(z_0) \beta(z_0)=1$ and the orientation-preserving condition yields $ \alpha(z_0) \beta(z_0)=-1$. This provides the desired contradiction.


Once we know that ${\overline{f}}$ is one-to-one in restriction to $\tilde{\Sigma}_{\tau_m}$, we get that ${\overline{f}}(\tilde{\Sigma}_{\tau_m})$ is a Lorentz surface of $M$, to which we can again apply the normal flow. This results into an extension of ${\overline{f}}$ to a smooth immersion defined on a domain $(0,\tau_m+\epsilon) \times \Sigma_0$. If ${\overline{f}}(\tilde{\Sigma}_{\tau_m})$ does not meet $\Sigma_0$, it is easily checked that for $\epsilon>0$ small enough, ${\overline{f}}$ is one-to-one on $(0,\tau_m+\epsilon) \times \Sigma_0$, contradicting the definition of $\tau_m$.

We infer that there exist $z_1$ and $z_2$ in $\Sigma_0$ such that ${\overline{f}}(\tau_m,z_1)={\overline{f}}(0,z_2)$. Observe that $D_{(\tau_m,z_1)}{\overline{f}}(T\tilde{\Sigma}_{\tau_m})$ and $D_{(0,z_2)}{\overline{f}}(T \tilde{\Sigma}_0)$ can not intersect transversely, otherwise ${\overline{f}}$ would not be one-to-one on $(0,\tau_m) \times \Sigma_0$. We can recast this property saying that $\gamma_{z_1}(\tau_m)=z_2$, and $\gamma_{z_1}'(\tau_m)$ is orthogonal to $T_{z_2} \Sigma_0$. In other words, $z_1$ belongs to the first return set $\Omega_r$, and the return time is $\tau_r=\tau_m$. It means in particular that for all $z \in \Omega_r$, ${\overline{f}}(\tau_m,z) \in \Sigma_0$. By density of $\Omega_r$ in $\Sigma_0$, we finally get that ${\overline{f}}$ maps $\tilde{\Sigma}_{\tau_m}$ isometrically and diffeomorphically onto $\Sigma_0$.
\end{proof}



\subsection{End of proof of Theorem~\ref{theo.hyperbolic}}\label{sec.conclusionhyperbolic}
Let us just recollect what we did so far. First, showing that the normal flow $\phi^t$ is defined on $(-\epsilon,\tau_m+\epsilon)$ with $\phi^{\tau_m}(\Sigma_0)=\Sigma_0$ immediately implies that $\phi^t$ is defined for every $t \in \RR$. Equivalently, the map $f$ is defined on all of $N=\RR \times \Sigma_0$.

Next, Proposition~\ref{fact.lorentz} implies that $(\phi^{\tau_m})^*g_0=a(\tau_m)g_0$, with $a(\tau_m)>0$. Because the global Lorentz volume of $\Sigma_0$ must be preserved, we get $a(\tau_m)=1$. The transformation $\phi^{\tau_m}$ is a Lorentz isometry of $(\Sigma_0,g_0)$ commuting with $h$: It must be either $\pm \id$ or a linear hyperbolic transformation. The possibility $\phi^{\tau_m}=- \id$ is ruled out by the assumption that $(M,g)$ is time-orientable.

In the following, we denote by $A$ the transformation $\phi^{\tau_m}$. We just showed that $t \mapsto a(t)$ is $\tau_m$-periodic, and thanks to Propositions~\ref{fact.lorentz} and~\ref{prop.extension}, we get that\linebreak $f: (N,\tilde{g}) \to (M,g)$ is an isometric immersion. Let us call $\varphi : N \to N$ the transformation $\varphi(t,x)=(t+1,A^{-1}x)$. Then $\varphi$ acts isometrically for $\tilde{g}$, and $f \circ \varphi = f$. Calling $\Gamma$ the cyclic group generated by $\varphi$, we finally see that $f$ induces an isometry between $\Gamma \backslash N$ (endowed with the metric induced by ${\tilde g}$) and $(M,g)$. This shows the topological part of Theorem~\ref{theo.hyperbolic}.

Since $\Sigma_0$ is a flat torus, the universal cover $(\tilde{M},\tilde{g})$ is isometric to $\RR^3$ endowed with the metric $\dt^2+2a(t)\du\dv$. Affine transformations preserving the planes $t=t_0$ and acting by Lorentz isometries on the Minkowski $(u,v)$-plane, provide an isometric action of $\Sol$ on $(\tilde{M},\tilde{g})$. This shows points~\eqref{list5.2.2} and~\eqref{list5.2.3} of Theorem~\ref{theo.hyperbolic}.



\section{The local geometry of manifolds with no hyperbolic component}\label{sec.local.geometry}

We keep going in our study of closed $3$-dimensional Lorentz manifolds $(M,g)$, such that $\Iso(M,g)$ is not compact. Thanks to Sections~\ref{sec.homogeneous} and~\ref{sec.hyperbolic}, we can prove Theorem~\ref{thm.topological} when all the components of the integrability locus $\mint$ are locally homogeneous, or when there exists at least one hyperbolic component. Looking at the possibilities for the different components listed in Section~\ref{sec.components}, it only remains to investigate the case where all the components are either parabolic or of constant curvature, and there is at least one non locally homogeneous component. This section is devoted to a careful geometric study of such manifolds, and our aim is to prove the


\begin{theo}\label{thm.conf.flat}
Let $(M,g)$ be a $3$-dimensional Lorentz manifold. If all the components of $\mint$ are of constant curvature or parabolic, and if $(M,g)$ is not locally homogeneous, then $(M,g)$ is conformally flat.
\end{theo}


Recall that $(M,g)$ is said to be conformally flat if each sufficiently small open neighbourhood of $M$ is conformally diffeomorphic to an open subset of Minkowski space. Hence, Theorem~\ref{thm.conf.flat} tells us that \emph{at the conformal level}, our structure $(M,[g])$ is locally homogeneous. This local information will be decisive to recover the global properties of $(M,g)$, both topologically and geometrically, a task that will be carried over in Section~\ref{sec.global}.



\subsection{More on the geometry of parabolic components}
Parabolic components split into two categories, the locally homogeneous ones for which the Killing algebra is $4$-dimensional, and the others for which it is $3$-dimensional.

\subsubsection{Locally homogeneous parabolic components}\label{sec.lochomoparabolic}
The study of locally homogeneous parabolic components was made in~\cite{frances.open}, and can be summarized as follows:

\begin{prop}[{\cite[Proposition~4.3]{frances.open}}]\label{prop.lochomoparabolique}
Let $\calm$ be a component of the integrability locus $\mint$ of a $3$-dimensional Lorentz manifold $(M,g)$, which is locally homogeneous and parabolic. Then:
\begin{enumerate}
\item \label{6.2.1} If the scalar curvature of $\calm$ is $0$, the Lie algebra $\kiloc(\calm)$ is isomorphic to a semi-direct product $\RR \ltimes \heis,$ where $\heis$ stands for the $3$-dimensional Heisenberg Lie algebra.
\item \label{6.2.2} If the scalar curvature is nonzero on $\calm$, then $\kiloc(\calm)$ is isomorphic to $\sld \oplus \RR$.
\end{enumerate}
\end{prop}


Actually, the statement of~\cite{frances.open} is slightly more precise since it describes which semi-direct products $\RR \ltimes \heis$ can occur. However, we will not need this extra information here. For the sequel, it will be important to notice that in the first case of Proposition~\ref{prop.lochomoparabolique}, the Lie subalgebra $\heis$ contains the isotropy algebra at each points, hence acts with $2$-dimensional pseudo-orbits (this follows from the computations done in~\cite[Section~4.3.3]{frances.open}).

\subsubsection{Parabolic components which are not locally homogeneous}
We investigate now the geometry of parabolic components which are not locally homogeneous.

\begin{prop}\label{prop.heis}
Let $\calm \subset \mint$ be a parabolic component which is not locally homogeneous.
\begin{enumerate}
\item \label{6.3.1}The Lie algebra $\kiloc(\calm)$ is isomorphic to the $3$-dimensional Heisenberg algebra $\heis$. 
\item \label{6.3.2}The $\kiloc$-orbits on $\calm$ are totally geodesic, lightlike surfaces.
\item \label{6.3.3}The scalar curvature $\sigma$ vanishes on $\calm$.
\end{enumerate}
\end{prop}


The proof of Proposition~\ref{prop.heis} involves quite a bit of computations, that we defer to Appendix~\ref{AnexA}, at the end of the text. Let us mention an important corollary which will be crucial later on.

\begin{coro}\label{coro.islocorbites}
Let $\calm$ be a parabolic component which is not locally homogeneous. Let $x \in \calm$. Then the $\Is^{\loc}$-orbit of $x$ is a submanifold of $\mint$ which is closed in $\mint$.
\end{coro}

\begin{proof}
We already know (Theorem~\ref{thm.integrabilite}) that the $\Is^{\loc}$-orbit of $x$ is a submanifold $\Sigma$ of $\mint$. We have to show that $\Sigma$ is closed in $\mint$. We thus consider a sequence $(x_k)$ of $\Sigma$ converging to a point $x_{\infty} \in \mint$. Let $\hx$ be a lift of $x$ in $\hm$. We recall the generalized curvature map $\dk: \hm \to \mathcal{W}$ (see Section~\ref{sec.curvaturemap}). Let us call $w=\dk(\hx)$, and $\mathcal{O}.w$ the orbit of $w$ under the action of $\Omp(1,2)$ on $\mathcal{W}$. Since $\calm$ is a parabolic component, $\mathcal{O}.w$ is $2$-dimensional and the isotropy at $w$ is a $1$-parameter unipotent subgroup of $\Omp(1,2)$. Let $(\hx_k)$ be a sequence of $\hm$ lifting $(x_k)$, such that $\hx_k \to \hx_{\infty}$. Since $x_k \in \Sigma$ for all $k$, we have $\dk(\hx_k) \in \mathcal{O}.w$ for all $k$, and in particular $\dk(\hx_{\infty})$ belongs to the closure $\oldoverline{\mathcal{O}.w}$. The action of $\Omp(1,2)$ on $\mathcal{W}$ is algebraic. For algebraic actions, the closure of an orbit $\oldoverline{\mathcal{O}.w}$ is made of $\calo.w$ and (maybe) other orbits of dimension smaller than that of $\calo.w$. In particular orbits in $\oldoverline{\mathcal{O}.w} \setminus \calo.w$ have dimension $<2$. Since there are no $1$-dimensional orbits in finite dimensional representations of $\Omp(1,2)$, we conclude that if $\dk(\hx_{\infty})$ belongs to $\oldoverline{\mathcal{O}.w} \setminus \calo.w$, then the stabilizer of $\dk(\hx_{\infty})$ is $3$-dimensional. Since $\hx_{\infty} \in \hmint$, this means that the isotropy algebra $\mathfrak{Is}(x_{\infty})$ is $3$-dimensional, isomorphic to $\oo(1,2)$ (Fact~\ref{rk.isotropie}). Let $\calm'$ the component containing $x_{\infty}$. The points $x_k$ belong to $\calm'$ for $k$ large enough, what shows $\kiloc(\calm') \simeq \kiloc(\calm)\simeq \heis(3)$. This contradicts $\mathfrak{Is}(x_{\infty}) \simeq \oo(1,2)$.

We infer that $\dk(\hx_{\infty}) \in \calo.w$. Hence, replacing the sequence $\hx_k$ by $\hx_k.p_k$ for a bounded sequence $(p_k)$ of $P$, we may assume that $\dk(\hx_k)=w$ for all $k$. By the discussion following Theorem~\ref{thm.integrabilite}, $\dk^{-1}(w) \cap \hmint$ is a submanifold, the connected component of which are $\kiloc$-orbits. We conclude that for $k$ large enough, $\hx_{\infty}$ and $\hx_k$ are in the same $\kiloc$-orbit. The same is thus true for $x_{\infty}$ and $x_k$, and the Corollary~\ref{coro.islocorbites} is proved.
\end{proof}


\subsection{Conformal flatness}\label{sec.conformal.flatness}
Under the standing assumptions stated at the beginning of Section~\ref{sec.local.geometry}, the only non locally homogeneous components in $M$ are parabolic. It follows from Proposition~\ref{prop.heis} that the scalar curvature of $g$ is constant on each component, and equal to zero on the non locally homogeneous ones. As a consequence, the scalar curvature vanishes identically on $M$, which implies that components of constant sectional curvature are actually flat, hence conformally flat. It thus remains to show that all parabolic components (locally homogeneous or not) are conformally flat. Observe that conformal flatness is given by a tensorial condition, namely the vanishing of the Cotton--York tensor in dimension~$3$, so that $(M,g)$ will be conformally flat as soon as a dense open subset of $M$ is. Observe also that the vanishing of the scalar curvature says that locally homogeneous parabolic components are exactly those described by the first point of Proposition~\ref{prop.lochomoparabolique}. This fact, together with Proposition~\ref{prop.heis} and the remark after Proposition~\ref{prop.lochomoparabolique} reduces the proof of Theorem~\ref{thm.conf.flat} to the following general observation:

\begin{prop}\label{prop.heisenberg.flat}
Let $(N,h)$ be a $3$-dimensional Lorentz manifold. Assume that there exists on $N$ a Lie algebra $\lien$ of Killing fields which is isomorphic to $\heis(3)$, and whose pseudo-orbits have dimension $\leq 2$. Then all pseudo-orbits are $2$-dimensional and lightlike, and $(N,h)$ is conformally flat.
\end{prop}

\begin{proof}
Our hypothesis that all pseudo-orbits have dimension $\leq 2$ implies that the isotropy algebra at each point $x \in N$ is a nontrivial subalgebra of $\lien$. This isotropy algebra is thus isomorphic to $\RR$, $\RR^2$ or $\heis(3)$. Because there is no subalgebra of ${\oo}(1,2)$ isomorphic to $\RR^2$ or ${\heis}(3)$, the isotropy algebra of $\lien$ is $1$-dimensional at each point, and all pseudo-orbits of $\lien$ have dimension~$2$. Let us consider $X,Y,Z$ three Killing fields generating $\lien$, and satisfying the relations $[X,Y]=-Z$ and $[X,Z]=[Y,Z]=0$. We are going to look at the subalgebra ${\mathfrak a}$ spanned by $Y$ and $Z$. Because no subalgebra of ${\oo}(1,2)$ is isomorphic to $\RR^2$, pseudo-orbits of ${\mathfrak a}$ have dimension~$1$ or $2$. We claim that the open subset $\Omega$ where the pseudo-orbits of $\liea$ are $2$-dimensional is dense in $N$. To see this, let us consider $\Delta$ a $1$-dimensional pseudo-orbit of $\liea$, and let $x \in \Delta$. The isotropy, in ${\mathfrak a}$, of the point $x$ is spanned by an element $U=aY+bZ$. There is another vector field $V=cY+dZ$ such that $v:=V(x) \not =0$. Since $U$ and $V$ commute, $U$ actually vanishes at each point of $\Delta$. Let $t \mapsto \gamma(t)$ be a geodesic for the metric $h$, satisfying $\gamma(0)=x$, $h(\gamma'(0),v) \not =0$, and $\gamma'(0) \not \in \RR.v$. Clairault's equation ensures that for $t>0$ small enough, $h(\gamma'(t),U(\gamma(t)))=0$ and $h(\gamma'(t),V(\gamma(t))) \not =0$. Observe that $U(\gamma(t)) \not = 0$, because locally, the zero set of a nontrivial Killing field on a $3$-dimensional manifold is a submanifold of dimension $\leq 1$. We thus get that $\gamma(t) \in \Omega$ for $t>0$ small, ensuring the density of $\Omega$.

This density property shows that we will be done if we show that $\Omega$ is conformally flat. To this aim, we consider a point $x_0 \in \Omega$. Since, $[Y,Z]=0$ and $Y,Z$ span a $2$-dimensional space at each point of $\Omega$, there exist local coordinates $(x_1,x_2,x_3)$ around $(0,0,0)$ such that $Z=\delx$ and $Y=\dely$. Because the orbits of $\lien$ are $2$-dimensional, $X$ is of the form $\lambda \delx + \mu \dely$ for some functions $\lambda$ and $\mu$. The bracket relations $[X,Z]=0$ and $[X,Y]=-Z$ lead to $0=\frac{\partial \lambda}{\partial x_1}=\frac{\partial \mu}{\partial x_1}=\frac{\partial \mu}{\partial x_2}$ and $\frac{\partial \lambda}{\partial x_2}=1$. Hence we can~write
\[
X=(x_2+a(x_3))\delx + b(x_3)\dely.
\]
Observe that replacing $X$ by $X-a(0)Z-b(0)Y$ will not affect the bracket relations between $X$, $Y$ and $Z$, so that we will assume in the following that $a(0)=b(0)=0$.

Let us consider a point $p=(p_1,p_2,p_3)$. The vector field $U=X-(p_2+a(p_3))Z-b(p_3)Y$ is nonzero and vanishes at $p$. We compute that at $p$:
\[
\left[U,\delx\right]=0, \quad \left[U,\dely\right]=- \delx,\quad \left[U,\delz\right]=-a'(p_3)\delx-b'(p_3)\dely.
\]
Since $U$ belongs to $\lien$, hence is Killing for the metric $h$, we infer that the matrix $A=\begin{smallpmatrix}
0&-1&a'(p_3)\\
0&0&b'(p_3)\\
0&0&0\\
\end{smallpmatrix}$, which is the matrix of $\nabla U(p)$, must be antisymmetric for the Lorentz scalar product $h_p$. It is readily checked that a rank $1$ nilpotent matrix never has this property (basically because $\exp(tA)$ would be a nontrivial $1$-parameter group in (a conjugate of) $\Omp(1,2)$ fixing pointwise a $2$-plane, which is impossible). We thus infer that the derivative $b'$ is nowhere $0$. In particular, there exist a smooth map $\psi $ defined around $0$, such that $\psi(0)=0$ and $b(\psi(x_3))=x_3$. The transformation
\[
\varphi: (x_1,x_2,x_3) \mapsto (x_1,x_2 -a(x_3),\psi(x_3)x_3).
\]
then yields a local diffeomorphism fixing the origin. Applying $\varphi^*$ to $X,Y,Z$, we get three vector fields
\[
X'=x_2 \delx+x_3 \dely, \quad Y'=\dely,\ \text{ and }\ Z'=\delx
\]
which are Killing for the metric $h'=\varphi^*h$.

Let again $p=(p_1,p_2,p_3)$ be a point in our coordinate chart. The vector field $U'=X'-p_2 Z'-p_3Y'$ vanishes at $p$, and is a Killing field for $h'$. A straightforward computation yields
\[
[U',\delx]=0, \quad [U',\dely]=-\delx,\ \text{ and }\ [U',\delz]=-\dely
\]
everywhere. It follows that the matrix $
\begin{smallpmatrix}
0&1&0\\
0&0&1\\
0&0&0\\
\end{smallpmatrix} $ must be antisymmetric with respect to $h_p'$. This allows us to see that the matrix of $h_p'$ in the frame $(\delx,\dely,\delz)$ is of the form $
\begin{smallpmatrix}
0&0&-\beta(p)\\
0&\beta(p)&0\\
-\beta(p)&0&\gamma(p)\\
\end{smallpmatrix}$, with $\beta(p) >0$. Now, $Z'$ and $Y'$ being Killing fields for $h'$, we see that $\beta$ and $\gamma$ only depend on the variable $x_3$, and we conclude that the metric $h'$ writes as :
\[
-2 \beta(x_3)\dx_1\dx_3+\beta(x_3)\dx_2^2+\gamma(x_3)\dx_3^2.
\]
Now, if $x_3 \mapsto \zeta(x_3)$ is a primitive of $\frac{-\gamma(x_3)}{2 \beta(x_3)}$, a change of coordinates
\[
(x_1,x_2,x_3) \mapsto (x_1+\zeta(x_3),x_2,x_3)
\]
shows that $h'$ is locally isomorphic to
$-2\beta(x_3)\dx_1\dx_3+\beta(x_3)\dx_2^2$, hence is conformally~flat.
\end{proof}

\section{Geometry on Einstein's universe}\label{sec.geomeinstein}

Lorentz conformally flat structures in dimension $n=3$ are examples of $(G,X)$-structures in the sense of Thurston. In particular, there is a universal space among those structures, called Einstein's universe $\ein$, such that if $(M,g)$ is Lorentz and conformally flat, there exists a conformal immersion $\delta: \tilde{M} \to \ein$, which is equivariant under a representation of $\pi_1(M)$ into $\Conf(\ein)$ (see Section~\ref{sec.developing} below). The proof of Theorem~\ref{thm.topological} for manifolds $(M,g)$ satisfying hypotheses of Theorem~\ref{thm.conf.flat} and with a noncompact isometry group, will rely in a crucial way on the study of this developing map $\delta$. This study will be carried over in the next Section~\ref{sec.global}, and it will require a deeper knowledge of the geometry of $\ein$. That's why we dedicate the present section to studying $\ein$ in more details. The reader eager to learn more about the geometry of $\ein$ is referred to~\cite{charlesthese} or~\cite{primer}.


\subsection{Basics on Einstein's universe}\label{sec.description.einstein}
Einstein's universe is the Lorentz analogue of the Riemannian conformal sphere. We recall its construction, sticking to dimension~$3$, which is the relevant one for our purpose.

Let $\RR^{2,3}$ be the space $\RR^{5}$ endowed with the quadratic form
\[
Q^{2,3}(x_0, \ldots, x_{4}) = 2x_0x_{4} + 2x_{1}x_{3} + x_2^2
\]
We consider the null cone
\[
\mathcal{N}^{2,3} = \{ x \in \RR^{2,3} \ | \ Q^{2,3}(x)=0 \}
\]
and denote by $\mathcal{\widehat{N}}^{2,3}$ the cone $\mathcal{N}^{2,3}$ with the origin removed. The projectivization ${\BP}(\widehat{\mathcal{N}}^{2,3})$ is a smooth submanifold of $\R\BP^{4}$, and inherits from the pseudo-Riemannian structure of $\RR^{2,3}$ a Lorentz conformal class (more details can be found in~\cite{charlesthese}, \cite{primer}). We call the $3$-dimensional \emph{Einstein universe}, denoted $\ein$ this compact manifold ${\BP}(\widehat{\mathcal{N}}^{2,3})$ with this conformal structure. One can check that a $2$-fold cover of $\ein$ is conformally diffeomorphic to the product $({\bf S}^1 \times {\bf S}^{2},-g_{{\bf S}^1} \oplus g_{{\bf S}^2})$.

The orthogonal group of $Q^{2,3}$, isomorphic to $\mbox{O}(2,3)$, acts naturally on the $4$-dimensional projective space, preserving $\ein$ and its conformal structure. It turns out (see Theorem~\ref{th.liouville} below) that $\mbox{PO}(2,3)$ is the full conformal group of $\ein$. Observe that $\ein$ is homogeneous under the action of $\mbox{PO}(2,3)$.

\subsubsection{Photons and lightcones}
It is a remarkable fact of Lorentz geometry that all the metrics of a given conformal class have the same lightlike geodesics (as sets but not as parametrized curves). In the case of Einstein's universe, the lightlike geodesics are the projections on $\ein$ of totally isotropic 2-planes $P \subset \RR^{2,3}$ (namely planes $P$ on which $Q^{2,3}$ vanishes identically). We will rather use the term \emph{photon} for the lightlike geodesics of Einstein's universe. Observe that all photons of $\ein$ are simple closed curves.

Given a point $p$ in $\ein$, the \emph{lightcone with vertex $p$}, denoted by $C(p)$, is the union of all photons containing $p$. If $p \in \ein$ is the projection of $u \in \mathcal{\widehat{N}}^{2,3}$, the lightcone $C(p)$ is just $\BP({ u}^{\bot} \cap \mathcal{\widehat{N}}^{2,3})$. The lightcone $C(p)$ is singular (from the differentiable viewpoint) at its vertex $p$, and $C(p) \backslash \{p \}$ is topologically a cylinder. The entire cone $C(p)$ has the topology of a $2$-torus pinched at $p$.


\subsubsection{Stereographic projection}\label{sec.stereographic}
There is for $\ein$ a generalized notion of stereographic projection, which shows that $\ein$ is a conformal compactification of the Minkowski space.


Let us call $\RR^{1,2}$ the space $\RR^3$ endowed with the quadratic form $Q^{1,2}(x,x)=2x_1x_3+x_2^2$. Consider $\varphi : \RR^{1,2} \to \ein$ given in projective coordinates of $\BP(\RR^{2,3})$ by
\begin{equation}\label{eq.stereographic}
\varphi: x=(x_1,x_2,x_3) \mapsto \left[-\frac{1}{2}Q^{1,2}(x,x): x_1 : x_2 : x_3 : 1\right]
\end{equation}
Then $\varphi$ is a conformal embedding of $\RR^{1,2}$ into $\ein$, called the inverse \emph{stereographic projection } with respect to $p_0:=[e_0]$. The image $\varphi(\RR^{1,2})$ is a dense open set of $\ein$ with boundary the lightcone $C(p_0)$. Observe that this proves the fact (rather hard to visualize): The complement of a lightcone $C(p)$ in $\ein$ is connected.

\subsubsection{Developing conformally flat structures into Einstein's universe}\label{sec.developing}
It is a standard fact that Einstein's universe satisfies an analogue of the classical Liouville's theorem on the sphere (see for instance~\cite[Chapter~7, Corollary~3.5]{sharpe} for a proof which generalizes to all signatures $(p,q)$). Namely:

\begin{theo}[Liouville's Theorem for $\ein$]\label{th.liouville}
Let $U \subset \ein$ be a connected nonempty open set. Let $f: U \to \ein$ be a conformal immersion. Then $f$ is the restriction to $U$ of an unique element of $\mbox{PO}(2,3)$.
\end{theo}

The existence of the stereographic projection~\eqref{eq.stereographic}, and the transitivity of the action of $\mbox{PO}(2,3)$ on $\ein$ shows that $\ein$ is conformally flat. Liouville's Theorem~\ref{th.liouville} shows that any $3$-dimensional, conformally flat Lorentz structure $(M,g)$ is actually a $(\mbox{PO}(2,3),\ein)$-structure, in the sense of Thurston.

As a consequence, for every conformally flat Lorentz structure $(M,g)$, there exists a conformal immersion
\[
\delta: (\tm,\tilde{g}) \to \ein
\]
called \emph{the developing map} of the structure. Here, $\tm$ is the universal cover of the manifold $M$, and $\tilde{g}$ is the lifted metric. This developing map comes with a \emph{holonomy morphism} $\rho: {\mbox{Conf}(\tm,\tilde{g})} \to \mbox{PO}(2,3)$ satisfying the equivariance relation:
\begin{equation}\label{eq.equiholonomy}
\delta \circ h = \rho(h) \circ \delta
\end{equation}
available for every $h \in {\mbox{Conf}(\tm,\tilde{g})}$. Notice that the terminology \emph{holonomy morphism} is usually used for the restriction of $\rho$ above to the group $\pi_1(M)$, seen as the group of deck transformations of $\tilde{M}$. However, we will really need the extension of $\rho $ to ${\mbox{Conf}(\tm,\tilde{g})}$ in the following.



\subsection{More geometry on \texorpdfstring{$\ein$}{ein}}\label{sec.more.geometry}

\subsubsection{The foliation \texorpdfstring{$\calfd$}{calfd}}\label{sec.foliation-einstein}
We refer here to the notations introduced in Section~\ref{sec.description.einstein}. Let $P$ be the plane in $\RR^{2,3}$ spanned by the vectors $e_0$ and $e_1$. The form $Q^{2,3}$ vanishes identically on $P$, hence the projection of $P$ on $\ein$ defines a photon that we will denote by $\Delta$. The open subset obtained by removing $\Delta$ to $\ein$ will be called $\Omega_{\D}$.

Given a point $p \in \D$, we consider the lightcone $C(p)$ with vertex $p$. Since $\D$ is a photon, we have $\D \subset C(p)$. Now, the intersection of $C(p)$ with $\Omega_{\D}$, namely $C(p) \setminus \Delta$ is a lightlike hypersurface of $\Omega_{\D}$, diffeomorphic to a plane. We call it $F_{\D}(p)$. We now make the observation that in $\ein$, there is no nontrivial lightlike triangle, namely if two photons $\D_1$ and $\D_2$ intersect $\Delta$ transversely at two distinct points, then $\D_1 \cap \D_2 = \emptyset$. This is the geometric counterpart of the following algebraic fact: In $\RR^{2,3}$, there are no $3$-dimensional spaces on which $Q^{2,3}$ vanish identically. It follows that if $p \not= p'$ are points of $\D$, $C(p) \cap C(p')=\Delta$, or in other words $F_{\Delta}(p) \cap F_{\Delta}(p^{\prime}) = \emptyset$. This shows that $\{ F_{\Delta}(p) \}_{p \in \D}$ are the leaves of a codimension~$1$ lightlike foliation of $\Omega_{\D}$, that we will call $\calfd$. Actually, there is a smooth submersion $\pi_{\Delta} : \Omega_{\Delta} \to \Delta$, which to a point $x \in \Omega_{\Delta}$ associates $p=\pi_{\Delta}(x) \in \Delta$ such that $x \in C(p)$. The fibers of $\pi_{\Delta}$ are precisely the leaves of $\calfd$, and the space of leaves is naturally identified with $\Delta$. For $x \in \Omega_{\Delta}$, we will adopt the notation $F_{\Delta}(x)$ for the leaf of $\calfd$ containing $x$.

\subsubsection{Symmetries of the foliation \texorpdfstring{$\calfd$}{calfd}}\label{sec.symetries}
Let us call $G_{\Delta}$ the stabilizer of $\D$ in $\PO$. Obviously, $G_{\Delta}$ preserves $\Omega_{\Delta}$ and the foliation $\calfd$.

It is readily checked that this group is a semi-direct product
\[
G_{\Delta} \simeq \PGL(2,\RR) \ltimes N,
\]
where the group $N$ is isomorphic to the $3$-dimensional
Heisenberg group $\Heis(3)$, and given in $\PO$ by the matrices:
\begin{equation}\label{eq.formule.heisenberg}
N(x,y,z):=
\begin{pmatrix}
1&0&-x&\, \, {-(z+xy)}&\, -\frac{x^2}{2}\\
0&1&-y&-\frac{y^2}{2}&z\\
0&0&1&y&x\\
0&0&0&1&0\\
0&0&0&0&1\\
\end{pmatrix} \ \ x,y,z \in \RR.
\end{equation}

The factor $\PGL(2,\RR)$ is the subgroup of $\PO$ corresponding to matrices:
\begin{equation}
R_A:= \label{eq.formule.pgl2}
\begin{pmatrix}
A &0&0\\
0&1&0\\
0&0&\frac{A}{\operatorname{det}(A)}\\
\end{pmatrix}  \ \ \det(A)=\pm 1.
\end{equation}

Observe that $\Delta$ being obtained as the projectivization of a null plane of $\RR^{2,3}$, it is naturally identified with ${\mathbb{RP}^1}$. The action of $G_{\Delta}$ on the space of leaves of $\calfd$ corresponds to the projective action of the factor $\PGL(2,\RR)$ on $\Delta$. The subgroup $S_{\Delta} \subset G_{\Delta}$ which preserves individually all the leaves of $\calfd$ is a semi-direct product
\[
S_{\Delta} \simeq \RR_+^* \ltimes N,
\]
where the factor $\RR_+^*$ corresponds to matrices:
\begin{equation}
R_{\lambda}:= \label{eq.facteurR}
\begin{pmatrix}
\lambda&0& & & \\
0& \lambda & & & \\
& & 1 & & \\
& & & \frac{1}{\lambda}&0\\
& & & 0& \frac{1}{\lambda}\\
\end{pmatrix}, \ \lambda \in \RR_+^*.
\end{equation}


Let us end this algebraic parenthesis by giving more details about the action of the group $N$. Obviously, $N$ fixes the point $p_0=[e_0] \in \ein$, hence if we perform a stereographic projection given by formula~\eqref{eq.stereographic}, the group $N$ becomes a subgroup of conformal transformations of $\RR^{1,2}$. These transformations are affine, given by
\begin{equation}\label{eq.formuleaffine}
N(x,y,z)=
\begin{pmatrix}
1&-y&-\frac{y^2}{2}\\ 0&1&y\\0&0&1\\
\end{pmatrix} + 
\begin{pmatrix} z\\ x\\ 0 \\
\end{pmatrix}
\end{equation}

Inside the group $N$, there is a $2$-dimensional subgroup of translations, denoted $T$, comprising all transformations of the form
\[
T(x,z):=\Id + 
\begin{pmatrix} z\\ x\\ 0 \\
\end{pmatrix}, \ x,z \in \RR.
\]
In $\PO$, such transformations take the matricial form:
\[
T(x,z)=
\begin{pmatrix}
1&0&-x&\, {-z}& -\frac{x^2}{2}\\
0&1&0&0&z\\
0&0&1&0&x\\
0&0&0&1&0\\
0&0&0&0&1\\
\end{pmatrix}.
\]

From this matricial representation, it is straightforward to check the following

\begin{fact}\label{fact.action}\ \\*[-1.3em]
\begin{enumerate}
\item \label{7.2.1} The set of fixed points for the action of the group $N$ (resp. $T$) on $\ein$ is exactly $\Delta$.
\item \label{7.2.2} For every $x \in \Omega_{\Delta}$, the $N$-orbit of $x$ is the leaf $F_{\Delta}(x)$
\item \label{7.2.3} The action of $T$ is free on $\Omega_{\Delta} \setminus F_{\Delta}(p_0)$, and orbits of $T$ on this open set coincide with leaves of $\calfd$.
\item \label{7.2.4} On $ F_{\Delta}(p_0)$, orbits of $T$ are $1$-dimensional and coincide with the photons of $C(p_0)$, with $p_0$ removed.
\end{enumerate}
\end{fact}

In the rest of the paper, we will adopt the notations $\lieg_{\Delta}, \lies_{\Delta}, \lien, \mathfrak{t}$ for the Lie subalgebras of $\oo(2,3)$ corresponding to the groups $G_{\Delta}, S_{\Delta}, N, T$.


\subsection{Standard Heisenberg algebras in \texorpdfstring{$\oo(2,3)$}{oo(2,3)}}\label{sec.standard}

The Lie group $N$ admits a Lie algebra $\lien \subset \oo(2,3)$ that will be called the \emph{standard Heisenberg algebra} of $\oo(2,3)$.

It is not true that all subalgebras of $\oo(2,3)$ which are isomorphic to $\heis(3)$ are conjugated to the standard algebra $\lien$. There is however the following useful characterization:

\begin{lemm}\label{lem.standard}
Let $\lieh \subset \oo(2,3)$ be a Lie subalgebra isomorphic to $\mathfrak{heis}(3)$, and $H \subset \PO$ the corresponding connected Lie subgroup. Assume there exists a nonempty open set of $\ein$ where the orbits of $H$ are $2$-dimensional and lightlike. Then $\lieh$ is conjugated in $\PO$ to the standard Heisenberg algebra $\lien$.
\end{lemm}


\begin{proof}
As any solvable Lie subalgebra of $\oo(2,3)$, $\lieh$ must leave invariant a line $\RR.v$ or a 2-plane $P$ in $\RR^{2,3}$. Such a vector $v$ can not be timelike or spacelike, otherwise the decomposition $\RR^{2,3}=\RR.v \oplus v^{\perp}$ would lead to an embedding of $\lieh$ in one of the Lie algebras $\RR \oplus \oo(1,3)$ or $\RR \oplus \oo(2,2) \simeq \RR \oplus \mathfrak{sl}(2,\RR) \oplus \mathfrak{sl}(2,\RR)$. But none of those algebras contains a subalgebra isomorphic to $\heis(3)$. Similarly, $P$ can not be of signature $(+,+)$, $(+,-)$ or $(-,-)$, otherwise the decomposition $\RR^{2,3}=P \oplus P^{\perp}$ would lead to an embedding of $\lieh$ into $\oo(2) \oplus \oo(2,1) \simeq \RR \oplus \oo(1,2)$, $\oo(1,1) \oplus \oo(1,2) \simeq \RR \oplus \oo(1,2)$ or $\oo(2) \oplus \oo(3) \simeq \RR \oplus \oo(3)$. One checks as above that this is not possible. The only possibilities are then:
\begin{enumerate}\alphenumi
\item \label{7.3.1}The vector $v$ is lightlike or $P$ has signature $(0,+)$ (resp. $(0,-)$). This means that $H$ has a global fixed point in $\ein$, that we can assume to be $p_0$ after conjugating within $\PO$.
\item \label{7.3.2}The form $Q^{2,3}$ vanishes identically on $P$, in which case $H$ has an invariant photon that we can assume to be $\Delta$.
\end{enumerate}



We first deal with case $a)$. After considering a stereographic projection of pole $p_0$, $\lieh$ becomes a subalgebra of $\mathfrak{Conf}(\RR^{1,2}) \simeq (\RR \oplus \oo(1,2)) \ltimes \RR^3$. Here the normal subalgebra $\RR^3$ integrates into the subgroup of translations. Let us consider the projection $\pi: (\RR \oplus \oo(1,2)) \ltimes \RR^3 \to \oo(1,2)$. Since $\oo(1,2)$ does not have any subalgebra isomorphic to $\heis(3)$ or $\RR^2$, the rank of $\pi_{|\lieh}$ is $0$ or $1$. Because $\RR \ltimes \RR^3$ (with $\RR$ acting by homothetic transformations on $\RR^3$) does not contain a copy of $\heis(3)$, this rank is actually $1$, hence the kernel of $\pi_{|\lieh}$, denoted $\liea$, has dimension~$2$ in $\lieh$, hence is abelian. The only subalgebras isomorphic to $\RR^2$ in $\RR \ltimes \RR^3$ are actually contained in $\RR^3$.

Our hypothesis on the orbits of the group $H$ implies that the translation vectors in $\liea$ span a lightlike plane, hence after conjugating within $\Conf(\RR^{1,2})$, we can assume $\liea=\mathfrak{t}$, where $\mathfrak{t}$ was introduced at the end of Section~\ref{sec.symetries}.


The first point of Fact~\ref{fact.action} implies that since $H$ centralizes $\mathfrak{t}$, $H \subset G_{\Delta}$. The hypothesis on the orbits of $H$ says that on some open set, $H$-orbits and $T$-orbits coincide. Points $3$ and $4$ of Fact~\ref{fact.action} imply that the action of $H$ on $\Delta$ is trivial on some nonempty open set, hence trivial. This yields $H \subset S_{\Delta}$. Because the normalizer of $\mathfrak{t}$ in $S_{\Delta}$ is $N$, we finally get $H=N$, and the proof is completed in this case.



Consider now case $b)$. Because $H$ leaves $\Delta$ invariant, $H$ is a subgroup of $G_{\Delta}$. As above, we can look at the morphism
\[
\pi: \lieg_{\Delta} \simeq (\RR \oplus \mathfrak{sl}(2,\RR)) \ltimes \lien \to \mathfrak{sl}(2,\RR).
\]
The same arguments as above show that the kernel of $\pi_{| \lieh}$ is a $2$-dimensional abelian Lie subalgebra $\liea \subset \lieh$. Observe that $\liea \subset \lies_{\Delta}$, and the only $2$-dimensional abelian subalgebras of $\lies_{\Delta}$ are contained in $\lien$. After conjugating within $G_{\Delta}$, we can ensure $\liea={\mathfrak t}$. We then finish the proof as in the first case.
\end{proof}



\section{The global geometry of manifolds without hyperbolic components}\label{sec.global}
This section is devoted to establishing Theorem~\ref{thm.topological} in the only remaining case to be studied, namely that of closed $3$-dimensional Lorentz manifolds $(M,g)$ which are not locally homogeneous, such that $\mint$ does not admit any hyperbolic component, and with a noncompact isometry group $\Iso(M,g)$. By Theorem~\ref{thm.conf.flat}, those manifolds are conformally flat.

What we will really show in this section is:
\begin{theoreme}\label{th.theoplat}
Let $(M,g)$ be a closed, orientable and time-orientable, $3$-dimensional Lorentz manifold, such that $\Iso(M,g)$ is noncompact. We assume that $(M,g)$ is not locally homogeneous, and that $\mint$ does not admit any hyperbolic component. Then:
\begin{enumerate}
\item \label{8.1.1} The manifold $M$ is homeomorphic to a $3$-torus, or a parabolic torus bundle~${\bf T}_A^3$.
\item \label{8.1.2} There exists a metric $g'=e^{2\sigma}g$ in the conformal class of $g$ which is flat, and which is preserved by $\Iso(M,g)$.
\item \label{8.1.3} There exists a smooth, positive, periodic function $a:\RR \to (0,\infty)$ such that the universal cover $(\tilde{M},{\tilde{g}})$ is isometric to $\RR^3$ endowed with the metric
\[
{\tilde{g}}=a(v)(\dt^2+2\du\dv).
\]
\item \label{8.1.4} There is an isometric action of $\Heis$ on $(\tm,\tilde{g})$.
\end{enumerate}
\end{theoreme}


This result clearly implies Theorem~\ref{thm.topological} in the case under study. Its proof will be the aim of Sections~\ref{sec.approximately} to~\ref{sec.conclusion} below. In all those sections, $(M,g)$ satisfies the assumptions of Theorem~\ref{th.theoplat}.

\subsection{Approximately stable foliation on \texorpdfstring{${M}$}{M}}\label{sec.approximately}
So far, we saw that $(M,g)$ is an aggregate of (possibly infinitely many) components, the local geometry of which we understand fairly well. But we need a global object which allows one to understand how those components fit together. This global object turns out to be a foliation provided by the noncompactness of $\Iso(M,g)$ as follows.



Consider a sequence $(f_n)$ in $\Iso(M,g)$ which tends to infinity, and call $AS(f_n)$ the subset of $TM$ comprising all vectors $v \in TM$ for which there exists a sequence $(v_n)$ in $TM$ converging to $v$, such that $|Df_n(v_n)|$ is bounded (where $|.|$ is the norm associated to an auxiliary Riemannian metric on $M$). In~\cite{zeghibisom1}, A. Zeghib proved the following result :

\begin{theo}[{\cite[Theorem~1.2]{zeghibisom1}}]\label{theo.approximatively}
Let $(M,g)$ be a closed Lorentz manifold, and $(f_n)$ a sequence of $\Iso(M,g)$ tending to infinity. Replacing if necessary $(f_n)$ by a subsequence, the set $AS(f_n)$ is a codimension~$1$, lightlike, Lipschitz distribution in $TM$, which integrates into a codimension~$1$, totally geodesic, lightlike foliation.
\end{theo}

The foliation given by Theorem~\ref{theo.approximatively} is called \emph{the approximately stable foliation of~$(f_n)$}.

In the particular case of a $3$-dimensional manifold, codimension~$1$, totally geodesic, lightlike foliations have very nice properties that were studied by A. Zeghib in~\cite{zeghibfoliation}. He proved in particular:

\begin{theo}[{\cite[Theorem~11]{zeghibfoliation}}]\label{th.zeghibfol}
Let $(M,g)$ be a $3$-dimensional closed Lorentz manifold. Let $\calf$ be a $C^0$, codimension~$1$, totally geodesic, lightlike foliation of $M$. Then:
\begin{enumerate}
\item \label{8.3.1}A leaf of $\calf$ is homeomorphic to a plane, a cylinder or a torus.
\item \label{8.3.2}The foliation $\calf$ has no vanishing cycles.
\end{enumerate}
\end{theo}

A consequence of the non-existence of vanishing cycles is that loops of a leaf $F$ representing a nontrivial element in $\pi_1(F)$ also represent a nontrivial element in $\pi_1(M)$.



We now choose a sequence $(f_n)$ tending to infinity in $\Iso(M,g)$, and after considering a suitable subsequence, we denote the approximatively stable foliation of $(f_n)$ by $\calf$. By Theorem~\ref{th.zeghibfol}, the leaves of $\calf$ are planes, cylinders or tori. Our main aim, and a decisive step to prove Theorem~\ref{th.theoplat} will be to show that \emph{all leaves of $\calf$ are tori}, yielding the torus bundle structure of $M$. It will be convenient in the sequel to consider the lift of $\calf$ to the universal cover $\tm$. We will call $\tcalf$ this lifted foliation.



\subsection{Fixing a component \texorpdfstring{$\calm$}{calm} and an adapted developing map}\label{sec.adapted}
We now fix, until the end of Section~\ref{sec.global}, a parabolic component $\calm \subset \mint$ \emph{which is not locally homogeneous}. We lift this component to the universal cover $\tm$ of $M$, and consider $\tcalm$ a connected component of this lift (that will be also fixed once and for all in the following). Recall the developing map $\delta: \tilde{M} \to \ein$ and the associated holonomy morphism $\rho: \Conf(\tilde{M}, \tilde{g}) \to \PO$. We pick $x_0 \in \tcalm$ and $U_0$ a $1$-connected neighbourhood of $x_0$ on which the developing map $\delta$ is injective. If $U_0$ is chosen small enough, the Lie algebra $\kil(U_0)$ of Killing fields on $U_0$ coincides with $\kiloc(x_0)$. Einstein's universe $\ein$ satisfies a generalization of Liouville's theorem (Theorem~\ref{th.liouville}): Any conformal Killing field defined on some connected open set of $\ein$ is the restriction of a global one. Thus the algebra $\delta_*(\kil(U_0))$ is a subalgebra of $\oo(2,3)$ isomorphic to the $3$-dimensional Heisenberg algebra. The pseudo-orbits of $\delta_*(\kil(U_0))$ on $\delta(U_0)$ are $2$-dimensional and lightlike by the second point of Proposition~\ref{prop.heis}. Lemma~\ref{lem.standard} applies and says that post-composing $\delta$ by an element of $\PO$, we may assume $\delta_*(\kil(U_0))=\lien$. We will now work with a developing map $\delta$ having this property, and say that \emph{$\delta$ is adapted to $\calm$}. We will consider the associated holonomy morphism $\rho: \Conf(\tilde{M}, \tilde{g}) \to \PO$ satisfying $\delta \circ h = \rho(h) \circ \delta$ for all $h \in \Conf(\tilde{M}, \tilde{g})$.

\subsection{Leaves of \texorpdfstring{$\calf$}{calf} coincide with \texorpdfstring{$\kiloc$}{kiloc}-orbits on \texorpdfstring{$\calm$}{calm}}\label{sec.calfetkill}
The aim of this section is to show:

\begin{prop}\label{prop.coincidencecalf}
The $\kiloc$-orbits in the component $\calm$ coincide with leaves of $\calf$. In particular, $\calm$ is saturated by leaves of $\calf$.
\end{prop}

Observe that by this proposition, the trace of the foliation ${\calf}$ on parabolic components which are not locally homogeneous, actually does not depend on the\linebreak sequence~$(f_n)$.

\subsubsection{The pullback foliation \texorpdfstring{$\tcalfd$}{tcalfd} and its geometric properties}\label{sec.pulbackfolia}

We consider the developing map $\delta : \tm \to \ein$, which we recall is adapted to $\calm$, and take the pullback by $\delta$ of the foliation ${\calf}_{\Delta}$ defined in Section~\ref{sec.foliation-einstein}. We get in this way a (singular) foliation $\tcalfd$ on $\tm$. Actually, $\tcalfd$ is a genuine foliation by lightlike hypersurfaces on the open set $\tilde{\Omega}_{\Delta}=\delta^{-1}(\Omega_{\Delta})$. Singularities occur on the complement of $\tilde{\Omega}_{\Delta}$ in $\tm$, namely $\tilde{\Delta}:=\delta^{-1}(\Delta)$. This singular set is either empty (in which case $\tcalfd$ is a regular foliation on $\tm$), or a $1$-dimensional lightlike manifold.

Let us emphasize the fact that a priori, we don't have any invariance property for $\tcalfd$ under the action of the fundamental group $\pi_1(M)$. In particular, there is no reason for $\tcalfd$ to define any foliation on $M$.

In the following, we will identify $\oo(2,3)$ with the Lie algebra of conformal Killing fields of $\ein$ (see Theorem~\ref{th.liouville}). We can pull back the vector fields of the Lie algebra $\lien$ by the developing map $\delta:\tm \to \ein$, getting a Lie algebra $\tlien$ of conformal Killing fields on $(\tm,{\tilde{g}})$. By Fact~\ref{fact.action}, the pseudo-orbits of $\tlien$ coincide with the leaves of $\tcalfd$.

\subsubsection{Foliation \texorpdfstring{$\tcalfd$}{tcalfd} and \texorpdfstring{$\kiloc$}{kiloc}-orbits}\label{sec.folorbits}
A first important feature of the foliation $\tcalfd$ is its relation to the $\kiloc$-orbits in~$\tcalm$.


\begin{lemm}\label{lem.coincidence}
The restriction to $\tcalm$ of any vector field of $\tlien$ is a Killing field for $\tilde{g}$. Conversely, any local Killing field defined on some open set $U \subset \tcalm$ is the restriction of a vector field in $\tlien$.
\end{lemm}

\begin{proof}
Recall the point $x_0 \in \tcalm$ and the $1$-connected open subset $U_0$ introduced in Section~\ref{sec.adapted}. By the fact that our developing map $\delta$ is adapted to $\calm$, any Killing field on $U_0$ is the restriction of a vector field of $\tlien$. Since $\tlien$ and $\kil(U_0)$ have same dimension, the restriction to $U_0$ of \emph{any} vector field of $ \tlien$ must be Killing. Let $X$ be a vector field of $\tlien$, and let us call $Y=X_{|U_0}$. Let us pick an arbitrary $y \in \tcalm$, and draw a simple curve $\gamma$ joining $y$ to $x_0$ inside $\tcalm$. Let us consider $V$ a \linebreak$1$-connected open neighbourhood of $\gamma$ contained in $\tcalm$ and containing $U_0$. Because the dimension of $\kiloc(z)$ is constant on $\tcalm$, the vector field $Y$ can be extended by analytic continuation to a Killing field for $\tilde{g}$ (still denoted $Y$) defined on $V$. But now, $Y$ and $X_{|V}$ are two conformal Killing fields on $V$, which coincide on $U_0$. They must then coincide on $V$, showing that $X$ is Killing for $\tilde{g}$ in a neighbourhood of~$y$. We have thus proved that the restriction of $X$ to $\tcalm$ is Killing. A dimensional argument as above shows that conversely, a Killing field defined on some connected open subset of $\tcalm$ is the restriction of a field in $\tlien$.
\end{proof}


\begin{coro}\label{coro.nonsingulier}
The component $\tcalm$ is contained in $\tilde{\Omega}_{\Delta}$.
\end{coro}

\begin{proof}
Points of $\tilde{\Delta}$ are singularities for the vector fields of $\tlien$. Hence if a point $x \in \tilde{\Delta}$ belongs to $\tcalm$, Lemma~\ref{lem.coincidence} will provide a Lie subalgebra of Killing fields vanishing at $x$ and isomorphic to ${\mathfrak{heis}}(3)$. The isotropy representation then yields an embedding of Lie algebras ${\mathfrak{heis}}(3) \to \oo(1,2)$. This is impossible.
\end{proof}


We conclude this paragraph with the following important lemma.


\begin{lemm}\label{lem.saturation1}
Let $x$ be a point of $\tcalm$, and $\tilde{F}_{\Delta}(x)$ the leaf of $\tcalfd$ through $x$. Then $\tilde{F}_{\Delta}(x)$ is contained in $ \tcalm$, and coincides with the $\kiloc$-orbit of $x$.
\end{lemm}

\begin{proof}
Let us consider a leaf $\tFd$ having a nonempty intersection with $\tcalm$.\linebreak Assume for a contradiction that $V=\tFd \cap \tcalm$ is not all of $\tFd$. It means that $V$ is an open subset of $\tFd$ having a nontrivial boundary $\partial{V}$ inside $\tFd$. Of course, $\partial V \subset \partial \tcalm$ (this last boundary is taken in $\tm$). Since $\tFd$ is a pseudo orbit of $\tlien$, it is easy to show that there exists $y \in \partial{V}$, a vector field $X \in \tlien$ and a point $x \in {V}$ such that the local orbit $t \mapsto \phi_{X}^t.x$ is defined on $[0,1]$, $\phi_{X}^t.x$ belongs to ${V}$ for $t \in [0,1/2)$ but $\phi_{X}^{1/2}.x \in \partial {V}$. We denote by $\hat{R}$ the bundle of frames on $\tm$, and exceptionally in this proof, we adopt the notation $\hm$ for the bundle of orthonormal frames of $\tm$ (and not of $M$). The local action of $\phi_{X}^t$ lifts naturally to $\hat{R}$. We pick $\hx \in \hm$ in the fiber of $x$, and look at the orbit $t \mapsto \phi_{X}^t.\hx$ in $\hat{R}$. Because $X$ is Killing on $\tcalm$ (Lemma~\ref{lem.coincidence}), this orbit is contained in $\hm$ for $t \in [0,1/2)$, and the same is true for $t \in [0,1/2]$ because $\hm$ is closed in $\hat{R}$. We now look at the generalized curvature map $\dk : \hm \to \mathcal{W}$, and its derivative that we see as a map $D \dk : \hm \to Hom(\lieg, \mathcal{W})$. The map $t \mapsto D \dk(\phi_{X}^t.\hx)$ makes sense for $t \in [0,1/2]$, and is constant on this interval because $X$ is Killing on $\tcalm$. In particular, the kernel of $D \dk(\phi_{X}^t.\hx)$ is the same for all $t \in [0,1/2]$, hence the rank of $\dk$ is the same at $\hx$ and at $\phi_{X}^{1/2}.\hx$. We get that the rank of $\dk$ at $\phi_{X}^{1/2}.x$ is $3$, but we already observed in the proof of Lemma~\ref{lem.rank3}, that all points where $\dk$ has rank $3$ are contained in $\tm^{\intl}$. We infer $\phi_{X}^{1/2}.x \in \tm^{\intl}$, contradicting $\phi_{X}^{1/2}.x \in \partial \calm$.

The last part of the lemma follows easily. Lemma~\ref{lem.coincidence}, together with Corollary~\ref{coro.nonsingulier} ensures that for every $x \in \tcalm$, the $\kiloc$-orbit of $x$ coincides with $\tFd(x) \cap \tcalm$. But $\tFd(x) \cap \tcalm=\tFd(x)$ by the first part of the proof.
\end{proof}



\subsubsection{Proof of Proposition~\ref{prop.coincidencecalf}}

We keep the notations of the previous paragraph. We also lift the foliation $\calf$ to a foliation $\tcalf$ on the universal cover $\tm$. For each $x \in \tm$, we denote by $\tF(x)$ the leaf of $\tcalf$ containing $x$.

Thanks to Lemma~\ref{lem.saturation1}, Proposition~\ref{prop.coincidencecalf} will be a simple consequence of:

\begin{lemm}\label{lem.saturation2}
For every $x \in \tcalm$, one has $\tFd(x)= \tF(x)$.
\end{lemm}

This shows that the foliation $\tcalf$ which is a priori only transversally Lipschitz, is transversally smooth in restriction to $\tcalm$.


\begin{proof}
We work on $\tcalm$, and we consider the two $1$-dimensional lightlike distributions $\tcald_{\Delta}=T\tcalfd^{\perp}$ and $\tcald=T\tcalf^{\perp}$. Our aim is to show that those distributions coincide on $\tcalm$. For every $x \in \tcalm$, let us introduce the set ${\mathcal C}(x)$, comprising all lightlike directions $u \in {\mathbb P}(T_x\tm)$ such that there exists a lightlike totally geodesic hypersurface $\Sigma$ through $x$, with $T_x\Sigma^{\perp}=u$. Let us recall a key observation made in~\cite{zeghibisom2}:

\begin{lemm}[{\cite[Proposition~2.4]{zeghibisom2}}]\label{lem.sec.curvature}
If the set ${\mathcal C}(x)$ spans $T_x\tm$, then the sectional curvature at $x$ is constant.
\end{lemm}

If at some point $x$ of $\tcalm$, the directions $\tcald_{\Delta}$ and $\tcald$ do not coincide, then Lemma~\ref{lem.sec.curvature} implies that they must both be fixed by the local flow generated by the isotropy algebra ${\mathfrak Is}(x)$. But a nontrivial parabolic $1$-parameter flow in $\Omp(1,2)$ has only one invariant direction: Contradiction.


We are thus led to $\tFd(x) \subset \tF(x)$ for every $x \in \tcalm$. This inclusion can not be proper, otherwise $\tcalfd$ could be extended in a smooth way to points of $\tilde{\Delta}$.
\end{proof}



\begin{rema}\label{rem.codimension1}
The previous proof shows actually that on $\tcalm$, any lightlike, totally geodesic, codimension~$1$ foliation has to coincide with $\tcalf$.
\end{rema}


\subsection{Existence of toral leaves for \texorpdfstring{${\mathcal F}$}{F}}\label{sec.toral}
We keep going in the geometric study of the foliation $\calf$ (and simultaneously in the understanding of the developing map $\delta$), by proving the existence of one toral leaf for $\calf$.
\subsubsection{Injectivity properties of \texorpdfstring{$\delta$}{delta}}\label{sec.injectivity-leave}
We keep the notations of Section~\ref{sec.pulbackfolia}, and recall the open set $\tilde{\Omega}_{\Delta} \subset \tm$ where the foliation $\tcalfd$ is defined.
\begin{lemm}\label{lem.injective}
Let $x \in \tilde{\Omega}_{\Delta}$, and assume that the leaves $\tF(x)$ and $\tFd(x)$ coincide. Then $\delta$ is injective in restriction to $\tF(x)$.
\end{lemm}

\begin{proof}
Considering if necessary a finite cover of $M$ (which will not change $\tm$), there exists $W$ a vector field on $M$, tangent to $\calf$ and satisfying $g(W,W)=1$. We lift $W$ to a vector field $\tilde{W}$ on $\tm$, which is tangent to $\tcalf$. Notice that $\tilde{W}$ is complete. By assumption, $\tilde{W}$ is tangent to $\tF(x)$. The proof follows now closely the arguments of~\cite[Proposition~6.5]{zeghibflot}. Let us call $\tilde{\cald}$ the $1$-dimensional foliation integrating $\tcalf^{\perp}$. A fundamental remark made in~\cite[Proposition~2]{zeghibfoliation} is that because $\tF(x)$ is totally geodesic, any vector field $U$ tangent to $\tcald$ acts as a Killing field on the degenerate surface $(\tF(x),\tilde{g})$. It follows easily that if $\gamma$ and $\eta$ are two curves on $\tF(x)$ parametrized by $[0,1]$, such that $\gamma(0)$ and $\eta(0)$ belong to the same leaf of $\tcald$, $\gamma'(0)$ and $\eta'(0)$ point to the same side, and if $\gamma$ and $\eta$ have the same length with respect to $\tilde{g}$, then $\gamma(1)$ and $\eta(1)$ also belong to a same leaf of $\tcald$. Applying this remark to the integral curves of the flow $\{ \psi^t \}$ generated by $\tilde{W}$, we obtain that $\psi^t$ maps leaves of $\tilde{\cald}$ to leaves of $\tilde{\cald}$. Given $\tD_0$ a leaf of $\tcald$ in $\tF(x)$, the union ${\mathcal U}(\tD_0)=\bigcup_{t \in \RR} \psi^t(\tD_0)$ is open in $\tF(x)$, and two such open sets either coincide, or are disjoint, so that $\tF(x)={\mathcal U}(\tD_0)$. By hypothesis, the leaves $ \tFd(x)$ and $\tF(x)$ coincide, so that the developing map $\delta$ sends $ \tF(x)$ to $F_{\Delta}=F_{\Delta}(\delta(x)) \subset \ein$. Let $\gamma: I \to \tF(x)$ be an injective parametrization of the leaf of $\tcald$ through $x$. We observe that for every $s \in I$, $\delta$ is injective on the curve $t \mapsto \psi^t(\gamma(s))$, because in $F_{\Delta}$, there is no closed curve transverse to photons of $F_{\Delta}$. Also, for every $t \in \RR$, $\delta$ is injective in restriction to $s \mapsto \psi^t(\gamma(s))$, because no photon in $F_{\Delta}$ is closed. The injectivity of $\delta$ on $\tF(x)$ follows.
\end{proof}

\subsubsection{The group \texorpdfstring{$\Iso(M,g)$}{Iso(M,g)} is not a torsion group}\label{sec.notorsion}
Our noncompactness hypothesis on the group $\Iso(M,g)$ does not prevent a priori $\Iso(M,g)$ from being a torsion group. In particular, we still don't know if there exists a single element $h \in \Iso(M,g)$ such that $\{ h^k \}$ is infinite discrete. The aim of this paragraph is to show it is indeed the case, and to prove the stronger statement:

\begin{prop}\label{prop.pastorsion}
Let $F \subset \calm$ be a leaf of $\calf$ containing at least one recurrent point. Let $S_F$ be the stabilizer of $F$ in $\Iso(M,g)$. There exists $h \in S_{F}$ such that the group $\{ h^k \}$ is not relatively compact in $\Iso(M,g)$.
\end{prop}

Notice that there are examples of noncompact Lorentz surfaces admitting a noncompact isometry group which is a torsion group (see~\cite[Remarque 4.12]{mounoud}).


We recall (Proposition~\ref{prop.coincidencecalf}) that $\calm$ is saturated by the leaves of $\calf$.


The proof of Proposition~\ref{prop.pastorsion} will require the intermediate Lemmas~\ref{lem.stab-ferme} and~\ref{lem.holonomy} below. We lift $F$ to a leaf $\tF \subset \tcalm$ and call ${S}_{\tF}$ the stabilizer of $\tF$ in $\Iso(\tm,\tilde{g})$. Observe that ${S}_{\tF}$ projects surjectively on $S_F$ under the epimorphism $\Iso(\tm, \tilde{g}) \to \Iso(M,g)$.
\begin{lemm}\label{lem.stab-ferme}
For every leaf $F \subset \calm$ containing recurrent points, the groups $S_F$ and $S_{\tF}$ are closed, noncompact subgroups of $\Iso(M,g)$ and $\Iso(\tm,\tilde{g})$ respectively.
\end{lemm}

\begin{proof}
We first prove that $S_F$ is closed in $\Iso(M,g)$.
If $(f_k)$ is a sequence of $S_F$ which converges to $f_{\infty} \in \iso(M,g)$, then for $k$ very large, $f_k^{-1}f_{\infty}$ belongs to the identity component $\Iso^o(M,g)$. Because $F$ coincides with a $\kiloc$-orbit of $\calm$ (Lemma~\ref{lem.saturation2}), we thus have $f_k^{-1}f_{\infty}(F)=F$, which in turns implies $f_{\infty}(F)=f_k(f_k^{-1}f_{\infty}(F))=f_k(F)=F$.

Let us now check that $S_F$ is noncompact. By assumption on $F$, there is a recurrent point $x$ in $F$. It means that there exists a sequence $(f_k)$ tending to infinity in $\Iso(M,g)$ such that $f_k(x) \to x$. Because the $\\Is^{\loc}$-orbit of $x$ is a $2$-dimensional submanifold, the connected components of which are $\kiloc$-orbits (see Theorem~\ref{thm.integrabilite}), we get that $f_k(x) \in F$ for $k$ large enough. In particular, $S_F$ is a noncompact subgroup of $\Iso(M,g)$.

The corresponding assertions on $S_{\tF}$ are then straightforward.
\end{proof}



\begin{lemm}\label{lem.holonomy}
Let $F \subset \calm$ be a leaf of $\calf$, and $\tF$ a lift of $F$ to $\tm$.
\begin{enumerate}
\item \label{8.14.1}The holonomy morphism $\rho$ maps the group ${S}_{\tF}$ into the group $S_{\Delta}$. In particular, any element of $S_{\tF}$ leaves invariant the leaves of $\tcalf$ which are sufficiently close to $\tF$.
\item \label{8.14.2}The morphism $\rho : S_{\tF} \to S_{\Delta}$ is injective and proper.
\end{enumerate}
\end{lemm}

\begin{proof}
We heavily use the notations introduced in Section~\ref{sec.more.geometry}. We choose a transversal $I \subset \tcalm$ to the foliation $\tcalf$, that cuts $\tF$ at $x$. We assume that $I$ is small enough, so that $\delta$ sends $I$ injectively on a transversal $J$ of ${\mathcal F}_{\Delta}$. Shrinking $I$ if necessary, $J$ meets each leaf of ${\mathcal F}_{\Delta}$ at most once, so that by Lemma~\ref{lem.injective}, $I$ meets each leaf of $\tcalf$ at most once. We call $\mathcal{V}$ the open subset obtained by saturating $I$ by leaves of $\tcalf$. We recall the submersion $\pi_{\Delta}: \Omega_{\Delta} \to \Delta$ introduced in Section~\ref{sec.foliation-einstein}. By what we just said, $I$ is the space of leaves of $\mathcal{V}$, and the map $\varphi:=\pi_{\Delta} \circ \delta: I \to \Delta$ gives an identification of $I$ with $J'=\pi_{\Delta}(J)$, the space of leaves of the foliation induced by ${\mathcal F}_{\Delta}$ on $\delta({\mathcal V})$. Under this identification, the point $x$ is sent to a point $p \in J'$.

Let $\tth$ be an element of $S_{\tF}$. Its action on the space of leaves of $\tcalf$ yields a germ $\overline{h}$ of diffeomorphism of $I$ fixing $x$. The equivariance relation $\delta \circ \tth=\rho(\tth) \circ \delta$ shows that $\rho(\tth)$ permutes leaves of $\calfd$ near $\delta(\tF)$. In particular, $\rho(\tth)$ maps $J'$ to an interval of $\Delta$ containing $p$. We infer that $\rho(\tth)$ preserve $\Delta$, what yields $\rho(\tth) \in G_{\Delta}$. Moreover, denoting $l : G_{\Delta}\simeq \PGL(2,\RR) \ltimes N \to \PGL(2,\RR)$, we get the equivariance relation $\varphi \circ \overline{h}=l(\rho(\tth)) \circ \varphi$. Now, $l(\rho(\tth))$ acts as an element of $\PGL(2,\RR)$ on $\Delta$, admitting $p$ as fixed point. We know the local dynamics of a Möbius transformation around one of its fixed points: If $l(\rho(\tth))$ is nontrivial, we can choose $q \in J'$, $q \not = p$, such that $l(\rho({\tth}^k))(q)$ belongs to $J'$ for all $k \geq 0$, and $\lim_{k \to \infty}l(\rho({\tth}^k))(q)=p$. This means that if $\tF'$ is a leaf corresponding to $\varphi^{-1}(q)$, the iterates ${\tth}^k(\tF')$ will accumulate on $\tF$. But $\tF'$ is a $\kiloc$-orbit by Proposition~\ref{prop.coincidencecalf}, and closeness of the $\Is^{\loc}$-orbit of $\tF'$ in $\tm^{\intl}$ (Corollary~\ref{coro.islocorbites}) says that $\tF$ and all the $h^k(\tF')$, $k \in \NN$, belong to the same $\Is^{\loc}$-orbit. This accumulation phenomenon then contradicts the fact that $\\Is^{\loc}$-orbits are submanifolds in $\tm^{\intl}$ (see Theorem~\ref{thm.integrabilite}). We conclude that $l(\rho(\tth))$ is trivial, which implies that $\rho(\tth) \in S_{\Delta}$. Moreover, $\overline{h}$ is trivial, which means that all leaves of $\tcalf$ close to $\tF$ are left invariant by $\tth$.

We now prove the second point of the Lemma. Let $\tth \in S_{\tF}$ such that $\rho(\tth)=\id$. Equivariance relation $\delta \circ \tth = \rho(\tth) \circ \delta$, together with Lemma~\ref{lem.injective}, shows that the action of $\tth$ on $\tF$. The following fact then implies $\tth=\id$.
\begin{fact}\label{fact.injectiveproper}
Let $N$ be a Lorentz manifold and $\Sigma \subset N$ a lightlike hypersurface. denote by $S_{\Sigma}$ the stabilizer of $\Sigma$ in $\Iso(N)$. Then the restriction map $r: S_{\Sigma}\to \Homeo(\Sigma)$ is injective and proper.
\end{fact}

\begin{proof}
The proof relies on the fact that the map, which to an isometry associates its $1$-jet at a given point, is injective and proper, and that restricting elements of $\Omp(1,n-1)$ to a lightlike hyperplane is also injective and proper. Details can be found in~\cite[Proposition~3.6]{zeghibflot}, for instance.
\end{proof}

Properness of the map $\rho:S_{\tF} \to S_{\Delta}$ follows the same lines. If $(\tth_k)$ is a sequence of $S_{\tF}$ such that $\rho(\tth_k)$ is relatively compact in $S_{\Delta}$. Then $\rho(\tth_k)_{|\delta(\tF)}$ is relatively compact, hence the restriction of ${\tth_k}$ to $\tF$ is relatively compact by Lemma~\ref{lem.injective}. Fact~\ref{fact.injectiveproper} yields that $(\tth_k)$ is relatively compact in $S_{\tF}$.
\end{proof}

We can now proceed to the proof of Proposition~\ref{prop.pastorsion}. We know from Theorem~\ref{th.zeghibfol}, that the leaves of $\calf$ are discs, cylinders or tori, and there are no vanishing cycles. It means (see the comment right after Theorem~\ref{th.zeghibfol}) that the leaf $\tF$ is a disc, and the stabilizer $\Gamma_F$ of $\tF $ in $\pi_1(M)$ is either trivial, or a discrete subgroup isomorphic to $\ZZ$ or $\ZZ^2$. On the other hand, we also know that $\tilde{S}_F/\Gamma_F$ is noncompact, because of Lemma~\ref{lem.stab-ferme}.
\begin{enumerate}\alphenumi
\item \label{listprf8.15.a}\emph{Case where $F$ is a disk.} We choose a nontrivial $\tth \in S_{\tF}$. It restricts to a nontrivial transformation of $\tF$ (Fact~\ref{fact.injectiveproper}), hence $\rho(\tth)$ is a nontrivial element of $S_{\Delta}$, by Lemma~\ref{lem.holonomy}. Every nontrivial element of $S_{\Delta}$ generates an infinite discrete group. In particular, $\{\rho(\tth)^k\}_{k \in \ZZ}$ is not relatively compact in $S_{\Delta}$. Second point of Lemma~\ref{lem.holonomy} says that $\{h ^k \}$ is not relatively compact in $\Iso(\tm,\tilde{g})$. Fact~\ref{fact.injectiveproper} thus implies that $\{{\tth}^k_{|\tF}\}$ is not relatively compact in $\Homeo(\tF)$, hence the same is true for $\{h_{|F}^k \}$, because the projection $\pi: \tF\to F$ is a diffeomorphism in the case we are considering. Finally, $\{ h^k \}$ is not relatively compact in $\Iso(M,g)$.

\item \label{listprf8.15.b}\emph{Case where $F$ is a cylinder.} Because $\calf$ does not have vanishing cycles, $\Gamma_F:=S_{\tF} \cap \Gamma$ is nontrivial, generated by a single element ${\gamma}$. The automorphism group of $\Gamma_F$ is $\{ \pm 1\}$ and $\Gamma_F$ is normalized by $S_{\tF}$. Hence after considering an index $2$ subgroup of $S_{\tF}$, we may assume that $\gamma$ is centralized by all elements of $S_{\tF}$. We observe that $\rho(\gamma)$ is nontrivial by Lemma~\ref{lem.holonomy}, and consider its centralizer in $S_{\Delta}$. Two cases can then occur:
\begin{itemize}
\item \label{listprf8.15.b.1}The group $\rho(S_{\tF})$ is contained in a $1$-parameter subgroup of $S_{\Delta}$. In this case, $\rho(S_{\tF})/\langle\rho(\gamma)\rangle$ is relatively compact. This implies that $(S_F)_{|F}$ is relatively compact, hence $S_F$ is relatively compact in $\Iso(M,g)$ (again Fact~\ref{fact.injectiveproper}). This is ruled out by Lemma~\ref{lem.stab-ferme}.
\item \label{listprf8.15.b.2}If we are not in the previous case, we can find $\tth \in S_{\tF}$ such that the group generated by $\rho(\tth)$ and $\rho(\gamma)$ is discrete isomorphic to $\ZZ^2$. As above, applying second point of Lemma~\ref{lem.holonomy} and Fact~\ref{fact.injectiveproper}, one gets that $\tth$ projects to $h \in S_F$, such that $\{ h^k \}$ is infinite discrete in $\Iso(M,g)$.
\end{itemize}
\item \label{listprf8.15.c}\emph{Case where $F$ is a torus.} This time, $\Gamma_{F}$ is isomorphic to $\ZZ^2$ and generated by $\gamma_1$ and $\gamma_2$. Lemma~\ref{lem.holonomy} ensures that $\tau_1:=\rho(\gamma_1)$ and $\tau_2:=\rho(\gamma_2)$ generate a discrete subgroup of $S_{\Delta}$ isomorphic to $\ZZ^2$. Such a subgroup must be contained in $N$, and after conjugating $\rho$ within $G_{\Delta}$ (what amounts to post-compose $\delta$ by some element of $G_{\Delta}$), we have that $\langle\tau_1,\tau_2\rangle \subset T$. We must have $\rho(S_{\tF}) \subset N$ because $S_{\tF}$ normalizes $\Gamma_F$, and because $S_F$ is noncompact, $\rho(S_{\tF}) \not \subset T$. Picking $\tth \in S_{\tF}$ such that $\rho(\tth) \not \in T$, we get an element $h \in S_F$ which, by similar arguments as above, generates an infinite discrete group $\{ h^k \} \subset \Iso(M,g)$.
\end{enumerate}



\subsubsection{Existence of a toral leaf}\label{sec.toralproof}
We now consider an element $h \in \Iso(M,g)$ given by Proposition~\ref{prop.pastorsion}, namely $\{ h^k \}$ is not relatively compact in $\Iso(M,g)$. Theorem~\ref{theo.approximatively} provides an approximately stable foliation ${\mathcal F}_h$ associated to a subsequence of $\{ h^k \}$, and since all what we did before did not assume anything special on ${\mathcal F}$, we can decide that now $\calf={\mathcal F}_h$.

\begin{prop}\label{prop.existence.tori}
Every leaf $F$ of $\calf$ which is contained in $\calm$ is a torus.
\end{prop}

Let $F$ be a leaf of $\calf$ contained in $\calm$, such that almost every point of $F$ is recurrent for $\{ h^k \}$. We lift $F$ to $\tF \subset \tcalm$, and we will also assume that $\delta(\tF)$ is not contained in the leaf $F_{\Delta}(p_0)$ (see Fact~\ref{fact.action}). Observe that since $M$ is closed, Poincaré recurrence ensures that almost every point of $(M,g)$ is recurrent for $\{ h^k \}$. It follows that for almost every leaf of $\calf$, almost every point is recurrent (leaves of $\calm$ coincide with $\kiloc$-orbits hence are locally closed and transversally smooth on $\calm$. There is thus nothing tricky in disintegrating the volume form in $M$ along those leaves). Hence almost every leaf of $\calf$ in $\calm$ is an $F$ with the properties above. We claim that $F$ is a torus. To see this, we lift $h$ to an element $\tth \in S_{\tF}$, and our assumption is that the set of points in $\tF$, which are recurrent under the group $\langle \tth, \Gamma_{F}\rangle$ have full measure in $\tF$. By Lemma~\ref{lem.injective}, almost all points of $\delta(\tF)$ must be recurrent for the group $\rho(\langle \tth,\Gamma_{F}\rangle)$. We now go back to the analysis made at the end of Section~\ref{sec.notorsion}. If $F$ is not a torus, we are in cases $a)$ or $b)$ of this discussion. In case $a)$, $\Gamma_F$ is trivial. The action of $\rho(\tth)$ on $\delta(\tF)$ is conjugated to that of an affine transformation on (an open subset of) the plane (see Section~\ref{sec.more.geometry} and formula~\eqref{eq.formuleaffine}). The set of recurrent points of $\rho(\tth)$ has thus zero measure on $\delta(\tF)$. A contradiction.

In case $b)$, we saw that the group $\rho(\langle \tth,\Gamma_{F}\rangle )$ is conjugated to a lattice in the closed subgroup $T$. Since by assumption, $\delta(\tF)$ is not contained in the leaf $F_{\Delta}(p_0)$, point $3)$ of Fact~\ref{fact.action} shows that $T$ has no recurrent point on $\delta(\tF)$. We reach a new contradiction.


The arguments above show that almost all leaves $F \subset \calm$ are tori. Now, for a codimension~$1$ foliation on a closed manifold, the union of all compact leaves is itself compact (see~\cite[Chapter~II, Corollary~3.10]{godbillon}). Proposition~\ref{prop.existence.tori} follows.


\subsection{All leaves of \texorpdfstring{${\mathcal F}$}{F} are tori}

We keep the notations of the last section. We still consider $h \in \Iso(M,g)$ such that $\{ h^k \}$ is not relatively compact. We consider the approximately stable foliation $\calf$ associated to some sequence $(h^{n_k})$, with $n_k \to \infty$. Proposition~\ref{prop.existence.tori} and its proof provide us with a leaf $F_0 \subset \calm$ which is an $h$-invariant torus. We lift $F_0$ to $\tF_0 \subset \tm$, and $h$ to $\tth \in \Iso(\tm,\tilde{g})$ preserving $\tF_0$. The group $\Gamma_0=S_{\tF_0} \cap \pi_1(M)$ is discrete and isomorphic to $\ZZ^2$, generated by two elements $\gamma_1$ and $\gamma_2$. Lemma~\ref{lem.holonomy} ensures that $\tau_1:=\rho(\gamma_1)$ and $\tau_2:=\rho(\gamma_2)$ generate a discrete subgroup of $S_{\Delta}$ isomorphic to $\ZZ^2$. After conjugating by an element of $G_{\Delta}$, we may assume that $\tau_1$ and $\tau_2$ are elements of $T$.

We call in the sequel $\tilde{H}$ the subgroup generated by $\tth, \gamma_1$ and $\gamma_2$. Lemma~\ref{lem.holonomy} ensures that $\rho(\tth) \in S_{\Delta}=\RR_+^* \ltimes N$ (see Section~\ref{sec.symetries}). Elements of $S_{\Delta}$ which are not in $N$ act on $N$ with nontrivial dilation. They can not preserve any lattice in $T$. This says that because $\tth$ normalizes $\Gamma_0$, we must have $\rho(\tilde{H}) \subset N$. There are thus three elements $X,Y,Z$ in the Lie algebra $\lien$ such that $\tau_1=e^X$, $\tau_2=e^Z$ and $\rho(\tth)=e^Y$. The center of $\lien$ is contained in $\Span(X,Z)$, and we pick $Z_0 \not =0$ in this center.

We now pull back the four vector fields $X,Y,Z,Z_0$ of $\ein$ by the developing map $\delta: \tm \to \ein$. This way, we get vector fields $\tX,\tY, \tZ, \tZ_0$ on $\tm$. Observe that $\tau_1^*X=X$, $\tau_2^*Z=Z$, $\rho(\tth)^*Y=Y$, and $\rho(\tth)^*Z_0=\tau_1^*Z_0=\tau_2^*Z_0=Z_0, \ \tau_1^*Z=\tau_2^*Z=Z$ imply the relations
\begin{multline}\label{eq.vectorfields}
\gamma_1^*\tX=\tX,\ \gamma_2^*\tZ=\tZ, \ \tth^*\tY=\tY,\\
\text{ and }\ \tth^*\tZ_0=\gamma_1^*\tZ_0= \gamma_2^*\tZ_0=\tZ_0, \ \gamma_1^*\tZ= \gamma_2^*\tZ=\tZ
\end{multline}
on $\tm$. After introducing those notations, we can state what will be the last technical step of our study:

\begin{prop}\label{prop.ensembleE}
For every $x \in \tm$, we have:
\begin{enumerate}
\item \label{list8.17.1}The point $x$ belongs to $\tilde{\Omega}_{\Delta}$ and $\tilde{F}(x)=\tilde{F}_{\Delta}(x)$.
\item \label{list8.17.2}The leaf $\tilde{F}(x)=\tilde{F}_{\Delta}(x)$ is $\tilde{H}$-invariant.
\item \label{list8.17.3}The restriction of the $3$ vector fields $\tX,\tY,\tZ$ are complete on $\tilde{F}_{\Delta}(x)$, and the equalities $\phi_{\tilde{X}}^1=\gamma_1, \ \phi_{\tilde{Z}}^1=\gamma_2, \ \phi_{\tilde{Y}}^1=\tilde{h}$ hold on $\tilde{F}_{\Delta}(x)$.
\end{enumerate}
\end{prop}

The proposition will show that $\tcalf$ coincide $\tcalfd$ on $\tm$, and $\Gamma_{0}$-invariance of the leaves of $\tcalf$ will easily imply that leaves of $\calf$ are all tori. In Section~\ref{sec.conclusion} we will derive more consequences from this equality ${\mathcal E}=\tm$, and prove Theorem~\ref{th.theoplat}.

We are going to consider the set $\mathcal{E} \subset \tm$, comprising all points $x \in \tm$ satisfying the three conditions of Proposition~\ref{prop.ensembleE}, and show that $\mathcal{E}$ is nonempty, open and closed in $\tm$, yielding $\mathcal{E}=\tm$.



\subsubsection{The set \texorpdfstring{$\mathcal{E}$}{E} is nonempty}\label{sec.nonempty}
We check here that every point $x \in \tF_0$ belongs to $\mathcal{E}$. Recall $\rho: \Conf(\tm) \to \PO$ the holonomy morphism.

\begin{lemm}\label{lem.feuillecocompacte}
Let $x \in \tilde{\Omega}_{\Delta}$ such that $\tF(x)=\tFd(x)$. Assume moreover that $\tFd(x)$ is invariant by a subgroup $\Lambda \subset \pi_1(M)$, isomorphic to $\ZZ^2$ and such that $\rho(\Lambda) \subset T$. Then the map $\delta$ is a diffeomorphism from $\tFd(x)$ to $F_{\Delta}(\delta(x))$. Moreover $F_{\Delta}(\delta(x)) \not =F_{\Delta}(p_0)$.
\end{lemm}

\begin{proof}
Lemma~\ref{lem.injective} ensures that $\delta$ is a diffeomorphism from $\tFd(x)$ to an open subset $U \subset F_{\Delta}(\delta(x))$. The group $\Lambda$ is isomorphic to $\ZZ^2$ and acts properly discontinuously on the disk $\tF(x)=\tFd(x)$. By a cohomological dimension argument, the action must be cocompact. The group $\rho(\Lambda)$ is thus a lattice in $T$, and must act properly and cocompactly on $U$. Last point of Fact~\ref{fact.action} says that the action of $\rho(\Lambda)$ can not be proper on any open subset of $F_{\Delta}(p_0)$. We thus infer that $F_{\Delta}(\delta(x)) \not = F_{\Delta}(p_0)$. In particular, again by Fact~\ref{fact.action}, the action of $\rho(\Lambda)$ is proper and cocompact on $F_{\Delta}(\delta(x))$. We then must have $U=F_{\Delta}(\delta(x))$.
\end{proof}


The completeness of $\tX,\tY$ and $\tZ$ on $\tF_0$ follows from Lemma~\ref{lem.feuillecocompacte}, applied for $\Lambda=\Gamma_0$, because $X,Y,Z$ are complete on leaves of $\calfd$. The relations $\tau_1=e^X$, $\tau_2=e^Z$ and $\rho(\tth)=e^Y$ imply that the relations $\phi_{\tilde{X}}^1=\gamma_1,
\ \phi_{\tilde{Z}}^1=\gamma_2, \ \phi_{\tilde{Y}}^1=\tilde{h}$ hold on $\tF_0$. We infer that $\tF_0 \subset {\mathcal E}$.

\subsubsection{The set \texorpdfstring{$\mathcal{E}$}{E} is open}\label{sec.open}

We begin by stating a lemma that we will use repeatedly in the sequel.

\pagebreak
\begin{lemm}\label{lem.coincide}
Let $U \subset \tilde{\Omega}_{\Delta}$ be a connected open set. Let $f,g:U \to \tm$ two conformal immersions. Assume that for some $x \in \tilde{\Omega}_{\Delta}$, $\tF_{\Delta}(x) \cap U \not = \emptyset$, and that $f$ and $g$ coincide on $\tF_{\Delta}(x) \cap U$, then $f$ and $g$ coincide on $U$.
\end{lemm}

\begin{proof}
Shrinking $U$ if necessary and looking at $\delta(U) \subset \ein$, we are reduced to the situation of two transformations $g_1$ and $g_2$ of $\PO$ which coincide on some open subset of a lightcone in $\ein$. At level of linear algebra, it means that those two transformations of $\PO$ must coincide on a lightlike hyperplane of $\RR^{2,3}$. This easily implies $g_1=g_2$.
\end{proof}


Let us start with $x \in \mathcal{E}$. Vector fields $\tX,\tY$ and $\tZ$ are complete in restriction to $\tFd(x)=\tF(x)$, so given $\epsilon >0$, we can choose $U \subset \tilde{\Omega}_{\Delta}$ a small neighbourhood of $x$ such that $\phi_{\tX}^t.y$ is defined on $[-\epsilon, 1+\epsilon]$ for every $y \in U$, and for every $t \in [- \epsilon, 1+\epsilon]$ and every $y \in U$, $\phi_{\tZ}^s.\phi_{\tX}^t.y$ is defined for every $s \in [-\epsilon, 1+ \epsilon]$. Lemma~\ref{lem.coincide} says that identity $\phi_{\tX}^1.y=\gamma_1.y$ holds on $U$, because it holds on $U \cap \tFd(x)$. It follows easily from the property $\gamma_1^*\tX=\tX$ that for every $y \in U$, $\phi_{\tX}^t.y$ is defined for $t \in \RR$. Relation $\gamma_1^*\tZ=\tZ$ now implies that $\phi_{\tZ}^s.\phi_{\tX}^t.y$ makes sense for every $t \in \RR$, $s \in [- \epsilon, 1+\epsilon]$, and $y \in U$. Let us call $\mathcal{U}=\{ \phi_{\tX}^t.y \ | \ t \in \RR, y \in U\}$. This is an open set on which $\phi_{\tZ}^1$ is defined. Relation $\phi_{\tZ}^1=\gamma_2$ holds on ${\mathcal U} \cap \tFd(x)$, hence on ${\mathcal U}$ by Lemma~\ref{lem.coincide}. Together with the property $\gamma_2^* \tZ=\tZ$, this implies that $\phi_{\tZ}^s.\phi_{\tX}^t.y$ makes sense for every $y \in U$, and $s,t \in \RR$.

Now, Lemma~\ref{lem.feuillecocompacte} says that $F_{\Delta}(\delta(x))\not = F_{\Delta}(p_0)$. If $U$ was chosen small enough, $\delta(y) \not \in F_{\Delta}(p_0)$ for every $y \in U$. It follows that $(t,s) \mapsto e^{sZ}.e^{tX}.\delta(y)$ is a diffeomorphic parametrization of $F_{\Delta}(\delta(y))$. In other words, for every $y \in U$, $\{ \phi_{\tZ}^s.\phi_{\tX}^t.y \ | \ (s,t) \in \RR^2 \}$ coincides with the leaf $\tFd(y)$, and the developing map $\delta: \tFd(y) \to F_{\Delta}(\delta(y))$ is a diffeomorphism. Completeness of vector fields $\tX,\tY,\tZ$ on $\tFd(y)$ follows, because vector fields $X,Y,Z$ are complete on leaves of $\calfd$.

Moreover, Lemma~\ref{lem.coincide} says that relations $\phi_{\tX}^1=\gamma_1, \ \phi_{\tZ}^1=\gamma_2$ and $\phi_{\tY}^1=\tth$ hold on ${\mathcal U}$ because they hold on $\mathcal{U} \cap \tFd(x)$. In particular, for $y \in U$, the leaf $\tFd(y)$ is stable by $\gamma_1, \ \gamma_2$ and $\tth$, hence by $\tilde{H}$.


To conclude that $U \subset {\mathcal E}$, it remains to check that $\tFd(y)$ coincides with $\tF(y)$ for every $y \in U$. A first observation is that $\tFd(y)$ is diffeomorphic to $F_{\Delta}(\delta(y))$, hence is a disk. It follows from a cohomological dimension argument that because $\Gamma_{0} \simeq \ZZ^2$, the quotient $\Sigma=\Gamma_{0} \backslash
\tFd(y)$ is a torus in $M$. Recall from~\eqref{eq.vectorfields} the relation $\tth^*\tZ_0=\gamma_1^*\tZ_0=\gamma_2^* \tZ_0=\tZ_0$. Remember also that $\tZ_0$ is a linear combination of $\tX$ and $\tZ$, hence tangent to $\tFd(y)$. On the torus $\Sigma$, $\tZ_0$ thus induces a vector field $\overline{Z}_0$ which is $h$-invariant. Notice that $\overline{Z}_0$ is lightlike because $Z_0$ is lightlike on $\ein$, and $\overline{Z}_0$ is nonsingular because the singularities of $Z_0$ are exactly the points of $\Delta$, and $\tFd(y) \subset \tilde{\Omega}_{\Delta}$. Hence, for every $z \in \Sigma$, $\overline{Z}_0(z)^{\perp}=T_z \Sigma$. On the other hand, equality $D_zh^{n_k}(\overline{Z}_0(z))=\overline{Z}_0(h^{n_k}(z))$ shows that $\overline{Z}_0(z)$ belongs to the approximately stable distribution of $h^{n_k}$ (see the definition of this distribution in Section~\ref{sec.approximately}). The approximately stable distribution has codimension~$1$ and is lightlike, so that it coincides with $\overline{Z}_0(z)^{\perp}=T_z\Sigma$ for all $z \in \Sigma$. We conclude that $\Sigma$ is a leaf of $\calf$, which proves $\tF(y)=\tFd(y)$.



\subsubsection{The set \texorpdfstring{$\mathcal{E}$}{E} is closed}\label{sec.closed}
We consider a sequence $(x_k)$ of ${\mathcal E}$ converging to $x_{\infty} \in \tm$. The leaf $\tF(x_{\infty})$ is accumulated by the sequence of leaves $\tF(x_k)=\tF_{\Delta}(x_k)$. In particular, the vector fields $\tX,\tY,\tZ$ being tangent to $\tF_{\Delta}(x_k)$ for all $k$, they are also tangent to $\tF(x_{\infty})$. Point $2)$ of Fact~\ref{fact.action} then says that $\tF(x_{\infty}) \setminus \tilde{\Delta}$ is a union of leaves of $\tcalfd$. If the set $\tF(x_{\infty}) \cap \tilde{\Delta}$ is not empty, those leaves of $\tcalfd$ might be extended smoothly across the singular set $\tilde{\Delta}$, a contradiction. We infer that $\tF(x_{\infty}) \subset \tilde{\Omega}_{\Delta}$, and $\tF(x_{\infty})=\tF_{\Delta}(x_{\infty})$.

The union of the compact leaves of $\calf$ is a compact subset of $M$ (see~\cite[Chapter~II, Corollary~3.10]{godbillon}). Since $\calf$ has no vanishing cycles, $\tF(x_{\infty})$ is left invariant by a discrete subgroup $\Lambda_1 \subset \pi_1(M)$ which is isomorphic to $\ZZ^2$. We choose $I \subset M$, a small transversal to the foliation $\calf$ containing the point $\pi(x_{\infty})$. Following the loops of $F(\pi(x_{\infty}))$ defining $\Lambda_1$ in the neighboring leaves, we get a corresponding holonomy morphism for the leaf $F(\pi(x_{\infty}))$:
\[
\hol: \Lambda_1 \to \Diff_{\pi(x_{\infty})}^{\loc}(I).
\]
Here $\Diff_{\pi(x_{\infty})}^{\loc}(I)$ denotes the pseudo-group of local diffeomorphisms of $I$ fixing $\pi(x_{\infty})$. We refer to~\cite[Chapter~II]{godbillon} for more details on holonomy of a foliation. If $\gamma \in \Lambda_1$ and $z \in I$ is close enough to $\pi(x_{\infty})$, then ${\hol}(\gamma^2).z$ makes sense. If we have ${\hol}(\gamma^2).z \not = z$, then the pseudo-orbit ${\hol}(\gamma^n).z$ is defined for all $n \geq 0$ or all $n \leq 0$ and is infinite. In this case the leaf $F(z)$ cannot be closed. Because $F(\pi(x_k))$ is a torus for each $k \in \NN$, it follows that replacing $\Lambda_1$ by some index $2$ subgroup, we may assume that $\Lambda_1$ leaves invariant $\tF(x_k)$ for $k$ large enough.

This property also shows that $\rho(\Lambda_1)$, which is a priori not a subgroup of $G_{\Delta}$, leaves invariant infinitely many leaves of $\calfd$. Those leaves are traces on $\Omega_{\Delta}$ of lightcones in $\ein$ with vertex on $\Delta$. We infer that $\rho(\Lambda_1)$ fixes infinitely many points on $\Delta$. It follows that $\rho(\Lambda_1)$ leaves $\Delta$ invariant: $\rho(\Lambda_1) \subset G_{\Delta}$. Since a Möbius transformation fixing infinitely many points on the circle must be trivial, we actually have $\rho(\Lambda_1) \subset S_{\Delta}$.

The group generated by $\Lambda_1$ and $\Gamma_0$ is contained in $\pi_1(M)$, hence must act properly discontinuously on each $\tF(x_k)$. It follows that $\Lambda=\langle\Lambda_1,\Gamma_0\rangle$ is a discrete group isomorphic to $\ZZ^2$. In particular $\Lambda_1$ commutes with $\Gamma_0$, and $\rho(\Lambda_1) \subset T$. We can then apply Lemma~\ref{lem.feuillecocompacte} to $\Lambda$, and we get that $\delta$ is a diffeomorphism from $\tF(x_{\infty})=\tFd(x_{\infty})$ to $F_{\Delta}(x_{\infty})$, and $F_{\Delta}(\delta(x_{\infty})) \not = F_{\Delta}(p_0)$. It follows that $\tX,\ \tY$ and $\tZ$ are complete in restriction to $\tFd(x_{\infty})$.


Let $y_{\infty} \in \tF(x_{\infty})$, and let $U \subset \tilde{\Omega}_{\Delta}$ be a small neighbourhood of $y_{\infty}$. By completeness of $\tX,\ \tY$ and $\tZ$ in restriction to $\tFd(x_{\infty})$, and shrinking $U$ if necessary, the local diffeomorphisms $\phi_{\tX}^1$, $\phi_{\tY}^1$ and $\phi_{\tZ}^1$ are defined on $U$. For $k$ large, $U \cap \tFd(x_k) \not = \emptyset$, and identities $\phi_{\tX}^1=\gamma_1$, $\phi_{\tY}^1=\tth$ and $\phi_{\tZ}^1=\gamma_2$ hold on $U \cap \tFd(x_k)$. Lemma~\ref{lem.coincide} says that those identities hold on $U$. Finally $y_{\infty}$ was arbitrary in $\tFd(x_{\infty})$ so that these identities hold on $\tFd(x_{\infty})$. This proves $x_{\infty} \in {\mathcal E}$.



\subsection{Proof of Theorem~\ref{th.theoplat}}\label{sec.conclusion}
Let us draw further conclusions from Proposition~\ref{prop.ensembleE}. The coincidence of the foliations $\tcalf$ and $\tcalfd$ implies that $\tcalfd$ is $\pi_1(M)$-invariant. Moreover, the $\Gamma_0$-invariance of each leaf $\tFd$, together with Lemma~\ref{lem.feuillecocompacte} implies that $\delta$ is injective on each leaf $\tFd$, and that $\delta(\tm) \subset \Omega_{\Delta} \setminus F_{\Delta}(p_0)$.

Also, it follows from Proposition~\ref{prop.ensembleE} that $\Gamma_0$ is exactly the subgroup of $\pi_1(M)$ leaving each leaf of $\tcalf$ invariant. It follows that $\Gamma_0$ is normal in $\pi_1(M)$. We claim that $\Gamma_0$ is also normalized by $N_{\pi_1}$, the normalizer of $\pi_1(M)$ in $\Iso(\tm,\tilde{g})$. Indeed, if $f \in N_{\pi_1}$, then $f\Gamma_0 f^{-1}$ leaves each leaf of $f(\tcalf)$ invariant. Now $f(\tcalf)$ is a lightlike, totally geodesic, codimension~$1$ foliation. Remark~\ref{rem.codimension1} ensures that on any non locally homogeneous component $\tcalm$, $f(\tcalf)$ coincides with $\tcalf$. In particular, $f\Gamma_0 f^{-1}$ coincide with $\Gamma_0$ on $\tcalm$, hence $f\Gamma_0f^{-1}=\Gamma_0$.


The group $\rho(N_{\pi_1})$ normalizes $\rho(\Gamma_0)$, hence $T$ since $\rho(\Gamma_0)$ is Zariski-dense in $T$. By Fact~\ref{fact.action}, the lightcone $C(p_0)$ can be characterized as the set of points where the orbits of $\rho(\Gamma_0)$ are contained in a photon of $\ein$. It follows that $C(p_0)$ is left invariant by $\Nor(T)$, the normalizer of $T$ in $\PO$. Applying the stereographic projection $\varphi$ of pole $p_0$ (see Section~\ref{sec.stereographic}) we can see $\rho(N_{\pi_1})$ as a subgroup of $\Conf(\RR^{1,2})$. We then show:

\begin{lemm}\label{lem.normalisateur}
Seen in $\Conf(\RR^{1,2})$, the elements of $\rho(N_{\pi_1})$ are contained in the group
\[
G:=\left\lbrace
\begin{pmatrix}
\epsilon_1&- \epsilon_1 y&-\frac{\epsilon_1}{2 }y^2\\ 0&\;\epsilon_2 & \;\epsilon_2 y\\0&\;\;0&\;\;\epsilon_1\\
\end{pmatrix} +
\begin{pmatrix} z\\ x\\ t \\
\end{pmatrix}, \ \ x,z,y,t \in \RR, \ \epsilon_i=\pm 1 \right\rbrace.
\]

In particular, we have the inclusion $\rho(N_{\pi_1}) \subset \Iso(\RR^{1,2})$.
\end{lemm}

\begin{proof}
After performing the stereographic projection $\varphi$, the foliation $\calf_{\Delta}$ restricted to $\ein \setminus C(p_0)$ becomes a foliation of $\RR^{1,2}$. Formula~\eqref{eq.stereographic} for $\varphi$ readily shows that this is the foliation by affine planes of direction $\Span(e_1,e_2)$. Recall (see Section~\ref{sec.symetries}) that the group $T$ corresponds to the group of translations of vectors $v \in \Span(e_1,e_2)$. Since $\Nor(T)$, hence $\rho(N_{\pi_1})$, must preserve this foliation (this is a consequence of Fact~\ref{fact.action}), we see that elements of $\rho(N_{\pi_1})$ belong to the subgroup $G' \subset \Conf(\RR^{1,2})$ comprising all elements of the form:
\begin{equation}\label{eq.formuleG}
\begin{pmatrix}
\lambda \mu&-\lambda \mu y&-\frac{\lambda \mu}{2}y^2\\ 0& \;\;\mu & \;\;\frac{\mu}{\mu}y\\0&\;\;0&\;\frac{\mu}{\lambda}\\
\end{pmatrix} +
\begin{pmatrix} z\\ x\\ t \\
\end{pmatrix}, \ \ x,y,z,t \in \RR \ \ \lambda, \mu \in \RR^*.
\end{equation}

If a matrix $
\begin{smallpmatrix}
\lambda \mu&-\lambda \mu y&\;-\sfrac{\lambda \mu}{2}y^2\\ 0&\mu &\;\sfrac{\mu}{\mu}y\\0&0&\;\sfrac{\mu}{\lambda}\\
\end{smallpmatrix}$ normalizes a lattice in $\Span(e_1,e_2)$, then the determinant of its restriction to $\Span(e_1,e_2)$ is $\pm 1$. It follows that $\mu=\pm \frac{1}{\sqrt{|\lambda|}}$.

We saw that $\rho(\tth)$ belongs to the group $N$, hence has the form:
\[
\rho(\tth)= 
\begin{pmatrix}
1&- y&-\frac{\;\;y^2}{2 }\\ 0&\;\;1 & \;\;y\\0&\;\;0& \;\;1\\
\end{pmatrix} +
\begin{pmatrix}
 z\\ x\\ 0 \\
\end{pmatrix}, \ y \not = 0.
\]

In particular, because $\tth$ normalizes $\Gamma_0$, if $\tau \in \rho(\Gamma_0)$, and if we see $\tau$ as a translation of vector $v \in \Span(e_1,e_2)$, then $\rho(\Gamma_0)$ will also contain $v'=v-A.v$, where \linebreak$A=
\begin{smallpmatrix} 1&y\\0&1\end{smallpmatrix}$. In other words $\rho(\Gamma_0)$ contains a translation of vector $\alpha e_1$, $\alpha \not =0$. The fact that $\rho(N_{\pi_1})$ normalizes the discrete group $\rho(\Gamma_0)$, leads to the relation $\lambda \mu =\pm 1$ in~\eqref{eq.formuleG}. Together with the relation $\mu= \pm \frac{1}{\sqrt{|\lambda|}}$, this leads to $|\mu|=|\lambda|=1$, and the Lemma~\ref{lem.normalisateur} follows.
\end{proof}

Lemma~\ref{lem.normalisateur} says that our $(\ein,\PO)$-structure is actually a $(\RR^{1,2},\Iso(\RR^{1,2}))$-structure. We conclude that there exists $g'$ in the conformal class of $g$ which is flat, and which is preserved by $\Iso(M,g)$. We can thus apply the results of Section~\ref{sec.minkowski}. Theorem~\ref{thm.bieberbach.minkowski} and Proposition~\ref{prop.topologie.minkowski} say that $(M,g')$ is the quotient of $\RR^3$, $\Heis$ or $\Sol$ by a lattice. But $\rho(\pi_1(M)) \subset G$ by Lemma~\ref{lem.normalisateur}, and $G$ does not contain any subgroup isomorphic to $\Sol$. We thus get that $M$ is homeomorphic to $\TT^3$ or to a torus bundle $\TT_A^3$ with $A \subset \SL(2,\ZZ)$ parabolic. This proves points $1)$ and $2)$ of Theorem~\ref{th.theoplat}.



Finally, Carrière's completeness result~\cite{carriere} says that $\delta: \tm \to \RR^{1,2}$ is a conformal diffeomorphism. It follows that if the coordinates associated to $(e_1,e_2,e_3)$ in $\RR^{1,2}$ are $(u,t,v)$, the metric $\tilde{g}$ is of the form
\[
a(u,t,v)(\dt^2+2\du\dv).
\]


It remains to check that the function $a$ depends only on $v$. First, the foliation by planes with direction $\Span(e_1,e_2)$ is totally geodesic. If $\tilde{\nabla}$ denotes the Levi-Civita connection of $\tilde{g}$, we thus have:
\[
0=\tilde{g}(\tilde{\nabla}_{\partial_t}\partial_t,\partial_u)=
-\frac{1}{2}\partial_u.\tilde{g}(\partial_t,\partial_t)=
-\frac{1}{2}\frac{\partial a}{\partial u}.
\]

Identifying $\tth$ and $\rho(\tth)$, we saw that
\[
\tth=\begin{pmatrix}
1&- y&-\frac{\;\;y^2}{2 }\\ 0&\;\;1 &\;\;y\\0&\;\;0&\;\;1\\
\end{pmatrix} +
\begin{pmatrix} z_0 \\ x_0 \\ 0 \\
\end{pmatrix}
\]
where $y \not =0$.

It follows that $\rho(\tth)$ acts on each hyperplane $v=v_0$ by the affine transformation:
\[
\begin{pmatrix} u\\ t
\end{pmatrix} \mapsto 
\begin{pmatrix} u\;-yt + z(v_0)\\\;\;t+ x(v_0)
\end{pmatrix} ,
\]
where $z(v_0)=z_0 - \frac{y^2}{2}v_0$ and $x(v_0)=x_0 + y v_0$.

The group $\Gamma_0$ is generated by two translations $\tau_1,\tau_2$ of (linearly independent) vectors $
\begin{smallpmatrix} a\\ b
\end{smallpmatrix}$ and $\begin{smallpmatrix} c\\ d
\end{smallpmatrix}$ respectively.
The $w$-coordinate of $\tth^k\circ \tau_1^m \circ \tau_2 ^n
\begin{smallpmatrix} u\\ t
\end{smallpmatrix}$
is $t+k x(v_0)+mb+nd$. Because $\tth^k\circ \tau_1^m \circ \tau_2 ^n$ acts isometrically for $\tilde{g}$ this leads to $a(t,v)=a(t+k x(v)+mb+nd,v)$ for every $(k,m,n) \in \ZZ^3$. Since $b$ and $d$ can not be both zero (let say $b \not =0$), and because $x(v)$ and $b$ are rationally independent for almost every value of $w$ (because $y \not =0$), we get that for almost every $w$, $t \mapsto a(t,v)$ is constant. As a consequence, $a=a(v)$, and the fact that it is a periodic function follows easily from the compactness of $M$.

Finally the group $N$ which comprises transformations of the form
\[
\begin{pmatrix}
1&- y&-\frac{ y^2}{2 }\\ 0&1 & y\\0&0& 1\\
\end{pmatrix} +
\begin{pmatrix} z\\ x\\ 0 \\
\end{pmatrix}, \ \ x,y,z \in \RR
\]
is isomorphic to $\Heis$ and acts isometrically on $(\tm,\tilde{g})$. This concludes the proof of Theorem~\ref{th.theoplat}.


\section{Conclusions}\label{sec.fin}

The study made in Section~\ref{sec.homogeneous}, as well as Theorems~\ref{theo.hyperbolic} and~\ref{th.theoplat} provide all possible topologies for a closed $3$-dimensional, orientable and time-orientable, Lorentz manifold with a noncompact isometry group. Those are the $3$-dimensional torus, hyperbolic or parabolic torus bundles, and compact quotients $\Gamma \backslash \tilde{\PSL}(2,\RR)$. Together with the examples provided in Section~\ref{sec.panorama}, this yields Theorem~\ref{thm.topological}.

Let us now look at the geometries which can occur on those manifolds, and prove Theorem~\ref{thm.geometrical}. The manifolds $\Gamma \backslash \tilde{\PSL}(2,\RR)$ occur only in Proposition~\ref{prop.topologieads}. Hence the only metrics on such manifolds which admit a noncompact isometry group are covered by $\tilde{\PSL}(2,\RR)$, endowed with a Lorentzian, non-Riemannian, left-invariant metric. In particular those manifolds $(M,g)$ are locally homogeneous and $(\tm,\tilde{g})$ admits an isometric action of $\tilde{\PSL}(2,\RR)$.

Parabolic torus bundles appear in Proposition~\ref{prop.topologie.minkowski} and Theorem~\ref{th.theoplat}. We saw there that the universal cover is isometric to $\RR^3$ endowed with a metric
\[
a(v)(\dt^2+2\du\dv),
\]
with $a$ smooth and periodic. This universal cover admits an isometric action of $\Heis$. If the manifold $(M,g)$ is locally homogeneous, Proposition~\ref{prop.topologie.minkowski} ensures that $g$ is flat or locally isometric to the Lorentz--Heisenberg metric.

Hyperbolic torus bundles appear only in Proposition~\ref{prop.topologie.minkowski} Theorem~\ref{theo.hyperbolic}. We saw that the universal cover is isometric with $\RR^3$ endowed with a metric $\dt^2+2a(t)\du\dv$, with $a$ smooth and periodic. There is an isometric action of $\Sol$ on this universal cover. The manifold $(M,g)$ is locally homogeneous if and only if it is flat.

Finally, $3$-tori appear in Proposition~\ref{prop.topologie.minkowski} Theorem~\ref{theo.hyperbolic} and Theorem~\ref{th.theoplat}. The metric on the universal cover $\tm$ is provided by those two last theorems, and there is always an isometric action of $\Heis$ or $\Sol$ on $(\tm,\tilde{g})$. Finally, $(M,g)$ is locally homogeneous if and only if it is flat.


Those results altogether prove Theorem~\ref{thm.geometrical} and Corollary~\ref{coro.action.universel}.

%\pagebreak
\appendix
% provoque des avertissements d'overfull à chaque \section, va savoir pourquoi

\section{Some computations}\label{AnexA}
We present here the necessary computations leading to Proposition~\ref{prop.heis}.


\subsection{The curvature module}
We consider on $\RR^3$ the Lorentzian form, with matrix in a basis $e,h,f$ given by $ J=
\begin{smallpmatrix} 0&0&1\\ 0&1&0 \\ 1&0&0\\
\end{smallpmatrix}.$

We call ${\operatorname O}(1,2)$ the subgroup of $\GL(3,\RR)$ preserving the bilinear form determined by $J$. Its Lie algebra is denoted by $\oo(1,2)$, and admits the following basis :
\[
E=
\begin{pmatrix} 0&1&0\\ 0&0& -1\\ 0&0&0\\
\end{pmatrix},\quad
H=
\begin{pmatrix} 1&0&0\\ 0&0&0 \\ 0& 0& -1
\end{pmatrix},\quad
F=
\begin{pmatrix} 0&0&0\\ 1&0&0 \\ 0&-1& 0
\end{pmatrix}.
\]
We thus have the commutation relations $[H,E]=E, [H,F]=-F$ and $[E,F]=H$.

Let $(M,g)$ be $3$-dimensional Lorentz manifold, and denote by $\hm$ its bundle of orthonormal frames. At each $\hx \in \hm$, the curvature $\kappa(\hx)$ is an element of $\Hom(\wedge^2(\RR^3),\oo(1,2))$. This vector space $\Hom(\wedge^2(\RR^3),\oo(1,2))$ is naturally a ${\operatorname O}(1,2)$-module. Choosing $e \wedge h$, $e \wedge f$, $h \wedge f$ as a basis for $\wedge^2(\RR^3)$, and $E,H,F$ as a basis for $\oo(1,2)$, an element of $\Hom(\wedge^2(\RR^3),\oo(1,2))$ is merely given by a $3 \times 3$ matrix, and the action of ${\operatorname O}(1,2)$ on $\Hom(\wedge^2(\RR^3),\oo(1,2))$ corresponds to the conjugation on matrices. For this linear action of ${\operatorname O}(1,2)$, $\Hom(\wedge^2(\RR^3),\oo(1,2))$ can be decomposed as a sum $E_1 \oplus E_3 \oplus E_5$ of irreducible submodules of dimension, $1$, $3$ and $5$ respectively. Because of algebraic Bianchi's identities, the curvature takes values in the $6$-dimensional submodule $E_1 \oplus E_5$, that we call the \emph{curvature module} (see for instance~\cite[Chapter~6]{sharpe}). The submodule $E_1$ comprises scalar matrices, and $\kappa(\hx) \in E_1$ means that the sectional curvature is constant at $x$.



The other irreducible submodule $E_5$ of the curvature module is $5$-dimensional, spanned by the matrices:
\[
\begin{pmatrix} 0&0&1\\ 0&0&0\\ 0&0&0\\
\end{pmatrix},\quad
\begin{pmatrix} 0&1&0\\ 0&0&1\\ 0&0&0\\
\end{pmatrix}
\begin{pmatrix}1&0&0\\ 0&-2 & 0\\0& 0& 1 \\
\end{pmatrix},\quad
\begin{pmatrix} 0&0&0\\ 1&0 &0 \\ 0 & 1 & 0\\
\end{pmatrix},\quad
\begin{pmatrix}
0&0&0 \\ 0&0&0 \\1&0&0 \\
\end{pmatrix}.
\]



We call $\kappa_0$ the element of $\Hom(\wedge^2(\RR^3),\oo(1,2))$ corresponding to the identity matrix, namely $\kappa_0$ maps $e \wedge h$ to $E$, $e \wedge f$ to $H$ and $h \wedge f$ to $F$. We also call $\kappa_1$ the element of $\Hom(\wedge^2(\RR^3),\oo(1,2))$ corresponding to the matrix $\begin{smallpmatrix} 0&0&1\\ 0&0&0\\ 0&0&0\\
\end{smallpmatrix} $.

The two dimensional vector space spanned by $\kappa_0$ and $\kappa_1$ is the set of fixed points of the action of $\{ e^{tE} \}_{t \in \RR}$ on the curvature module.

\subsection{Identification of the \texorpdfstring{$\kiloc$}{kiloc}-algebra}
We consider a parabolic component $\calm$ which is not locally homogeneous. In such a component, the points are either parabolic, or points where the isotropy algebra is $3$-dimensional and the sectional curvature is constant. Constant curvature on a nonempty open subset $U \subset \calm$ would imply that the algebra of Killing fields is $6$-dimensional on $U$, hence on $\calm$. As it contradicts our assumption that $\calm$ is a parabolic component, we conclude that the set of parabolic points is a dense open set $\Omega \subset \calm$. Observe that at a parabolic point $x \in \Omega$, if $X$ a local Killing field around $x$, generating the isotropy ${\mathfrak Is}(x)$, the $1$-parameter group $D_x\varphi_X^t$ is unipotent in $\Omp(T_xM)$. In a suitable basis $(u_1,u_2,u_3)$ of $T_xM$ satisfying $g(u_1,u_3)=1=g(u_2,u_2)$ and all the other products are $0$, the matrix of $D_x\varphi_X^t$ reads
\[
\begin{pmatrix}
1&t&-t^2/2\\
0&1&-t\\
0&0&\;\;1\\
\end{pmatrix}
\]

We quickly check that the only $2$-plane stable by $D_x\varphi_X^t$ is spanned by $u_1$ and $u_2$, so that on $\Omega$, the $\kiloc$-orbits must be lightlike surfaces.


Let us now fix a point $x \in \Omega$. We work in the fiber bundle $\hm$ (and lift all local Killing fields on $M$ to local $\omega$-Killing fields on $\hm$). After multiplying $X$ by a suitable constant, we can find $\hx \in \hm$ in the fiber of $x$ such that $\omega(\hat{X}(\hx))=E$. We now choose $Z$ and $Y$ two local Killing fields around $x$ such that $Z(x)=u_1$ and $Y(x)=u_2$. After adding to $Z$ and $Y$ a suitable multiple of $X$, we can write, at $\hx$:
\[
\omega(\hat{Z})=e+\beta H + \gamma F \quad \text{and}\quad \omega(\hat{Y})=h+\alpha H +\nu F.
\]
The curvature $\kappa(\hx)$ is $\Ad(e^{tE})$-invariant, hence is of the form $\kappa=\sigma \kappa_0 +b \kappa_1$. In particular, the following identities hold at $\hx$:
\begin{equation}\label{eq.courbure}
\kappa(e \wedge h)= \sigma E, \quad \kappa(e \wedge f)=\sigma H, \quad \kappa(h \wedge f)=bE+\sigma F.
\end{equation}
Notice that $\sigma$, $b$, $\alpha$,$\beta$,$\gamma$,$\nu$ depend on $x$ and $\hx$, but since those points are fixed, there will be considered as constant in the sequel.



Cartan's formula $L_{\hat{U}} \omega= \iota_{\hat{U}} d \omega + d(\iota_{\hat{U}} \omega)$ shows that whenever $\hat{U},\hat{V}$ are two $\omega$-Killing fields on $\hm$, the following relation holds:
\begin{equation}\label{courbure.killing}
\omega([\hat{U},\hat{V}])=K(\hat{U},\hat{V})-[\omega(\hat{U}),\omega(\hat{V})].
\end{equation}
Here $K$ is the curvature of $\omega$, as defined in Section~\ref{sec.generalized.curvature}. We recall that it is linked to the curvature function $\kappa$ by the relation $K(\hat{U},\hat{V})=\kappa(\omega(\hat{U}),\omega(\hat{V}))$.

In the sequel, we will call ${\mathcal H}$ the span of $\omega(\hat{Z}),\omega(\hat{X}),\omega(\hat{Y})$ at $\hx$, and we are going to write Equation~\eqref{courbure.killing} at $\hx$, using identities~\eqref{eq.courbure}, when $\hat{U}$ and $\hat{V}$ range over $\hat{Z}$, $\hat{X}$, $\hat{Y}$. For instance, the first equation is:
\begin{align*}
\omega([\hat{Z},\hat{X}])&=-[\omega(\hat{Z}),\omega(\hat{X})] = -[e+ \beta H + \gamma F,E] \\
&= - \beta E + \gamma H \quad\;\,= \omega(- \beta \hat{X}) \;\:+ \gamma H.
\end{align*}

The fact~\ref{fait.nilpotent} that $\kiloc(x)$ is a Lie algebra, together with the property $H \not \in {\mathcal H}$ forces $\gamma$ to vanish. Next, two Killing fields which coincide at $\hx$ must be equal (by freeness of the action of isometries on the orthonormal frames), which implies $[\hat{Z},\hat{X}]= - \beta \hat{X}$. To summarize:
\begin{equation}\label{eq.zx}
\gamma =0\ \text{ and }\ [\hat{Z},\hat{X}]= - \beta \hat{X}.
\end{equation}

We proceed exactly in the same way for the two other equations:
\begin{align*}
\omega([\hat{X},\hat{Y}])& = -[\omega(\hat{X}),\omega(\hat{Y})]\:=-[E, h + \alpha H + \nu F] \\
&= - e + \alpha E - \nu H = \omega(-\hat{Z} + \alpha \hat{X}) +(\beta - \nu) H.
\end{align*}
leads to:
\begin{equation}\label{eq.xy}
\beta=\nu \quad \text{and}\quad [\hat{X},\hat{Y}]=-\hat{Z}+ \alpha \hat{X}.
\end{equation}

Finally
\begin{align*}
\omega([\hat{Z},\hat{Y}])& = \ka(e \wedge h)-[e+ \nu H, h + \alpha H + \nu F]=\sigma E + \alpha e + \nu h + \nu^2 F \\
&= \omega(\alpha \hat{Z} + \sigma \hat{X} + \nu \hat{Y}) -2 \alpha \nu H.
\end{align*}
implies:
\begin{equation}\label{eq.zy}
\alpha \nu =0 \quad\text{and}\quad [\hat{Z},\hat{Y}]= \alpha \hat{Z} + \sigma \hat{X} + \nu \hat{Y}.
\end{equation}

Notice that establishing~\eqref{eq.zx} and~\eqref{eq.xy}, we have actually shown that $\ad(X)$ is a nilpotent endomorphism of $\kiloc(x)$. This property did not use anything special on $x$, so that we actually have:
\begin{fact}\label{fait.nilpotent}
At each $z \in \Omega$, if $U$ is a local Killing field around $z$ generating the isotropy at $z$, then $\ad(U)$ is a nilpotent endomorphism of $\kiloc(z)$.
\end{fact}



At $x$, $Z(x)$ is lightlike and nonzero and $Y(x)$ is spacelike, orthogonal to $Z(x)$. The orthogonal to $Y(x)$ at $x$ is a Lorentzian plane spanned by $Z(x)$ and another vector $w \in T_xM$. Let us call $t \mapsto \gamma(t)$ the geodesic through $x$ satisfying $\dot{\gamma}(0)=w$. Clairault's equation ensures that the quantities $g(\dot{\gamma}(t),Z(\gamma(t)))$, $g(\dot{\gamma}(t),X(\gamma(t)))$ and $g(\dot{\gamma}(t),Y(\gamma(t)))$ do not depend on $t$. In particular, for $t>0$, both $Y(\gamma(t))$ and $X(\gamma(t))$ are orthogonal to $\dot{\gamma}(t)$ while $Z(\gamma(t))$ is not. For $t>0$ small enough, $Y(\gamma(t))$ is still spacelike, hence nonzero, and $\gamma(t)$ belongs to $\Omega$. In particular, the $\kiloc$-orbit at $\gamma(t)$ is $2$-dimensional, so that $Y$ and $X$ must be collinear at $\gamma(t)$. One then has $X(\gamma(t))=\lambda_t Y(\gamma(t))$, for some real $\lambda_t$. Observe finally that $w$ is not fixed by $D_x\phi_X^t$, hence is transverse to the set where $X$ vanishes. In particular, for $t \geq 0$ small, $X(\gamma(t))=0$ only for $t = 0$, and thus $\lambda_t \not =0$ if $t \not =0$.


We claim that those considerations lead necessarily to $\alpha=0$. Indeed, using the bracket relations~\eqref{eq.xy},\eqref{eq.zy} and~\eqref{eq.zx}, we compute
\[
\Trace(\ad(\lambda_t Y - X))=-2 \lambda_t \alpha.
\]
For $t \geq 0$ small, $X,Y,\ Z$ generate $\kiloc(\gamma(t))$, hence $-2 \lambda_t \alpha=0$ because of Fact~\ref{fait.nilpotent}. Since $\lambda_t \not = 0$ if $t \not =0$, we get $\alpha=0$. Injecting this data in equation~\eqref{eq.zy} and~\eqref{eq.zx}, we find that the matrix of $\ad(\lambda_t Y - X)$ in the basis $Y,Z,X$ is:
\[
\begin{pmatrix}
0 & -\lambda_t \nu &0\\
1 & 0 & \lambda_t\\
0 & -\lambda_t \sigma -\nu & 0\\
\end{pmatrix}.
\]
The characteristic polynomial of $\ad(\lambda_t Y - X)$ is
\[
Q(x)=-x^3-\lambda_t x(\lambda_t \sigma + 2 \nu).
\]
Hence, the nilpotency of $\ad(\lambda_t Y - X)$ (Fact~\ref{fait.nilpotent}) implies
\begin{equation}\label{eq.relation}
\lambda_t \sigma + 2 \nu=0,
\end{equation}
If $\sigma \not =0$, we get that $t \mapsto \lambda_t$ is constant, which is not the case since we observed that $\lambda_0=0$ but $\lambda_t \not =0$ for $t>0$ small. We end up with the equality $\sigma=\nu=0$. The vector fields $Z,X,Y$ then satisfy the bracket relations:
\[
[Z,X]=0=[Z,Y], \text{ and } [X,Y]=-Z,
\]
showing that Lie algebra $\kiloc(x)$ is isomorphic to $\heis$. We also proved that $\sigma$, the scalar curvature at $x$, vanishes, but since $x$ was arbitrary in the open set $\Omega$, we finally get the vanishing of the scalar curvature on $\Omega$, and then on $\calm$ by density.


\subsection{Description of the \texorpdfstring{$\kiloc$}{kiloc}-orbits}

The fact~\ref{fait.nilpotent} that the local Killing algebra is isomorphic to $\heis$ shows that no point in $\calm$ has a $3$-dimensional isotropy algebra. Indeed, the isotropy representation at those points would yield an embedding $\heis \to \oo(1,2)$, what is impossible. We thus get $\Omega=\calm$, and all the $\kiloc$-orbits on $\calm$ are $2$-dimensional and lightlike.

On the other hand, since the isotropy algebra ${\mathfrak{Is}}(x)$ generates a parabolic $1$-parameter subgroup of $\Omp(1,2)$ at each $x$, there is a totally geodesic lightlike hypersurface $F(x)$, whose tangent space is left invariant by the isotropy (see~\cite[Lemma~3.5]{sorin.ghani} and its proof). We already observed that at $x \in \calm$, the local isotropy preserves only one $2$-plane of $T_xM$. This implies that the $\kiloc$-orbits are everywhere tangent to a leaf of a totally geodesic foliation of $\calm$, hence the $\kiloc$-orbits are themselves totally geodesic. This concludes the proof of Proposition~\ref{prop.heisenberg.flat}.

\section{About the completeness of closed Lorentz--Heisenberg manifolds}\label{ann.completeness}

Our aim here is to explain how to adapt the proof of~\cite[Proposition~8.1]{sorin.ghani}, and get Theorem~\ref{thm.bieberbach.minkowski} for closed Lorentz--Heisenberg manifolds without using Theorem~\ref{thm.ghaniclassification}.

We call $\heis$ the $3$-dimensional Heisenberg Lie algebra, namely the Lie algebra generated by $Z,Y,X$, with relation $[Y,X]=Z$. The Lorentz--Heisenberg metric $g_{LH}$ on the Lie group $\Heis$ is the left-invariant Lorentz metric, which is given on $\heis$ by $\langle X,Y\rangle=1$, $\langle Z,Z\rangle=1$, and all other products are zero. It is explained in~\cite[Section~4.1]{sorin.ghani}, that the Lie algebra of Killing fields on $(\Heis,g_{LH})$ is $4$-dimensional, and is generated by $X,Y,Z$ as well as a fourth element $T$ satisfying the bracket relations $[T,Y]=Y$, $[T,X]=-X$ and $[T,Z]=0$. The identity component $G$ of $\Iso(\Heis, g_{LH})$ is thus isomorphic to a semi-direct product $\RR \ltimes \heis$, where the $\RR$-factor integrates into a group of hyperbolic automorphisms of $\heis$.

If $(M,g)$ is a closed Lorentz manifold locally modelled on $(\Heis, g_{LH})$, we consider a developing map $\delta: \tilde{M} \to \Heis$, and the corresponding holonomy morphism $\rho: \pi_1(M) \to \Iso(\Heis, g_{LH})$. Since $G$ has finite index in $\Iso(\Heis, g_{LH})$, we can replace $(M,g)$ by a finite cover and assume that $\Gamma:=\rho(\pi_1(M))$ is contained in $G$. Since $Z$ is central in $G$, the vector field $\delta^*(Z)$ projects on $(M,g)$ to a Killing vector field, the flow of which will be called \emph{the characteristic flow} on $(M,g)$, denoted ${\varphi}_Z^t$. The main part of the proof of~\cite[Proposition~8.1]{sorin.ghani} deals with the case where the characteristic flow ${\varphi}_Z^t$ is relatively compact. This part is independent of~\cite{zeghibflot}. Our goal is thus to provide a self-contained proof of the



\begin{prop}\label{prop.complete2}
Let $(M,g)$ be a closed Lorentz manifold locally modelled on $(\Heis,g_{LH})$. Then the characteristic flow on $M$ is relatively compact.
\end{prop}

\begin{proof}
The proof will be by contradiction. We thus assume in the following that ${\varphi}_Z^t$ is not relatively compact.

\begin{lemm}\label{lem.lightlike.complete}
Let $(M,g)$ be a closed Lorentz manifold locally modelled on $(\Heis,g_{LH})$. If the characteristic flow is not relatively compact, then lightlike geodesics which are orthogonal to the characteristic flow are complete.
\end{lemm}

\begin{proof}
We consider a vector field $E_3$ on $M$ which is tangent to the orbits of ${\varphi}_Z^t$ and satisfies $g(E_3,E_3)=1$ (recall that the orbits of ${\varphi}_Z^t$ are spacelike by definition of $g_{LH}$). We get a distribution $E_3^{\perp}$ of Lorentzian $2$-planes on $M$. After replacing $M$ by a finite cover, we may assume that there exist on $M$ two other vector fields $E_1,E_2$, spanning $E_3^{\perp}$ at each point, and satisfying the relations $g(E_1,E_1)=g(E_2,E_2)=0$, $g(E_1,E_2)=1$.

Because ${\varphi}_Z^t$ is a flow of isometries, and because it preserves $E_3$, there exists a smooth cocycle $\lambda: M \times \RR \to \RR_+^*$, such that $D_{x}\varphi_Z^{t}(E_1(x))= \lambda(x,t) E_1(x)$ (resp.~$D_{x}\varphi_Z^{t}(E_2(x))= \lambda(x,t)^{-1} E_2(x)$) for all $x \in M$ and $t \in \RR$.

Let us pick $x \in M$. If $(t_k)$ is any sequence of $\RR$ satisfying $\lim_{k \to \infty}|t_k|=+\infty$, we observe that $\lambda(x,t_k)$ can not stay in a compact subset of $\RR_+^*$. Indeed, this would mean that the $1$-jet of $\varphi_Z^{t_k}$ at $x$ remains bounded, implying that ${\varphi}_Z^t$ is relatively compact in $\Iso(M,g)$. This is in contradiction with our standing assumption.

Hence, for any sequence $(t_k)$ such that $\lim_{k \to \infty}|t_k|=+\infty$, we must have
\[
\lim_{k \to \infty}\lambda(x,t_k)=0 \text{ or } \lim_{k \to \infty}\lambda(x,t_k)=+ \infty.
\]
If only the second possibility occurs, we get that $\lim_{t \to \pm \infty} \lambda(x,t)= + \infty$. In particular, since the map $t \mapsto \lambda(x,t)$ is continuous, we can find two sequences $(t_k)$ and $(s_k)$ satisfying $\lim_{k \to \infty} t_k =\lim_{k \to \infty} s_k =+ \infty$ and $\lambda(x,t_k)=\lambda(x, -s_k)$. Let us call $y_k:={\varphi}_Z^{-s_k}(x)$ and $z_k:={\varphi}_Z^{t_k}(x)$. Considering subsequences if necessary, we may assume that $(y_k)$ and $(z_k)$ are convergent sequences. We have the relation $D_{y_k}{\varphi}_Z^{t_k+s_k}(E_1(y_k))=E_1(z_k)$ for all $k \in \NN$, which in turn forces $D_{y_k}{\varphi}_Z^{t_k+s_k}(E_2(y_k))=E_2(z_k)$. Hence the $1$-jet of $\varphi_Z^{s_k+t_k}$ is bounded along the sequence $(y_k)$. Since $\lim_{k \to \infty} t_k+s_k=+ \infty$, this leads to ${\varphi}_Z^t$ being relatively compact in $\Iso(M,g)$, and we get again a contradiction.

As a consequence, one can find a sequence $(t_k)$ such that $\lim_{k \to \infty}\lambda(x,t_k)=0$. By compactness of $M$, there exists $\tau>0$ such that the geodesics of direction~$E_1(y)$ or $E_2(y)$ are defined on $(-\tau, \tau)$, for all $y \in M$. The geodesic of direction~$E_1(x)$ is thus defined on $(-\lambda(x,t_k)^{-1} \tau, \lambda(x,t_k)^{-1} \tau)$ for every $k \in \NN$. It follows that this geodesic is complete.

A similar argument involving the cocycle $\lambda(x,t)^{-1}$ shows that the geodesic of direction~$E_2(x)$ is complete too.
\end{proof}


\begin{lemm}\label{lem.dev}
Let $(M,g)$ be a closed Lorentz manifold locally modelled on $(\Heis,g_{LH})$. If the characteristic flow is not relatively compact, then the developing map $\delta: \tilde{M} \to \Heis$ is a diffeomorphism.
\end{lemm}

\begin{proof}
For every $g \in \Heis$, the curves $t \mapsto g e^{tY}$ and $t \mapsto g e^{tX}$ are lightlike geodesics for the Lorentz--Heisenberg metric $g_{LH}$. Those geodesics are orthogonal to the central flow $e^{tZ}$, hence their inverse image by the map $\delta$ are lightlike geodesics on $(\tilde{M},\tilde{g})$ orthogonal to the lift of $\varphi_Z^t$. Lemma~\ref{lem.lightlike.complete} above, as well as the completeness of the flow $\varphi_Z^t$, ensures that any piecewise smooth curve of $\Heis$, obtained by successively flowing along the right multiplication under $e^{tY}$, $e^{tX}$ and $e^{tZ}$, can be lifted to $\tilde{M}$ through the map $\delta$, with arbitrary initial condition. Let us fix $\epsilon >0$ such that the map $f:(r,s,t) \mapsto \Heis$ defined by $f(r,s,t)=e^{rY}e^{sX}e^{tZ}$ yields a diffeomorphism from $(-\epsilon, \epsilon) \times (-\epsilon, \epsilon) \times (-\epsilon, \epsilon)$ to its image $B_{\epsilon}$. Our previous remark ensures that whenever $\tilde{x}$ is a point of $\tilde{M}$, such that $\delta(\tilde{x})=g$, then $gB_{\epsilon}$ can be lifted to an open set $U$ containing $\tilde{x}$, such that $\delta: U \to gB_{\epsilon}$ is a diffeomorphism. One easily deduces that $\delta: \tilde{M} \to \Heis$ has the path-lifting property, hence is a diffeomorphism.
\end{proof}

Lemma~\ref{lem.dev} ensures that $\Gamma:=\rho(\pi_1(M))$ is a discrete subgroup of $G \simeq \RR \ltimes \Heis$ that acts freely properly and cocompactly on $\Heis$. We call $\Gamma':=\Gamma \cap \Heis$.

If $\Gamma'= \{1 \}$, then $\Gamma$ projects injectively on the $\RR$-factor of $G\simeq \RR \ltimes \Heis$. It is thus abelian. It is easily checked that nontrivial connected abelian subgroups of $G$ have dimension~$1$ or $2$. Looking at the identity component of the Zariski closure of $\Gamma$ in $G$, we get that $\Gamma$ is isomorphic to $\ZZ$ or $\ZZ^2$. A cohomological dimension argument shows that the action of $\Gamma$ can not be cocompact on $\Heis$.

Discrete subgroups of $\Heis$ which are not cyclic must intersect nontrivially the center of $\Heis$. On the other hand, $\Gamma$ intersects the center of $\Heis$ trivially, otherwise $\varphi_Z^t$ would be a periodic flow, contradicting our assumption that it is not relatively compact. We conclude that $\Gamma'$ is infinite cyclic, and normalized by $\Gamma$. Considering a finite index subgroup, we get that $\Gamma$ centralizes $\Gamma'$. Because $\Gamma'$ is not contained in the center of $\Heis$, its centralizer in $G$ is $1$-dimensional. We conclude that $\Gamma \simeq \ZZ$, contradicting again that its action on $\Heis$ is cocompact.
\end{proof}





\begin{thebibliography}{MMMMMM}
\bibitem[Amo79]{amores} Amores, Angel M. Vector fields of a finite type G-structure, J. Differ. Geom., Volume 14 (1979) no. 1, pp. 1-6 \href{https://ahl.centre-mersenne.org/item/AHL_2020__3__407_0/#amores}{Publisher reference}.
\bibitem[AS97]{adams.stuck} Adams, Scot; Stuck, Garett The isometry group of a compact Lorentz manifold. I, II, Invent. Math., Volume 129 (1997) no. 2, p. 239-261, 263–287 \href{https://ahl.centre-mersenne.org/item/AHL_2020__3__407_0/#adams.stuck}{Publisher reference}.
\bibitem[BCD + 08]{primer} Barbot, Thierry; Charette, Virginie; Drumm, Todd; Goldman, William M.; Melnick, Karin A primer on the (2+1) Einstein universe, Recent developments in pseudo-Riemannian geometry (ESI Lectures in Mathematics and Physics), European Mathematical Society, 2008, pp. 179-229 \href{https://ahl.centre-mersenne.org/item/AHL_2020__3__407_0/#primer}{Publisher reference}.
\bibitem[BM16]{mounoud} Bavard, Christophe; Mounoud, Pierre Extension maximale et classification des tores lorentziens munis d’un champ de Killing (2016) (\url{https://arxiv.org/abs/1510.01253v2}) \href{https://ahl.centre-mersenne.org/item/AHL_2020__3__407_0/#mounoud}{Publisher reference}.
\bibitem[Car89]{carriere} Carrière, Yves Autour de la conjecture de L. Markus sur les variétés affines, Invent. Math., Volume 95 (1989) no. 3, pp. 615-628 \href{https://ahl.centre-mersenne.org/item/AHL_2020__3__407_0/#carriere}{Publisher reference}.
\bibitem[DG91]{dambragro} D’Ambra, Giuseppina; Gromov, Mikhael L. Lectures on transformations groups: geometry and dynamics, Surveys in Differential Geometry Vol. I: Proceedings of the conference on geometry and topology, held at Harvard University, (Cambridge, MA, 1990), American Mathematical Society, 1991, pp. 19-111 \href{https://ahl.centre-mersenne.org/item/AHL_2020__3__407_0/#dambragro}{Publisher reference}.
\bibitem[DM15]{sorin.karin} Dumitrescu, Sorin; Melnick, Karin Quasihomogeneous three-dimensional real-analytic Lorentz metrics do not exist, Geom. Dedicata, Volume 179 (2015), pp. 229-253 \href{https://ahl.centre-mersenne.org/item/AHL_2020__3__407_0/#sorin.karin}{Publisher reference}.
\bibitem[DZ10]{sorin.ghani} Dumitrescu, Sorin; Zeghib, Abdelghani Géométries Lorentziennes de dimension 3: classification et complétude, Geom. Dedicata, Volume 149 (2010), pp. 243-273 \href{https://ahl.centre-mersenne.org/item/AHL_2020__3__407_0/#sorin.ghani}{Publisher reference}.
\bibitem[D’A88]{dambra} D’Ambra, Giuseppina Sometry groups of Lorentz manifolds, Invent. Math., Volume 92 (1988) no. 3, pp. 555-565 \href{https://ahl.centre-mersenne.org/item/AHL_2020__3__407_0/#dambra}{Publisher reference}.
\bibitem[FG83]{fried.goldman} Fried, David; Goldman, William M. Three-dimensional affine crystallographic groups, Adv. Math., Volume 47 (1983) no. 1, pp. 1-49 \href{https://ahl.centre-mersenne.org/item/AHL_2020__3__407_0/#fried.goldman}{Publisher reference}.
\bibitem[Fra02]{charlesthese} Frances, Charles Géométrie et dynamique lorentziennes conformes, Ph. D. Thesis, ENS Lyon (France) (2002) \href{https://ahl.centre-mersenne.org/item/AHL_2020__3__407_0/#charlesthese}{Publisher reference}.
\bibitem[Fra18]{frances.open} Frances, Charles Variations on Gromov’s open-dense orbit theorem, Bull. Soc. Math. Fr., Volume 146 (2018) no. 4, pp. 713-744 \href{https://ahl.centre-mersenne.org/item/AHL_2020__3__407_0/#frances.open}{Publisher reference}.
\bibitem[FZ02]{fisher.zimmer} Fisher, David; Zimmer, Robert J. Geometric lattice actions, entropy and fundamental groups, Comment. Math. Helv., Volume 77 (2002) no. 2, pp. 326-338 \href{https://ahl.centre-mersenne.org/item/AHL_2020__3__407_0/#fisher.zimmer}{Publisher reference}.
\bibitem[GK84]{goldman.kami} Goldman, William M.; Kamishima, Yoshinobu The fundamental group of a compact flat Lorentz space form is virtually polycyclic, J. Differ. Geom., Volume 19 (1984) no. 1, pp. 233-240 \href{https://ahl.centre-mersenne.org/item/AHL_2020__3__407_0/#goldman.kami}{Publisher reference}.
\bibitem[God91]{godbillon} Godbillon, Claude Feuilletages. Études géométriques, Progress in Mathematics, 98, Birkhäuser, 1991 \href{https://ahl.centre-mersenne.org/item/AHL_2020__3__407_0/#godbillon}{Publisher reference}.
\bibitem[Gro88]{gromov} Gromov, Michael M. Rigid transformation groups, Géométrie Différentielle (Travaux en Cours), Volume 33, Hermann, 1988, pp. 65-141 \href{https://ahl.centre-mersenne.org/item/AHL_2020__3__407_0/#gromov}{Publisher reference}.
\bibitem[Kli96]{klingler} Klingler, Bruno Complétude des variétés lorentziennes à courbure constante, Math. Ann., Volume 306 (1996) no. 2, pp. 353-370 \href{https://ahl.centre-mersenne.org/item/AHL_2020__3__407_0/#klingler}{Publisher reference}.
\bibitem[KN63]{kobanomi} Kobayashi, Shoshichi; Nomizu, Katsumi Foundations of differential geometry I, Interscience Publishers, 1963 \href{https://ahl.centre-mersenne.org/item/AHL_2020__3__407_0/#kobanomi}{Publisher reference}.
\bibitem[KR85]{kulkarni.raymond} Kulkarni, Ravi; Raymond, Franck 3-dimensional Lorentz space-forms and Seifert fiber spaces, J. Differ. Geom., Volume 21 (1985) no. 2, pp. 231-268 \href{https://ahl.centre-mersenne.org/item/AHL_2020__3__407_0/#kulkarni.raymond}{Publisher reference}.
\bibitem[Mel11]{melnick} Melnick, Karin A Frobenius theorem for Cartan geometries, with applications, Enseign. Math., Volume 57 (2011) no. 1-2, pp. 57-89 \href{https://ahl.centre-mersenne.org/item/AHL_2020__3__407_0/#melnick}{Publisher reference}.
\bibitem[MS39]{myers-steenrod} Myers, Sumner Byron; Steenrod, Norman Earl The group of isometries of a Riemannian manifold, Ann. Math., Volume 40 (1939) no. 2, pp. 400-416 \href{https://ahl.centre-mersenne.org/item/AHL_2020__3__407_0/#myers-steenrod}{Publisher reference}.
\bibitem[Nom53]{nomizu2} Nomizu, Katsumi On the group of affine transformations of an affinely connected manifold, Proc. Am. Math. Soc., Volume 4 (1953), pp. 816-823 \href{https://ahl.centre-mersenne.org/item/AHL_2020__3__407_0/#nomizu2}{Publisher reference}.
\bibitem[Orl72]{orlik} Orlik, Peter Seifert Manifolds, Lecture Notes in Mathematics, 291, Springer, 1972 \href{https://ahl.centre-mersenne.org/item/AHL_2020__3__407_0/#orlik}{Publisher reference}.
\bibitem[Péc16]{vincent} Pécastaing, Vincent On two theorems about local automorphisms of geometric structures, Ann. Inst. Fourier, Volume 66 (2016) no. 1, pp. 175-208 \href{https://ahl.centre-mersenne.org/item/AHL_2020__3__407_0/#vincent}{Publisher reference}.
\bibitem[Sal99]{saleinthese} Salein, François Variétés anti-de Sitter de dimension 3, Ph. D. Thesis, ENS Lyon, France (1999) (\url{http://www.umpa.ens-lyon.fr/~zeghib/these.salein.pdf}) \href{https://ahl.centre-mersenne.org/item/AHL_2020__3__407_0/#saleinthese}{Publisher reference}.
\bibitem[Sco83]{scott} Scott, Peter The geometries of 3-manifolds, Bull. Lond. Math. Soc., Volume 15 (1983) no. 5, pp. 401-487 \href{https://ahl.centre-mersenne.org/item/AHL_2020__3__407_0/#scott}{Publisher reference}.
\bibitem[Sha97]{sharpe} Sharpe, Richard W. Differential Geometry: Cartan’s generalization of Klein’s Erlangen Program, Graduate Texts in Mathematics, 166, Springer, 1997 \href{https://ahl.centre-mersenne.org/item/AHL_2020__3__407_0/#sharpe}{Publisher reference}.
\bibitem[Tho14]{tholozan} Tholozan, Nicolas Uniformisation des variétés pseudo-riemanniennes localement homogènes, Ph. D. Thesis, Université de Nice Sophia-Antipolis, France (2014) \href{https://ahl.centre-mersenne.org/item/AHL_2020__3__407_0/#tholozan}{Publisher reference}.
\bibitem[Zeg96]{zeghibflot} Zeghib, Abdelghani Killing fields in compact Lorentz 3-manifolds, J. Differ. Geom., Volume 43 (1996) no. 4, pp. 859-894 \href{https://ahl.centre-mersenne.org/item/AHL_2020__3__407_0/#zeghibflot}{Publisher reference}.
\bibitem[Zeg98]{zeghib.isom} Zeghib, Abdelghani Sur les espaces-temps homogènes. [Homogeneous spacetimes], The Epstein birthday schrift (Geometry and Topology Monographs), Volume 1, Geometry and Topology Publications, 1998, pp. 551-576 \href{https://ahl.centre-mersenne.org/item/AHL_2020__3__407_0/#zeghib.isom}{Publisher reference}.
\bibitem[Zeg99a]{zeghibfoliation} Zeghib, Abdelghani Geodesic foliations in Lorentz 3-manifolds, Comment. Math. Helv., Volume 74 (1999) no. 1, pp. 1-21 \href{https://ahl.centre-mersenne.org/item/AHL_2020__3__407_0/#zeghibfoliation}{Publisher reference}.
\bibitem[Zeg99b]{zeghibisom1} Zeghib, Abdelghani Isometry groups and geodesic foliations of Lorentz manifolds. I. Foundations of Lorentz dynamics., Geom. Funct. Anal., Volume 9 (1999) no. 4, pp. 775-822 \href{https://ahl.centre-mersenne.org/item/AHL_2020__3__407_0/#zeghibisom1}{Publisher reference}.
\bibitem[Zeg99c]{zeghibisom2} Zeghib, Abdelghani Isometry groups and geodesic foliations of Lorentz manifolds. II. Geometry of analytic Lorentz manifolds with large isometry group, Geom. Funct. Anal., Volume 9 (1999) no. 4, pp. 823-854 \href{https://ahl.centre-mersenne.org/item/AHL_2020__3__407_0/#zeghibisom2}{Publisher reference}.
\bibitem[Zim86]{zimmer.iso} Zimmer, Robert J. On the automorphism group of a compact Lorentz manifold and other geometric manifolds, Invent. Math., Volume 83 (1986) no. 3, pp. 411-424 \href{https://ahl.centre-mersenne.org/item/AHL_2020__3__407_0/#zimmer.iso}{Publisher reference}.
\end{thebibliography}
\end{document}
